\documentclass[11pt,oneside]{book}

\usepackage[a4paper,margin=1in]{geometry}

\usepackage[T1]{fontenc}

\usepackage[utf8]{inputenc}

\usepackage{amsmath,amssymb}

\usepackage{booktabs,longtable,array}

\usepackage{xcolor}

\usepackage{titlesec}

\usepackage[most]{tcolorbox}

\usepackage{hyperref}

\hypersetup{colorlinks=true,linkcolor=blue!55!black,urlcolor=blue!55!black,citecolor=blue!55!black}

\pagestyle{plain}

\emergencystretch=3em

\definecolor{quantumgreen}{HTML}{3F6F5A}

\definecolor{softgray}{HTML}{F5F6F2}

\definecolor{ink}{HTML}{1C1E1A}

\titleformat{\chapter}[display]{\normalfont\huge\bfseries\color{ink}\raggedright}{\chaptertitlename\ \thechapter}{14pt}{\Huge}

\newtcolorbox{claimbox}{colback=softgray,colframe=quantumgreen!70!black,arc=2pt,boxrule=0.7pt,left=8pt,right=8pt,top=6pt,bottom=6pt}

\newenvironment{centeredalign}{\[\begin{gathered}[c]}{\end{gathered}\]}

\newcommand{\ket}[1]{\lvert #1\rangle}

\newcommand{\bra}[1]{\langle #1\rvert}

\begin{document}

\frontmatter

\hypersetup{pageanchor=false}

\begin{titlepage}

\centering

\vspace*{1.7cm}

{\Huge\bfseries Quantum Theory\par}

\vspace{0.12cm}

{\Huge\bfseries As A Mechanism Tree\par}

\vspace{0.3cm}

{\large Expository companion\par}

\vspace{1.2cm}

{\large Synthetix Institute\par}

{\large Edition dated 2026-09-26 10:36 UTC\par}

\vfill

\begin{claimbox}

\noindent Interactions generate phases, correlations and transitions. Their observable consequences connect the preparation of a quantum system to its later behaviour. This book develops those mechanisms through equations and follows their realization in different physical systems.

\end{claimbox}

\vspace{0.5cm}

\begin{minipage}{0.94\linewidth}\small\raggedright

\textbf{Scope and authorship.} This companion develops established quantum theory through its physical mechanisms. The mechanism tree and its cross-topic connections are the authors' organization of those mechanisms. References identify the particular relations used in each treatment.

\end{minipage}

\end{titlepage}

\hypersetup{pageanchor=true}

\tableofcontents

\chapter*{About This Edition}

This companion orders quantum descriptions by the states, interactions, constraints and measurements that enter their predictions. Topic-specific derivations develop individual mechanisms. The role division is the authors' organizing choice; the structure of the tree is not presented as a partition independently recovered from the archive. The interpretations and explanatory equations are editorial constructions, rather than new quantum laws inferred from the archive.

All 62 topic-specific treatments are retained. The reference index records 63 subjects whose full-edition entries repeat shared equations, 4 alternative names, and 17 historical or interpretive entries. An indexed subject is not counted as a developed treatment. The complete reference edition and topic files remain available separately.

References follow the calculations they support. Links headed Relations In The Original Papers identify particular displays in those articles; Equation Sources identifies relations aligned to corpus records. The authors develop the explanatory calculations in this volume.

\mainmatter

\part{Mechanisms And Predictions}

\chapter{Quantum Mechanisms And Their Predictions}\label{chap:mechanisms-predictions}

An interaction between two quantum systems can make a property of one depend
on their joint preparation. This dependence is responsible both for useful
operations, such as entangling gates, and for the loss of coherence when the
second system is uncontrolled. Explaining either effect requires following
what the interaction changes and how that change becomes observable. The same
calculation can then describe a controlled pair of spins, two coupled effective
two-level systems, or a spin interacting with a fluctuating neighbour, provided
their physical interactions realize the specified Hamiltonian.

Consider two spins with Hamiltonian \(H=\hbar g Z_1Z_2\), where \(g\) is an
angular frequency and \(X_j,Y_j,Z_j\) are Pauli operators acting on spin \(j\).
The eigenvalue of \(Z_2\) determines the sign of the field experienced by spin
1. A superposition of these eigenstates therefore subjects spin 1 to two
conditional rotations and can entangle the pair. For a joint density operator
\(\rho(t)\), the transverse polarization \(x(t)=\operatorname{Tr}[\rho(t)X_1]\)
couples to the correlation \(c(t)=\operatorname{Tr}[\rho(t)Y_1Z_2]\).
The Heisenberg equation gives
\begin{centeredalign}
\dot x=-2gc,\qquad \dot c=2gx,\\
x(t)=x(0)\cos(2gt)-c(0)\sin(2gt).
\end{centeredalign}
Thus a measurement on one spin can reveal a correlation established before
the interaction began. The oscillation is produced by the exchange between
two components of the joint state, even though only one of them is measured.

The dependence on the joint preparation is visible without assuming an
entangled initial state. Let \(I_4\) be the identity on the two-spin space,
and choose real numbers \(x_0,c_0\) satisfying \(x_0^2+c_0^2\leq1\). The states
\begin{centeredalign}
\rho_\pm=\tfrac14(I_4+x_0X_1\pm c_0Y_1Z_2),\\
\operatorname{Tr}_2\rho_\pm=\tfrac12(I_2+x_0X_1)
\end{centeredalign}
have identical reduced states for spin 1. They nevertheless give polarizations
whose difference is \(-2c_0\sin(2gt)\). Positivity follows because \(X_1\) and
\(Y_1Z_2\) anticommute and square to the identity: the eigenvalues of either
state are \((1\pm\sqrt{x_0^2+c_0^2})/4\). The correlation is consequently an
independent physical part of the preparation, not a relabelling of the initial
polarization.

Measuring \(Z_1\) instead would give a constant result, since
\([H,Z_1]=0\). The interaction has not disappeared; that observable is
insensitive to its action. A mechanism in this book means the physical
dependence connecting an interaction and an admissible preparation to a
specified observable consequence. It includes the degrees of freedom through
which the effect is transmitted. A Hamiltonian determines their evolution;
the preparation and measurement determine which part of that evolution is
tested. This distinction is essential when comparing experiments: equal
Hamiltonians with different preparations can produce different signals,
whereas different physical devices can implement the same dependence.

For the polarization above, retaining \((x,c)\) gives a complete pair of
evolution equations. Eliminating \(c\) expresses the same dependence as
\begin{centeredalign}
\dot x(t)=-2gc_0-4g^2\int_0^t x(s)\,\mathrm ds.
\end{centeredalign}
The integral records the correlation generated during the motion; the first
term retains the correlation present at preparation. This is where closure
enters the explanation. A description by \(x\) alone must carry both terms,
or restrict the preparations and observations so that the omitted information
cannot affect a prediction. Closure describes the sufficiency of those chosen
variables. The mechanism is the interaction that made the correlation
physically relevant in the first place.

This example specifies the operational language used in the book. The
carrier \(\Xi\) is the set of density operators on
\(\mathbb C^2\otimes\mathbb C^2\), and the operation
\(\Omega[\rho]=-i[H,\rho]/\hbar\) evolves them. Positivity and unit trace
are the admissibility conditions \(C\). The observable map
\(R_X(\rho)=\operatorname{Tr}(\rho X_1)\) gives the polarization, and the
protocol \(P\) specifies the joint preparation, evolution time and
measurement. A pair of coupled spins realizing \(H\) supplies the physical
implementation \(A\). Thus \(I_{\mathrm{op}}=((\Omega,\Xi);C,R_X,P)\)
records the ingredients needed to predict this signal; their roles follow
from the calculated evolution and measurement.

The reduction to \(x\) shows why these roles cannot be discarded
independently. Define \(\alpha_x(\rho)=R_X(\rho)\). For the two preparations
above, with \(c_0\ne0\),
\begin{centeredalign}
\alpha_x(\rho_+)=\alpha_x(\rho_-)=x_0,\qquad
\alpha_x(\Omega[\rho_\pm])=\mp 2gc_0.
\end{centeredalign}
No evolution law depending only on \(x_0\) can reproduce both initial slopes.
The correlation \(c\) is the extra state coordinate required for autonomous
evolution; when it is eliminated instead, its initial value and the history
integral in the preceding equation retain precisely the information that
\(x\) lost. A state map between descriptions must therefore preserve the
evolution and the specified observable consequence.

The organization of quantum theory follows these dependencies. A spin
Hamiltonian acts on a tensor product of spin spaces; exchange symmetry changes
the admissible many-particle states; a detector couples a particular
observable to an apparatus. These ingredients cannot be assembled in an
arbitrary order. An operator must be defined on the chosen states before its
evolution can be calculated. Predicting a later measurement then requires
both that evolution and the initial preparation. The nested description used below
expresses this dependence. Its roles identify the ingredients of a mechanism
after their physical function has been established.

Constructing another realization means finding an interaction and a
preparation that reproduce the relevant equations and measurement relation.
For the two-spin example, changes of basis are harmless when they transform
the Hamiltonian, state and observable together. Additional couplings require
new commutators: they can connect \(x\) and \(c\) to further correlations and
change the measured signal. The subsequent chapters develop this operation
across interference, quantum statistics, measurement, fields and information
processing. Explicit calculations show which physical dependence is preserved
and which additional term or degree of freedom becomes necessary.


\chapter{The Physical Identity Of A Quantum Mechanism}\label{chap:physical-identity}
\begin{claimbox}
Quantum predictions join an admissible state space to an operation, the conditions under which that operation is defined, and an observable consequence. The same equation can play different roles when its domain, state, or measurement changes.
\end{claimbox}
\section{Nested Construction}
\[(\Omega,\Xi)\longrightarrow M\longrightarrow I_{\mathrm{op}}=(M;C,R,P)\longrightarrow I_{\mathrm{real}}=(I_{\mathrm{op}};A)\]
The roles \(\Xi\), \(\Omega\), \(C\), \(R\), \(P\), and \(A\) are those introduced with the two coupled spins of Chapter~\ref{chap:mechanisms-predictions}; the chain nests them. The pair \(M=(\Omega,\Xi)\) identifies the law together with the state space on which it is meaningful. Closure, an observable map, and a protocol complete it to the operational identity, and the realization adds the material or field content, scales, parameter values, initial and boundary data, geometry, drives, and apparatus of a particular calculation or experiment.
\section{State Space And Operation Must Agree}
\begin{centeredalign}
\Xi_Q = (\mathcal H,\mathcal D,\mathcal S,\mathcal A_{\rm obs}),\\
q=\rho\in\mathcal S(\mathcal H),\\
\Omega_Q = \{G,\mathcal E,\{E_y\},\mathcal U,\ldots\},\\
M_Q = (\Omega_Q,\Xi_Q).
\end{centeredalign}
The state space \(\Xi_Q\) contains the Hilbert or Fock space, admissible density operators, operator domains, and, when needed, a subsystem or observable algebra. The operation \(\Omega_Q\) contains generators, channels, observables, symmetry actions, projections, or their compositions. Neither factor can be interpreted without the other: the state space fixes what operations are defined, and the operation selects which structure of the state enters a prediction.
\section{Completion And Realization}
Normalization, positivity, self-adjointness, domains, and gauge constraints form the closure \(C_Q\). Spectral measures, positive operator-valued measures, correlators, currents, and detector outcomes form the observable map \(R_Q\). Preparation, evolution, control, and measurement order form the protocol \(P_Q\). The realization \(A_Q\) then fixes the particles, fields, material, geometry, parameter values, drives, and apparatus.
\begin{centeredalign}
I_{\mathrm{op},Q} = ((\Omega_Q,\Xi_Q);C_Q,R_Q,P_Q),\\
I_{\mathrm{real},Q} = (I_{\mathrm{op},Q};A_Q),\\
I_{\mathrm{real},Q} \longmapsto p(y),\ \langle O\rangle,\ S(\omega),\ J,\ \text{or another stated consequence}.
\end{centeredalign}
Changing the state space can destroy self-adjointness, locality, or normalization. Changing the protocol can change the channel itself. The closure and observable must therefore be derived again whenever one of these physical roles changes.
\section{Transformations And Compatibility}
\begin{centeredalign}
I_i=((\Omega_i,\Xi_i);C_i,R_i,P_i)\xrightarrow{\tau}I_j=((\Omega_j,\Xi_j);C_j,R_j,P_j).
\end{centeredalign}
The transformation \(\tau\) relates two descriptions or realizations. It may change representation, complete a partial model, project onto a reduced description, transfer a law to another state space, compose operations, or take a controlled limit. Its physical content is the relation retained between the two constructions: an amplitude, expectation value, operator algebra, balance law, probability distribution, or approximation with a stated error.
The retained relation is tested by maps that carry states and outputs between the two descriptions.
\begin{centeredalign}
\alpha : \Xi_A\longrightarrow\Xi_B,\\
\beta : \Omega_A(\Xi_A)\longrightarrow\Omega_B(\Xi_B),\\
\Delta_{\alpha,\beta}=\Omega_B\alpha-\beta\Omega_A.
\end{centeredalign}
The map \(\alpha\) carries physical states from the source description to the target, while \(\beta\) carries the corresponding outputs. Exact compatibility means that acting first and then mapping gives the same result as mapping first and acting in the target. The residual \(\Delta_{\alpha,\beta}\) records the part of the source dynamics that the proposed target description cannot represent.
\[\begin{cases}\mathrm{transfer},&\Delta_{\alpha,\beta}=0,\\\mathrm{reject},&\Delta_{\alpha,\beta}\ne0\ \text{without stable closure},\\\mathrm{promote}(\Delta),&\Delta_{\alpha,\beta}\ne0\ \text{and closes}\end{cases}\]
A vanishing residual, together with agreement of the stated observables, identifies another realization of the same mechanism. A residual with no stable form rejects the proposed transfer. When \(\Delta_{\alpha,\beta}\) recurs, obeys its own closure relation, and changes a measurable consequence, it becomes a candidate correction operator, interaction, boundary term, or hidden state variable.

\chapter*{When External Conditions Become Quantum Physics}
\addcontentsline{toc}{chapter}{When External Conditions Become Quantum Physics}
Quantum theory changes when a quantity once treated as fixed becomes part of the physical state or its dynamics. A prescribed electromagnetic potential becomes a quantum field. A classical measuring device becomes an interacting quantum subsystem. A passive environment becomes a source of memory when its correlations influence later motion. Quantum gravity asks the same question of geometry itself.
In each case, the variables retained by the theory must determine future observable probabilities, and equivalent physical transformations must compose to the same result. We call this condition predictive closure.
\begin{claimbox}
At a chosen resolution, a physical theory is closed when its declared state fixes future observable probabilities for a specified subsequent experiment. Equivalent descriptions must give the same predictions under composition. A failure identifies a dependence that must be represented or controlled, such as an environmental correlation, a boundary condition or an omitted interaction.
\[q(h_1)=q(h_2)\Longrightarrow p(y,t\mid h_1)=p(y,t\mid h_2),\qquad T_{\gamma_1}=T_{\gamma_2}\ \text{for physically equivalent paths}\]
\end{claimbox}
The first relation compares two histories that produce the same declared state. Their later probabilities must agree. If they do not, the state has discarded a physically active correlation; an internal coordinate or memory kernel restores the missing dependence. The second relation compares two physically equivalent paths. Their transformations must agree when applied to the same state. Physical transport along different paths can instead give different states, as in an interference experiment with geometric phase; that difference is a prediction of the connection, not a failure of consistency.
Predictive closure also determines the boundary between a mechanism and its realization. A parameter remains external while it selects a member of a fixed theory. It enters the mechanism when it changes the state space, operator domain, dynamical map, admissibility condition, observable, or operation order.
\[\begin{aligned}a\in A\quad &\text{while}\quad (\Xi,\Omega,C,R,P)_a=(\Xi,\Omega,C,R,P),\\a\longrightarrow X\in\{\Xi,\Omega,C,R,P\}\quad &\text{when the physical object }X\text{ changes}.\end{aligned}\]
For a particle confined to an interval, both its length and the conditions at its ends affect the energy spectrum. A family of intervals can be mapped onto one reference interval by rescaling the coordinate, with the length then appearing in the kinetic coefficient. Dirichlet and Neumann conditions still define different operator domains on that reference interval. The realization therefore instantiates a specified family of operators; it is not a layer whose every change leaves the predictions fixed.
In the nested construction of Chapter~\ref{chap:physical-identity}, the dynamics and the observable consequence read
\[q\in\Xi,\qquad \dot q=\Omega[q;C,P],\qquad y=R(q)\]
where a physical state \(q\) contains the information required to determine later observables at the chosen resolution.
Each clause prevents a specific ambiguity. Without the state space, the operator has no domain. Without closure, its solutions or probabilities need not be admissible. Without an observable, the formal evolution has no predicted measurement. Without a protocol, noncommuting operations have no defined order.
For a quantum state, this identity reduces to a state space, an evolution map, and a positive measurement law.
\begin{centeredalign}
\Xi_Q = (\mathcal H,\mathcal D,\mathcal S),\qquad \rho_0\in\mathcal S(\mathcal H),\\
\rho_P = \mathcal E_P(\rho_0),\qquad \mathcal E_P\ \text{completely positive and trace preserving},\\
p(y\mid P) = \operatorname{Tr}(E_y\rho_P),\qquad E_y\ge0,\quad \sum_yE_y=I.
\end{centeredalign}
The Hilbert space, state class, and operator domain define the admissible states. The channel describes their evolution under protocol P. The positive operators E\_y define the possible outcomes, and the trace rule assigns their probabilities. Unitary dynamics, open-system evolution, measurement, and feedback differ in which physical structures enter these three lines.
Here \(\mathcal E_P\) describes a deterministic protocol, including measurement outcomes averaged or retained in a classical register. A selected measurement outcome is represented by a completely positive, trace-nonincreasing map; its trace gives the probability of that outcome. A reduced channel also presupposes a specified environment preparation compatible with the chosen input states.
These relations order the following 146 topics by the physical role each one fixes or changes. Entanglement belongs to composite state structure. Commutators belong to observable algebra. Bell experiments join the two through local measurements. Gauge constraints select physical states, while renormalization changes the effective law with scale. This ordering exposes transitions between theories that conventional subject headings place in separate fields.

\chapter{When External Structure Becomes Dynamical}
\begin{claimbox}
A background field, geometry, environment, boundary, constraint, or control sequence changes the theory when it alters the state space, dynamics, operator domain, observable, or operation order.
\end{claimbox}
\section{Backgrounds, Geometry, And Environment}
A background field supplies prescribed coefficients to a Hamiltonian. Quantizing that field enlarges the Hilbert space and gives the field its own conjugate variables, fluctuations, and quanta. Quantum electrodynamics makes this change for the electromagnetic field. Quantum gravity asks whether distances and causal relations must likewise participate in the quantum state.
The same change occurs in reduced dynamics. Tracing over an environment produces closed Markovian dynamics only when discarded correlations do not influence the future. When they return, the reduced state is incomplete. Retaining environmental coordinates or a memory kernel restores predictive closure. Decoherence, non-Markovian dynamics, and measurement back-action therefore differ by which environmental correlations remain physically active.
\section{Domains, Constraints, And Controlled Dynamics}
A boundary condition enters the physical law when it defines the operator domain. The same differential expression can then possess different self-adjoint extensions and different spectra. Confinement, tunnelling, scattering, and topological edge modes follow from this dependence on the domain.
Constraints alter the state space rather than the operator domain. Gauge redundancy and exchange symmetry begin as restrictions on allowed descriptions. Imposing them selects the physical Hilbert-space sectors. Gauss constraints remove gauge-equivalent vectors, while Bose or Fermi symmetry fixes the admissible many-particle states. Charges, particle statistics, and observable algebras follow from the surviving sectors.
A sequence of gates, measurements, and conditional controls becomes a single physical map after composition. Its order can change the resulting unitary or channel. In quantum computing, feedback, error correction, and Floquet control, the experimental sequence is therefore part of the dynamics implemented during the experiment.
\section{One Principle Across The Frontiers}
Each enlargement repairs the same failure. A smaller theory omitted information required for future probabilities or for consistent composition. The minimal enlargement identifies the missing physical object: a field, an environmental coordinate, an operator domain, a constrained state sector, or a composed channel. Quantum theories can therefore be compared by the physical role of their variables rather than by the historical field in which those variables were introduced.

\chapter{Global Composition: Curvature, Frustration, And Memory}
\begin{claimbox}
Parallel transport, bond-energy minimization and reduced evolution ask different questions about the information carried by a path. Their equations distinguish a geometric phase, incompatible local preferences and dependence on discarded correlations.
\end{claimbox}
\section{From Local Transformations To A Global State}
Let \(T_{ij}\) carry a state, basis, or constraint from one local description to the next. A path is an ordered composition of these maps. It can describe transport through parameter space, constraints around an interaction loop, or a sequence of interventions followed by relaxation. For a closed path \(\gamma\), the composite map is
\begin{equation*}
\mathfrak H_\gamma=T_{0n}T_{n,n-1}\cdots T_{21}T_{10}.
\end{equation*}
For pure changes of basis, returning to the initial basis gives \(\mathfrak H_\gamma=I\) when the transition maps are consistent. Parallel transport between physical points is a different operation and can have nontrivial holonomy in a fully specified theory. A physical protocol need not return the state when its external controls close a loop. The corresponding completeness test compares histories that end at the same retained state. If their later observables differ, the retained endpoint variables have omitted physically active path information.
\section{Curvature Records Transport Around A Loop}
A quantum state transported around a loop can return with a phase or an internal rotation. In gauge theory the path-ordered exponential of the connection gives the loop operator; for a small loop its departure from the identity is set by the field strength. Berry curvature gives the corresponding relation in parameter space.
\begin{centeredalign}
W(\gamma)=\operatorname{Tr}\,\mathcal P\exp\!\left(-ig\oint_\gamma A_\mu\,dx^\mu\right),\\
\mathfrak H_{\mu\nu}(\ell)=I-igF_{\mu\nu}\ell^2+O(\ell^3).
\end{centeredalign}
Here \(g\) is the gauge coupling and \(\ell\) the side of a positively oriented square in the \(\mu\nu\) plane. The convention is \(D_\mu=\partial_\mu+igA_\mu\), with \([D_\mu,D_\nu]=igF_{\mu\nu}\). The connection specifies how neighboring descriptions are compared; the curvature measures the infinitesimal failure of those comparisons to be path independent. The traced loop is invariant under a change of local gauge.
\section{Frustration Records Incompatible Local Preferences}
In a frustrated magnet, each bond can impose a well-defined local preference while no spin configuration satisfies every bond around a loop. For an Ising energy \(H=-\sum_{\langle ij\rangle}J_{ij}s_i s_j\), with \(s_i=\pm1\) and nonzero bonds, minimizing each bond separately requires \(s_i s_j=\operatorname{sign}J_{ij}\). Multiplying those requirements around a loop gives \(+1\) on the left. A negative product of bond signs therefore makes simultaneous minimization impossible:
\begin{equation*}
\prod_{\langle ij\rangle\in\gamma}\operatorname{sign}J_{ij}=-1.
\end{equation*}
The mismatch survives every local reassignment of spin signs and is therefore a property of the loop. In quantum magnets the same logic is carried by loop operators, flux sectors, or non-commuting bond terms. The signed bonds already determine this obstruction. A loop label summarizes their joint incompatibility; adding a label does not remove the unsatisfied bond or restore an unfrustrated ground state.
\section{Memory Records Hidden Paths Through State Space}
A reduced quantum state is complete only if equal present states give equal future observable statistics. Let \(h_1\) and \(h_2\) be two histories that produce the same retained state \(q\). The state is incomplete when hidden system--environment correlations make a later observable distinguish the histories:
\begin{centeredalign}
q(h_1)=q(h_2),\qquad R\Phi_t(h_1)\ne R\Phi_t(h_2),\\
V_{t+s}\ne V_tV_s.
\end{centeredalign}
For stationary preparations the second relation tests a time-homogeneous semigroup description. A time-dependent Markovian generator can also violate that one-parameter identity, so history dependence is decided by the first comparison under identical subsequent controls. Eliminating environmental coordinates can produce both a memory kernel and a term depending on their initial preparation. Both contributions must be included when assessing whether the retained state determines later observations. For order memory, two write operations followed by the same relaxation are distinguishable precisely when their ordered products leave different retained states.
\begin{equation*}
R\mathcal R_tW_BW_A(q_0)\ne R\mathcal R_tW_AW_B(q_0).
\end{equation*}
\section{The Shared Physical Content}
These calculations compare what local information determines about a larger construction. A connection determines parallel transport, and its curvature makes that transport depend on the path. Signed Ising bonds determine local energy preferences whose product can prevent simultaneous satisfaction. A reduced density operator can discard correlations that affect later evolution. The relevant physical information is respectively the connection along the path, the arrangement of bond signs, and the system--environment state. The comparison motivates searching for such missing dependencies when assembling mechanisms. It does not identify the three mathematical objects: curvature can occur in a complete theory, frustration can be fully specified by its bonds, and reduced memory need not be a geometric holonomy. Their equations establish the relation in each physical setting.

\chapter{Equivalence Across Quantum Descriptions}
The following relations compare descriptions through their dynamics and observable consequences. Their domain and measurement conditions determine whether the proposed correspondence holds.
\begin{claimbox}
Two quantum descriptions represent the same physical mechanism when a map between their states preserves the dynamics and the stated observable consequences.
\end{claimbox}
Quantum theory uses several forms of this equivalence. A change of representation carries states and observables together. A change of carrier embeds the same operator algebra in another Hilbert space. A reduced description projects onto fewer degrees of freedom, and a measurement or boundary condition completes the relation by fixing its observable content. In each case the equivalence is decided by the quantity that remains unchanged.
\section{One operator, two roles: generator and observable}
This relation connects Hamiltonian (quantum mechanics), Observable, Spectral theory.
The exponential of the Hamiltonian generates time evolution, whereas its spectral measure defines energy probabilities.
The retained physical relation is the self-adjoint operator and its domain.
\begin{centeredalign}
U(t)=e^{-iHt/\hbar}\\
H=\int_{\sigma(H)}E\,dP_H(E)
\end{centeredalign}
Equivalence requires a common operator domain and self-adjointness; the generated dynamics must be unitary and the spectral measure must reproduce the same energy statistics.
\section{Schrödinger, Heisenberg, and path-integral representations}
This relation connects Schrödinger picture, Heisenberg picture, Path integral formulation.
Time dependence may be carried by states or operators, and the same propagator may be represented by an action-weighted path integral.
The retained physical relation is transition amplitudes and expectation values.
\begin{centeredalign}
A_H(t)=U(t)^\dagger A_SU(t)\\
\langle q_f|U(t_f-t_i)|q_i\rangle=\int\mathcal Dq\,e^{iS[q]/\hbar}
\end{centeredalign}
Equivalence requires compatible domains, measures, boundary conditions, and regularization; the formulations coincide only where they yield the same amplitudes or correlators.
\section{Carrier factorization defines locality}
This relation connects Quantum entanglement, Quantum information, Quantum field theory.
Changing the tensor-product or algebraic decomposition changes which observables are local and which states are entangled.
The retained physical relation is the global state and algebraic predictions.
\begin{centeredalign}
\mathcal H=\mathcal H_A\otimes\mathcal H_B\\
\rho_A=\operatorname{Tr}_B\rho_{AB}
\end{centeredalign}
Equivalence requires a fixed subsystem algebra or tensor factorization because entanglement and locality are not invariant under arbitrary refactorization.
\section{Constraints define physical observables}
This relation connects Gauge theory, Quantum gravity, Measurement in quantum mechanics.
Gauge or diffeomorphism constraints select the physical state space before outcome operators acquire observable meaning.
The retained physical relation is predictions on the physical state space.
\begin{centeredalign}
\widehat{\mathcal C}_a|\psi_{\rm phys}\rangle=0\\
[\widehat O_{\rm phys},\widehat{\mathcal C}_a]|\psi_{\rm phys}\rangle=0
\end{centeredalign}
Physical effects must descend to the constrained or quotient state space; gauge-dependent quantities do not define physical records.
\section{Boundary conditions are spectral control variables}
This relation connects Particle in a box, Quantum tunnelling, Scattering, Spectral theory.
The operator domain, boundary conditions, and asymptotic channels determine the spectrum and transmission amplitudes.
The retained physical relation is the differential operator family and probability conservation.
\begin{centeredalign}
H_D=-\frac{\hbar^2}{2m}\Delta_D+V\\
S:\mathcal H_{\rm in}\to\mathcal H_{\rm out}
\end{centeredalign}
Equivalence requires the same operator domain and conserved flux; identical differential expressions on different domains can have different spectra.
\section{Measurement, channels, and error correction share an instrument calculus}
This relation connects Quantum channel, Measurement in quantum mechanics, Quantum error correction.
A quantum instrument produces both an outcome probability and a conditional state; error correction uses the recorded outcome to select a recovery channel.
The retained physical relation is complete positivity and total probability.
\begin{centeredalign}
\mathcal I_i(\rho)=\sum_\alpha K_{i\alpha}\rho K_{i\alpha}^\dagger\\
\mathcal E=\sum_i\mathcal I_i\\
\mathcal R_i\circ\mathcal I_i
\end{centeredalign}
Equivalence requires complete positivity, normalization, and recovery fidelity with a reference system; a state-update rule alone is insufficient.
\section{Duality and simulation require an intertwining map}
This relation connects Quantum simulator, AdS/CFT correspondence, Quantum error correction.
An encoding represents one state space in another only when it intertwines the relevant dynamics and observables.
The retained physical relation is a selected operator algebra and its correlators.
\begin{centeredalign}
VH_{\rm target}\simeq H_{\rm carrier}V\\
VO_{\rm target}\simeq O_{\rm carrier}V
\end{centeredalign}
Equivalence requires agreement of a generating set of observables or correlators within a stated approximation regime and error bound; state overlap alone is insufficient.

\chapter{Theory Extension By Transfer And Obstruction}
\begin{claimbox}
A physical relation can survive in another state space, or its failure can identify the term missing from the target theory. Exact compatibility and structured incompatibility are therefore two outcomes of the same transformation.
\end{claimbox}
\section{Incomplete Mechanisms}
A law can be stated before the physical degrees of freedom on which it acts are known. Conversely, a state space can be specified before its governing operation is fixed. These one-sided descriptions occur at the boundary between an abstract theory and a physical realization.
\begin{centeredalign}
(\Omega,0_{\Xi})+(0_{\Omega},\Xi')\longrightarrow(\Omega,\Xi')\\
\longrightarrow(\Omega,\Xi';\widehat C,\widehat R,\widehat P)\\
\longrightarrow I_{\mathrm{real}}'.
\end{centeredalign}
The combined theory exists only if the operation has a domain on the target state space. Its closure must admit physical states, and its observable must produce normalized probabilities. These requirements determine the additional terms introduced by the new realization.
\section{Exact Transfer}
\[\Delta_{\alpha,\beta}=\Omega_B\alpha-\beta\Omega_A\]
The maps alpha and beta carry source states and outputs into the target description. A vanishing residual means that evolution followed by translation gives the same state as translation followed by target evolution. Agreement of the corresponding observables then establishes another physical realization of the retained mechanism.
Schrödinger, Heisenberg, and path-integral descriptions provide standard examples. Their variables and intermediate objects differ, while amplitudes and expectation values agree on the common domain. Quantum simulation uses the same structure when an encoded target algebra reproduces selected correlators in another physical device.
\section{Structured Failure}
A nonzero residual separates two physical possibilities. An irregular residual reflects an incompatible comparison. A residual that recurs across states or scales and obeys a closure relation has acquired the form of a physical operator.
\begin{equation*}
\Delta\longmapsto\Omega',\ \Xi',\ C',\ R',\ P',\ \text{or }A'.
\end{equation*}
Promotion into Omega introduces a correction law or interaction. Promotion into Xi adds a hidden degree of freedom or state sector. Promotion into C changes the admissible domain. Promotion into P makes operation order part of the dynamics. The enlarged theory is fixed by the smallest promotion that restores predictive closure.
\section{Observable Consequences}
The promoted term becomes physical when it changes an observable that was not used to define it. Removing the obstructing coupling must remove that consequence. Curvature, memory, edge flux, and effective interactions are familiar outcomes of this logic: each retains information that the smaller state or transformation law had discarded.

\part{State Spaces And Dynamical Laws}

\chapter{State space, domain, and representation}

\begin{claimbox}

The Hilbert or Fock space, operator domain, basis, gauge sector, and subsystem decomposition determine which states and operations exist. Changing this structure can change the theory before any Hamiltonian parameter is varied.

\end{claimbox}

A Hilbert space and an operator domain determine which states and operations exist. Gauge sectors, bases, and subsystem decompositions refine that physical domain.

\section*{Topics}

\begin{longtable}{p{0.92\linewidth}}
\toprule
Page \\
\midrule
\endhead
\phantomsection\label{page:hilbert-space}Hilbert space \\
\phantomsection\label{page:quantum-mechanics}Quantum mechanics \\
\phantomsection\label{page:old-quantum-theory}Old quantum theory \\
\phantomsection\label{page:fourier-transform}Fourier transform \\
\bottomrule
\end{longtable}

\clearpage
\section{Hilbert space}
\label{topic:hilbert-space}

\subsection*{Mechanism}
Hilbert space is the admissible state carrier of quantum theory: it supplies the space in which states, operators, bases, spectra, and probabilities become legally defined.

Hilbert space is not merely a container for wave functions. Its inner product defines amplitudes and orthogonality, while the domains of unbounded operators decide whether expressions for position, momentum, and energy are mathematically and physically admissible.

Hilbert space is not physical space and not a geometric background in this book. It is the legal carrier of quantum identity. A state is a vector or density operator on it; the inner product gives amplitudes and norms; observables are self-adjoint operators on it; spectral projectors define possible answers; and unitary evolution preserves norm and probability. Hilbert space is therefore central because it binds state, probability, operator spectrum, and identity preservation into one formal carrier.

\subsection*{Physical Construction}
The state carrier is a complex Hilbert space, or a density-operator state space built on it. The calculation involves self-adjoint observables, unitary maps, spectral projectors, and domain-restricted generators defined on the carrier. Inner-product structure, normalization, positivity for density states, and operator-domain conditions make states and observables legal. The calculated quantities include Born probabilities, spectral projectors, expectation values, and preserved norms.

\subsection*{Topic Equations}
\begin{centeredalign}
\ket{\psi}\in\mathcal H,\qquad \langle\psi|\psi\rangle=1\\
\rho\in\mathcal S(\mathcal H),\qquad \rho\ge0,\quad \operatorname{Tr}\rho=1\\
A=A^\dagger,\qquad A=\int_{\sigma(A)}\lambda\,dE_A(\lambda)\\
\Pr(\Delta\mid \rho,A)=\operatorname{Tr}\!\left(\rho E_A(\Delta)\right)\\
\rho_t=U(t)\rho U(t)^\dagger,\qquad U^\dagger U=I
\end{centeredalign}

\subsection*{Physical Meaning}
The linear structure permits superposition, the inner product converts pairs of states into amplitudes, and completeness guarantees that convergent sequences of approximations remain inside the space. Operator domains must be carried with the operators; the same differential formula on another domain can describe a different physical system.

A qubit lives in a two-dimensional complex space, whereas a particle on a line is described in a space of square-integrable functions. Both obey the same Hilbert-space logic, but their operators, spectra, and boundary conditions are different.

After the arena is fixed, a quantum state selects one preparation or statistical ensemble within it.

\subsection*{Invariance And Realization}
The rule connecting prepared states, observables, and spectral probability measures across wave, matrix, path-integral, circuit, or field notation. The operator-to-spectrum relation: admissible observations are represented through eigenvalues, projections, modes, or outcome channels. The dependence of admissible observable on measurement context or boundary condition. The non-commuting compatibility structure, which survives changes of representation.

The name of the carrier: particle, wave, field, qubit, or excitation. Where time dependence is represented: on the state, on the operator, or in a path weight. The coordinate system, basis, or geometric picture used to display the same relation. The physical implementation of detector, boundary, preparation, or observable.

\subsection*{Discriminating Consequences}
Unitary changes of basis preserve Born probabilities; if probabilities change, the page has changed the physical context rather than only the representation. The operator domain and normalization conditions determine which questions are legal on the selected Hilbert space. Basis names are bookkeeping; inner products, spectra, and probabilities are the physical content that must survive the rewrite.




\clearpage
\section{Quantum mechanics}
\label{topic:quantum-mechanics}

\subsection*{Mechanism}
Quantum mechanics is the baseline constructor: states live in Hilbert space, physical questions are represented by operators, and probabilities are assigned to spectral projectors.

A quantum equation has no fixed meaning until its state space, inner product, representation, and operator domains have been specified. These choices decide which states are admissible and which apparent changes are only changes of coordinates.

The page supplies the general quantum assembly. A preparation gives a state vector or density operator. A self-adjoint observable or measurement operator family gives the possible outcome channels. The Born or trace rule assigns probabilities, while Hamiltonian evolution transports the state between preparation and observable.

\subsection*{Physical Construction}
The state carrier is a Hilbert, Fock, or function space together with the operator domains and representation used in the calculation. The calculation involves a unitary or isometric change of basis, Fourier transform, coordinate map, or representation equivalence. Inner products, domains, normalization, and completeness relations must be preserved by a purely representational change. The calculated quantities include transition amplitudes, expectation values, spectra, and probabilities that remain invariant under an admissible representation change.

\subsection*{Topic Equations}
\begin{centeredalign}
\rho\ge 0,\qquad \operatorname{Tr}\rho=1\\
A=\sum_a aP_a,\qquad p(a)=\operatorname{Tr}(\rho P_a)\\
\rho(t)=U(t)\rho(0)U(t)^\dagger,\qquad U(t)=e^{-iHt/\hbar}\\
[A,B]\ne0\quad\Rightarrow\quad \text{no generic common sharp eigenbasis}
\end{centeredalign}

\subsection*{Physical Meaning}
A unitary or isometric change of representation carries the state and operator together. Amplitudes, expectation values, and spectra agree. If they do not, the physical model has changed rather than merely its notation.

A quantum state belongs to this state space, and every Hamiltonian and observable must act on its stated domain. These domain relations determine whether the resulting amplitudes and probabilities are defined.

\subsection*{Invariance And Realization}
The relation between prepared states, observables, and spectral probability measures. The use of eigenvalues, projectors, modes, or outcome channels to represent admissible observations. The dependence of the observable on basis, domain, potential, preparation, or measurement context. The commutator structure that limits which observables can be jointly diagonalized.

The physical carrier: particle, wave, field mode, spin, qubit, detector, or excitation. The representation: wave mechanics, matrix mechanics, density matrices, path integrals, circuits, or fields. Where time dependence is placed: on the state, on the operator, in a propagator, or in a path weight. The implementation of preparation, boundary condition, detector, or outcome channel.

\subsection*{Discriminating Consequences}
A transfer target provides a state space, a transformation law, and a spectral or categorical observable, with one compatibility relation experimentally unresolved. A useful validation varies the basis, domain, or measurement context and measures whether the allowed observable changes while the underlying transformation law remains identifiable. A stronger validation contains two candidate observables whose predicted commutator controls joint resolvability.




\clearpage
\section{Old quantum theory}
\label{topic:old-quantum-theory}

\subsection*{Mechanism}
Old quantum theory is the semiclassical quantization precursor to modern quantum mechanics: classical periodic motion is retained, while only selected actions and transition frequencies are admitted.

A quantum equation has no fixed meaning until its state space, inner product, representation, and operator domains have been specified. These choices decide which states are admissible and which apparent changes are only changes of coordinates.

Old quantum theory predates the Hilbert-space formalism. Its carrier is classical phase space, its motion follows Hamiltonian trajectories, and its quantum restriction is imposed on periodic action integrals. The resulting discrete energies and transition frequencies anticipated quantum spectra but did not provide a general theory of states, observables, or noncommuting transformations.

\subsection*{Physical Construction}
The state carrier is classical phase space, especially periodic or multiply periodic Hamiltonian trajectories described by action variables. The calculation involves Hamiltonian flow together with Bohr--Sommerfeld action quantization; no general Hilbert-space operator calculus is assumed. Only trajectories satisfying quantized action conditions are retained, with correspondence to classical frequencies required at large quantum number. The calculated quantities include discrete energies, transition frequencies, and the spectral regularities inferred from them.

\subsection*{Topic Equations}
\begin{centeredalign}
J_k=\oint p_k\,dq_k=n_k h\\
E_m-E_n=h\nu_{mn}\\
\frac{\partial E}{\partial J_k}=\nu_k^{\mathrm{cl}}\quad\text{in the correspondence regime}
\end{centeredalign}

\subsection*{Physical Meaning}
A unitary or isometric change of representation carries the state and operator together. Amplitudes, expectation values, and spectra agree. If they do not, the physical model has changed rather than merely its notation.

A quantum state belongs to this state space, and every Hamiltonian and observable must act on its stated domain. These domain relations determine whether the resulting amplitudes and probabilities are defined.

\subsection*{Invariance And Realization}
Hamiltonian trajectories and their action variables provide the classical carrier of the construction. Action quantization selects a discrete subset of otherwise continuous classical motions. Energy differences determine transition frequencies, with correspondence to classical orbital frequencies at large quantum number.

The orbit family, potential, number of action variables, and applicable quantization condition depend on the physical system. Modern semiclassical theory adds phase and turning-point corrections absent from the earliest quantization rules. The construction ceases to be adequate for generic nonintegrable motion, many-electron spectra, and intrinsically noncommuting questions.

\subsection*{Discriminating Consequences}
The theory recovers the hydrogenic level and line spectrum in the regime where action quantization applies. A comparison of the semiclassical prediction with the corresponding Schrödinger spectrum and identify its controlled limit. Treat failures outside that limit as evidence for the modern state-and-operator formalism, rather than as adjustable orbit parameters.




\clearpage

\section{Fourier transform}

\label{topic:fourier-transform}

A free particle propagates most simply in momentum space, where each momentum component accumulates its own phase. A localized detector, however, asks for a position probability. The Fourier transform connects these descriptions while preserving the physical state. On the real line, for a square-integrable wave function \(\psi\) and momentum \(p\), take

\begin{centeredalign}
\widetilde\psi(p)=\frac{1}{\sqrt{2\pi\hbar}}\int_{\mathbb R}
e^{-ipx/\hbar}\psi(x)\,dx,\qquad
\int|\widetilde\psi(p)|^2dp=\int|\psi(x)|^2dx.
\end{centeredalign}

For wave functions on a common suitable domain, integration by parts carries the position-space momentum operator \(-i\hbar\partial_x\) into multiplication by \(p\). Multiplication by position becomes \(i\hbar\partial_p\). The kinetic equation and the detector must both be transformed:

\begin{centeredalign}
\mathcal F(-i\hbar\partial_x)\mathcal F^{-1}=p,\qquad
\mathcal F x\mathcal F^{-1}=i\hbar\partial_p,\qquad
\widetilde H=\frac{p^2}{2m}+V(i\hbar\partial_p).
\end{centeredalign}

The last expression is immediately useful for polynomial potentials; for a general potential, multiplication in position space is represented by an integral kernel in momentum space. A local interaction can therefore become nonlocal in the new coordinates without changing the theory. Similarity of displayed formulas is neither necessary nor sufficient for physical equivalence.

For a normalized Gaussian packet with initial position variance \(\sigma_x^2\), zero position-momentum covariance and minimum uncertainty, the free evolution multiplies each momentum amplitude by \(\exp[-ip^2t/(2m\hbar)]\). Transforming back gives

\begin{centeredalign}
\operatorname{Var}x(t)=\sigma_x^2+
\frac{\hbar^2t^2}{4m^2\sigma_x^2}.
\end{centeredalign}

The packet broadens because its different momentum components separate during propagation. The momentum distribution itself remains unchanged. This provides a direct comparison between descriptions: one representation makes the conservation of momentum probabilities explicit, while the other displays the spatial spreading produced by those same components.

Boundaries change this correspondence. Periodic functions on a finite interval have a discrete Fourier series; Dirichlet eigenfunctions lead instead to a sine expansion. Surface terms must vanish under the actual boundary conditions before the integration-by-parts identity defines the required operator map. The operation of Fourier transformation is transferable, but the domain and measure determine which transform and which spectrum describe the physical system.

\subsection*{References For The Physical Derivation}
\begin{itemize}
\item \href{https://arxiv.org/abs/1508.06951}{V. Moretti, Mathematical Foundations of Quantum Mechanics: An Advanced Short Course; states, operators and symmetry.}
\end{itemize}

\subsection*{Relations In The Original Papers}

\href{https://arxiv.org/html/1508.06951#S2.EGx7}{arXiv:1508.06951, S2.EGx7}. The Fourier transformation expresses the position wavefunction in wavevector coordinates. This source uses wavevector k; the chapter uses momentum p = hbar k with the corresponding normalization.

\chapter{Quantum states and subsystem structure}

\begin{claimbox}

A state is the probability-bearing element of the selected Hilbert space or its density-operator state space. Wave functions, density matrices, superpositions, and coherent states are different representations of this predictive carrier.

\end{claimbox}

The state must retain every coordinate needed to determine later probabilities. Entanglement enters here because it is a property of a joint state relative to a stated subsystem decomposition.

\section*{Topics}

\begin{longtable}{p{0.92\linewidth}}
\toprule
Page \\
\midrule
\endhead
\phantomsection\label{page:density-matrix}Density matrix \\
\phantomsection\label{page:quantum-decoherence}Quantum decoherence \\
\phantomsection\label{page:wave-function}Wave function \\
\phantomsection\label{page:quantum-state}Quantum state \\
\phantomsection\label{page:two-state-quantum-system}Two-state quantum system \\
\phantomsection\label{page:qubit}Qubit \\
\phantomsection\label{page:quantum-entanglement}Quantum entanglement \\
\bottomrule
\end{longtable}

\clearpage
\section{Density matrix}
\label{topic:density-matrix}

\subsection*{Mechanism}
Density matrix is the mixed-state constructor: it keeps probabilistic preparation, entanglement with unobserved degrees of freedom, and partial information in the same state formalism.

The density matrix is required whenever the preparation is statistical, part of an entangled system is ignored, or environmental coupling makes a state-vector description of the subsystem incomplete.

Density matrices generalize pure states without changing the state-to-spectrum probability assignment rule. They are the correct carrier when the preparation is statistical, when a subsystem is traced out, or when decoherence is being described.

\subsection*{Physical Construction}
The state carrier is the mathematical state object: vector, wavefunction, density operator, coherent state, field state, or register. Operators, maps, and observables become meaningful only after this carrier and its domain have been fixed. Normalization, positivity, inner product, representation, tensor factorization, or superselection conditions define legal states. The calculated quantities include probability distributions obtained by applying the appropriate observables or measurement maps to the carrier.

\subsection*{Topic Equations}
\begin{centeredalign}
\rho=\sum_a p_a\ket{\psi_a}\bra{\psi_a},\qquad p_a\ge0,\quad \sum_a p_a=1\\
\rho\ge0,\qquad \operatorname{Tr}\rho=1\\
p(i)=\operatorname{Tr}(\rho P_i)\\
\rho_A=\operatorname{Tr}_B(\rho_{AB})
\end{centeredalign}

\subsection*{Physical Meaning}
Diagonal elements in a selected basis give populations and off-diagonal elements carry coherence in that basis. Positivity and unit trace make the Born probabilities legal. Partial tracing produces the state of a subsystem without assigning a pure vector to it.

Either randomly preparing two spin directions or discarding one member of an entangled pair can produce a mixed state. The density matrix records the observable statistics, even though the physical origins of the mixture differ.

With the state representation established, unitary and non-unitary maps specify how it changes.

\subsection*{Invariance And Realization}
The rule connecting prepared states, observables, and spectral probability measures across wave, matrix, path-integral, circuit, or field notation. The operator-to-spectrum relation: admissible observations are represented through eigenvalues, projections, modes, or outcome channels. The dependence of admissible observable on measurement context or boundary condition. The non-commuting compatibility structure, which survives changes of representation.

The name of the carrier: particle, wave, field, qubit, or excitation. Where time dependence is represented: on the state, on the operator, or in a path weight. The coordinate system, basis, or geometric picture used to display the same relation. The physical implementation of detector, boundary, preparation, or observable.

\subsection*{Discriminating Consequences}
A usable state gives normalized probabilities for every complete observable attached to the selected Hilbert space. Vector, wave-function, density-matrix, and reduced-state forms can describe the same preparation when connected by the appropriate representation map. Physical state changes preserve positivity and trace, or norm in the pure closed-system limit.




\clearpage
\section{Quantum decoherence}
\label{topic:quantum-decoherence}

\subsection*{Mechanism}
Quantum decoherence is the loss of observable phase coherence in a subsystem when information about alternative amplitudes becomes encoded in environmental correlations.

Interference is lost when alternative system amplitudes become correlated with environmental states that can, even in principle, distinguish them. The missing coherence has then moved into correlations of the larger state rather than being erased by the subsystem equation alone.

The joint system and environment may evolve unitarily while the reduced system loses interference. Interaction correlates different system alternatives with distinguishable environmental states; tracing over the environment then suppresses off-diagonal terms in the selected pointer basis. Decoherence does not select one outcome, and it is not identical to memory. Markovian decoherence is possible when the reduced state still determines its future. Memory begins when discarded correlations return and two preparations with the same reduced state acquire different later statistics.

\subsection*{Physical Construction}
The state carrier is a joint system-environment state together with the reduced density operator accessible to the observer. The calculation involves system-environment unitary evolution followed by partial trace, or the corresponding reduced dynamical map. The joint state remains normalized and positive; a Markovian reduction additionally requires future maps to be fixed by the present reduced state. The calculated quantities include off-diagonal coherence, interference visibility, purity, and history-dependent response.

\subsection*{Topic Equations}
\begin{centeredalign}
\rho_S(t)=\operatorname{Tr}_E[U_{SE}(t)\rho_{SE}(0)U_{SE}^{\dagger}(t)]\\
\frac{\ket0\ket{E_0}+\ket1\ket{E_1}}{\sqrt2},\qquad (\rho_S)_{01}\propto\langle E_1|E_0\rangle\\
\dot\rho_S=-\frac{i}{\hbar}[H_S,\rho_S]+\sum_k\gamma_k\mathcal D[L_k]\rho_S\\
V_{t+s}=V_tV_s\quad\text{only in the time-homogeneous Markov limit}
\end{centeredalign}

\subsection*{Physical Meaning}
Unitary system-environment evolution can suppress the off-diagonal elements of the reduced density matrix after the environment is traced out. The overlap of the corresponding environmental states fixes the remaining visibility. Decoherence can be Markovian; memory requires the stronger condition that hidden correlations later alter evolution from the same reduced state.

In a two-path interferometer, scattering one distinguishable environmental state from each path lowers the fringe visibility in proportion to their overlap. Erasing the distinguishing record can restore interference in suitable conditional measurements.

The distinction between lost visibility and returning correlations leads directly to quantum channels, non-Markovian dynamics, and the composition test for memory.

\subsection*{Consequences Forced By The Relation}
Interference visibility falls with the overlap of the environmental states correlated with the interfering alternatives. The phase information can remain in the joint state even when it is absent from all reduced-system observables. When hidden correlations later alter the reduced dynamics, an auxiliary state coordinate or a memory kernel is required for predictive closure.
\subsection*{Domain Of The Construction}
The preferred basis and decoherence rate depend on the interaction, environmental spectrum, and initial system-environment state. Suppression of off-diagonal terms does not by itself derive a unique recorded outcome.

\subsection*{Invariance And Realization}
The joint state retains phase information transferred from the subsystem into system-environment correlations. Reduced interference visibility is fixed by the overlap of environmental states correlated with the alternatives. Equivalent system-environment dilations give the same reduced channel and system observables.

The preferred basis, decay rate, and recoherence depend on the interaction, environmental spectrum, and initial correlations. Changing the system-environment partition changes which correlations are hidden. A Markov approximation removes returning correlations; retaining them introduces auxiliary coordinates or a memory kernel.

\subsection*{Discriminating Consequences}
Interference visibility is compared with a control in which the environment cannot distinguish the alternatives. Two preparations with the same reduced state determine whether hidden correlations alter later observables. Reversing or decoupling the interaction determines whether coherence remains recoverable in the joint state.

\subsection*{Equation Sources}
\begin{itemize}
\item \href{https://arxiv.org/abs/quant-ph/0306087}{arXiv:quant-ph/0306087}
\item \href{https://arxiv.org/abs/quant-ph/0008131}{arXiv:quant-ph/0008131}
\item \href{https://arxiv.org/abs/quant-ph/0211096}{arXiv:quant-ph/0211096}
\end{itemize}



\clearpage
\section{Wave function}
\label{topic:wave-function}

\subsection*{Mechanism}
Wave function is a basis-dependent representative of a pure-state ray; it is not identical to the abstract state or to physical configuration space.

The wave function is a representation of a state in a chosen continuous basis. Treating it as the state itself can obscure that a Fourier transform changes its shape while preserving all predictions.

For a configuration space Q with measure mu, the position wave function is the generalized-basis representative psi(x)=<x|psi> of a ray [psi] in L2(Q,mu). Vectors that differ by a nonzero global phase represent the same pure state. Its modulus squared is a probability density only relative to the stated position measure; spin and particle statistics enlarge or constrain the carrier.

\subsection*{Physical Construction}
The state carrier is the mathematical state object: vector, wavefunction, density operator, coherent state, field state, or register. Operators, maps, and observables become meaningful only after this carrier and its domain have been fixed. Normalization, positivity, inner product, representation, tensor factorization, or superselection conditions define legal states. The calculated quantities include probability distributions obtained by applying the appropriate observables or measurement maps to the carrier.

\subsection*{Topic Equations}
\begin{centeredalign}
\mathcal H=L^2(Q,d\mu),\qquad \psi(x)=\langle x|\psi\rangle\\
\int_Q |\psi(x)|^2\,d\mu(x)=1,\qquad \ket\psi\sim e^{i\alpha}\ket\psi\\
\Pr(X\in\Delta\mid\psi)=\langle\psi|E_X(\Delta)|\psi\rangle=\int_\Delta|\psi(x)|^2\,d\mu(x)\\
\mathcal H_{\mathrm{spin}\,s}=L^2(Q,d\mu)\otimes\mathbb C^{2s+1}
\end{centeredalign}

\subsection*{Physical Meaning}
Its complex value is a probability amplitude. The squared magnitude gives a position density only in the position representation; phase differences remain essential because they control interference and momentum content. Normalization fixes the total probability.

A localized wave packet and its momentum-space Fourier transform emphasize different features of one preparation. Their forms differ, but either representation reproduces the same expectation values when the observables are transformed consistently.

Density operators extend this description to mixtures, subsystems, and open-system dynamics.

\subsection*{Invariance And Realization}
The rule connecting prepared states, observables, and spectral probability measures across wave, matrix, path-integral, circuit, or field notation. The operator-to-spectrum relation: admissible observations are represented through eigenvalues, projections, modes, or outcome channels. The dependence of admissible observable on measurement context or boundary condition. The non-commuting compatibility structure, which survives changes of representation.

The name of the carrier: particle, wave, field, qubit, or excitation. Where time dependence is represented: on the state, on the operator, or in a path weight. The coordinate system, basis, or geometric picture used to display the same relation. The physical implementation of detector, boundary, preparation, or observable.

\subsection*{Discriminating Consequences}
A usable state gives normalized probabilities for every complete observable attached to the selected Hilbert space. Vector, wave-function, density-matrix, and reduced-state forms can describe the same preparation when connected by the appropriate representation map. Physical state changes preserve positivity and trace, or norm in the pure closed-system limit.




\clearpage

\section{Quantum state}

\label{topic:quantum-state}

Preparing a spin along one axis fixes probabilities for measurements along every other axis. The information needed for those predictions can be collected in a density operator \(\rho\). For a two-level system let \(\boldsymbol\sigma=(\sigma_x,\sigma_y,\sigma_z)\) be the Pauli matrices and \(\mathbf r\) the vector of their expectation values. The three components of \(\mathbf r\) determine the state:

\begin{centeredalign}
\rho=\frac12(I+\mathbf r\cdot\boldsymbol\sigma),\qquad |\mathbf r|\leq1.
\end{centeredalign}

The bound follows from positivity: the two eigenvalues are \((1\pm|\mathbf r|)/2\). A unit vector describes a pure preparation. An interior point describes a mixed state, which may arise from uncontrolled preparation or from ignoring a correlated partner. Different ensembles can produce the same density operator. A later experiment acting only on this spin cannot distinguish those ensembles unless additional information about the preparation is accessible.

A Stern--Gerlach measurement along a unit vector \(\mathbf n\) has projectors \(E_\pm=(I\pm\mathbf n\cdot\boldsymbol\sigma)/2\). The probability of each outcome is consequently

\begin{centeredalign}
p_\pm=\operatorname{Tr}(\rho E_\pm)=\frac12(1\pm\mathbf r\cdot\mathbf n).
\end{centeredalign}

The same state gives certainty along its preparation axis and an equal distribution along a perpendicular axis. The uncertainty therefore belongs to the relation between preparation and measurement. It is not evidence that a pure state is an unspecified classical direction. Measurements along three independent axes reconstruct the Bloch vector, while measurements along only one axis leave a disk of compatible states.

For a composite system the density operator also contains correlations. The singlet and an equal mixture of oppositely aligned spins both give \(\rho_A=I/2\) for either individual spin, but they predict different joint measurements. A partial trace retains every expectation of an operator acting on subsystem \(A\) alone. It need not retain the information required for subsequent interacting evolution, because the interaction can convert a joint correlation into a local expectation value.

\begin{centeredalign}
\langle O_A\rangle=\operatorname{Tr}_{AB}[\rho_{AB}(O_A\otimes I_B)]
=\operatorname{Tr}_A(\rho_A O_A),\qquad \rho_A=\operatorname{Tr}_B\rho_{AB}.
\end{centeredalign}

Thus the state required by a mechanism is fixed by the available measurements and subsequent interactions together. Local spin probabilities require \(\rho_A\) at the present time; a prediction after coupling to \(B\) generally requires the joint state or an independently specified preparation of \(B\). This distinction is the physical reason to keep the carrier, interaction and preparation connected in the description.

\subsection*{References For The Physical Derivation}
\begin{itemize}
\item \href{https://arxiv.org/abs/1508.06951}{V. Moretti, Mathematical Foundations of Quantum Mechanics: An Advanced Short Course; states, operators and symmetry.}
\end{itemize}

\subsection*{Relations In The Original Papers}

\href{https://arxiv.org/html/1508.06951#S3.Ex131}{arXiv:1508.06951, S3.Ex131}. The spectral probability of an observable is the trace of the density operator with its spectral projector. Density operators and projection-valued measurements; the Bloch-state example is derived in the chapter.

\clearpage

\section{Two-state quantum system}

\label{topic:two-state-quantum-system}

An avoided crossing couples two states that would otherwise evolve independently. In their basis, a constant energy offset contributes only an overall phase, whereas the remaining Hamiltonian determines both transition probabilities and phase evolution. Let \(\Delta\) be the angular-frequency detuning between the uncoupled levels and \(\Omega\) a real coupling frequency. With \(|0\rangle\) and \(|1\rangle\) the eigenstates of \(\sigma_z\), choose

\begin{centeredalign}
H=\frac{\hbar}{2}(\Delta\sigma_z+\Omega\sigma_x),\qquad \omega_R=\sqrt{\Delta^2+\Omega^2}.
\end{centeredalign}

The Pauli algebra gives \((\Delta\sigma_z+\Omega\sigma_x)^2=\omega_R^2I\). Every higher power in the exponential therefore reduces to either the identity or the Hamiltonian. The evolution follows without fitting a transition curve:

\begin{centeredalign}
U(t)=\cos\frac{\omega_Rt}{2}\,I-i\sin\frac{\omega_Rt}{2}\,
\frac{\Delta\sigma_z+\Omega\sigma_x}{\omega_R}.
\end{centeredalign}

For the preparation \(|0\rangle\), the off-diagonal matrix element of \(U\) gives the transition probability

\begin{centeredalign}
P_{0\to1}(t)=\frac{\Omega^2}{\Delta^2+\Omega^2}
\sin^2\!\left(\frac{t}{2}\sqrt{\Delta^2+\Omega^2}\right).
\end{centeredalign}

On resonance the coupling can transfer the entire population; a pulse of duration \(\pi/|\Omega|\) interchanges the levels. Detuning increases the oscillation frequency but reduces the greatest achievable population transfer. These two consequences distinguish a changed level splitting from a changed coupling even when a short segment of an oscillation looks similar.

The Bloch vector rotates about the effective axis \((\Omega,0,\Delta)\). A second pulse about another axis can convert a relative phase into a population difference. Reversing two such pulses generally changes the result because their Hamiltonians do not commute. The state space is the same in both experiments; the difference lies in the sequence of interactions, and the measured population exposes that difference.

This calculation applies to any isolated pair of levels governed by the stated Hamiltonian, including two coupled modes or selected atomic states. Its use for a larger physical system requires that other levels remain unpopulated and environmental relaxation be negligible over the pulse duration. If the Hamiltonian was obtained in a rotating frame, the rotating-wave approximation and the transformation of the measured operator are additional parts of the correspondence. Specifying a two-dimensional matrix does not by itself establish those conditions in a device.

\subsection*{References For The Physical Derivation}
\begin{itemize}
\item \href{https://arxiv.org/abs/1508.06951}{V. Moretti, Mathematical Foundations of Quantum Mechanics: An Advanced Short Course; states, operators and symmetry.}
\end{itemize}



\clearpage
\section{Qubit}
\label{topic:qubit}

\subsection*{Mechanism}
Qubit is the two-dimensional state-carrier constructor used when the admissible state space is \textbackslash{}(\textbackslash{}mathbb C\textasciicircum{}2\textbackslash{}).

Quantum probabilities are calculated from a state. A Hamiltonian or an observable does not make a prediction by itself; it must act on a normalized state vector, density operator, or statistical sector that records the preparation.

A qubit is the minimal quantum state space with a basis, amplitudes, unitary control, and measurement observable. Bloch-vector language is a representation of the same two-dimensional carrier.

\subsection*{Physical Construction}
The state carrier is the mathematical state object: vector, wavefunction, density operator, coherent state, field state, or register. Operators, maps, and observables become meaningful only after this carrier and its domain have been fixed. Normalization, positivity, inner product, representation, tensor factorization, or superselection conditions define legal states. The calculated quantities include probability distributions obtained by applying the appropriate observables or measurement maps to the carrier.

\subsection*{Topic Equations}
\begin{centeredalign}
\ket{\psi}=\alpha\ket{0}+\beta\ket{1},\qquad |\alpha|^2+|\beta|^2=1\\
\rho=\frac12(I+\mathbf r\cdot\boldsymbol\sigma),\qquad |\mathbf r|\le1\\
p(0)=|\langle0|\psi\rangle|^2,\qquad p(1)=|\langle1|\psi\rangle|^2
\end{centeredalign}

\subsection*{Physical Meaning}
Normalization guarantees that the probabilities sum to one, while positivity prevents negative probabilities. Pure vectors and density operators are not competing theories: the density-operator form also represents mixtures and reduced states obtained when unobserved degrees of freedom are traced out.

The Hamiltonian or channel evolves the prepared state. An observable and measurement map then convert that evolved state into outcome probabilities.

\subsection*{Invariance And Realization}
The rule connecting prepared states, observables, and spectral probability measures across wave, matrix, path-integral, circuit, or field notation. The operator-to-spectrum relation: admissible observations are represented through eigenvalues, projections, modes, or outcome channels. The dependence of admissible observable on measurement context or boundary condition. The non-commuting compatibility structure, which survives changes of representation.

The name of the carrier: particle, wave, field, qubit, or excitation. Where time dependence is represented: on the state, on the operator, or in a path weight. The coordinate system, basis, or geometric picture used to display the same relation. The physical implementation of detector, boundary, preparation, or observable.

\subsection*{Discriminating Consequences}
A usable state gives normalized probabilities for every complete observable attached to the selected Hilbert space. Vector, wave-function, density-matrix, and reduced-state forms can describe the same preparation when connected by the appropriate representation map. Physical state changes preserve positivity and trace, or norm in the pure closed-system limit.




\clearpage
\section{Quantum entanglement}
\label{topic:quantum-entanglement}

\subsection*{Mechanism}
Quantum entanglement is a property of a composite state that cannot be written as a classical mixture of product states across a chosen subsystem decomposition; for a pure state, this reduces to failure to factor.

Entanglement is required when the state of a composite system cannot be assembled from independent states of its parts. It changes which correlations are possible and makes subsystem states mixed even when the total state is pure.

The tensor-product decomposition specifies what counts as subsystem A and subsystem B. A pure state is entangled when its Schmidt rank exceeds one. Each subsystem can then be mixed even though the joint state is pure, because partial tracing discards the correlations that purify it. Local basis changes preserve the Schmidt coefficients, while a different physical factorization can change whether the same vector is called entangled. Bell measurements test whether the resulting correlations admit a local hidden-variable model.

\subsection*{Physical Construction}
The state carrier is a composite Hilbert space with a physically specified subsystem algebra or tensor-product factorization. The calculation involves Schmidt decomposition, partial trace, local observables, and joint correlation operators. The joint density operator is positive and normalized; separability is defined relative to the chosen subsystem structure. The calculated quantities include reduced-state spectra, entanglement entropy, correlation witnesses, and Bell parameters.

\subsection*{Topic Equations}
\begin{centeredalign}
\mathcal H_{AB}=\mathcal H_A\otimes\mathcal H_B\\
\ket{\Psi}_{AB}=\sum_k\sqrt{\lambda_k}\ket{k_A}\ket{k_B},\qquad \sum_k\lambda_k=1\\
\rho_A=\operatorname{Tr}_B\ket{\Psi}\bra{\Psi},\qquad S_A=-\operatorname{Tr}(\rho_A\log\rho_A)\\
|S_{\mathrm{CHSH}}|\leq2,\qquad |S_{\mathrm{CHSH}}|_{\mathrm{QM}}\leq2\sqrt2
\end{centeredalign}

\subsection*{Physical Meaning}
The tensor-product structure defines the proposed subsystems. A state is entangled when it is not separable across that factorization. Partial traces describe each part, while joint observables expose correlations that no product state can reproduce.

In a Bell pair, each spin alone is maximally mixed, yet measurements on the pair are strongly correlated. The information resides in the relation between subsystems rather than in either local state.

Bell's theorem and contextuality determine which entangled correlations are incompatible with classical joint assignments.

\subsection*{Consequences Forced By The Relation}
A pure bipartite state with more than one nonzero Schmidt coefficient gives mixed reduced states and correlations unavailable to product preparations. Local unitary changes of basis preserve the Schmidt coefficients and the pure-state entanglement entropy. Suitable local measurements can violate a Bell inequality even though each subsystem alone carries no corresponding pure state.
\subsection*{Domain Of The Construction}
Entanglement is defined relative to a physical subsystem algebra or tensor-product factorization. For mixed states, nonfactorization of one decomposition is insufficient; separability requires testing all convex product decompositions.

\subsection*{Invariance And Realization}
Schmidt coefficients and pure-state entanglement entropy are invariant under local unitary changes of basis. The reduced-state spectra preserve the amount of bipartite pure-state entanglement. Correlations remain joint properties even when neither subsystem has a pure local state.

The subsystem factorization and accessible observable algebra determine which correlations count as entanglement. Noise, loss, and coarse graining can convert pure-state entanglement into mixed-state correlations. Photons, spins, atoms, modes, and encoded qubits can realize the same Schmidt structure.

\subsection*{Discriminating Consequences}
Local measurements reconstruct a correlation witness or Bell parameter that product preparations cannot reproduce. Independent local-unitary rotations leave the inferred Schmidt spectrum unchanged. A separable-state control fixes the correlation background of the apparatus.

\subsection*{Equation Sources}
\begin{itemize}
\item \href{https://arxiv.org/abs/math-ph/0211065}{arXiv:math-ph/0211065}
\end{itemize}



\chapter{Dynamics and transformations}

\begin{claimbox}

Hamiltonians, unitary maps, equations of motion, and path weights describe the lawful transport of the state before a question is resolved.

\end{claimbox}

Hamiltonians, channels, and path amplitudes determine how a prepared state changes. Equivalent formulations belong together when they preserve amplitudes or expectation values on a common domain.

\section*{Topics}

\begin{longtable}{p{0.92\linewidth}}
\toprule
Page \\
\midrule
\endhead
\phantomsection\label{page:unitary-operator}Unitary operator \\
\phantomsection\label{page:path-integral-formulation}Path integral formulation \emph{(see Path Integral)} \\
\phantomsection\label{page:path-integral}Path integral \\
\phantomsection\label{page:hamiltonian-quantum-mechanics}Hamiltonian (quantum mechanics) \emph{(see Hamiltonian Mechanics)} \\
\phantomsection\label{page:schr-dinger-equation}Schrödinger equation \\
\phantomsection\label{page:symmetry-in-quantum-mechanics}Symmetry in quantum mechanics \\
\phantomsection\label{page:heisenberg-picture}Heisenberg picture \\
\phantomsection\label{page:creation-and-annihilation-operators}Creation and annihilation operators \\
\bottomrule
\end{longtable}

\clearpage
\section{Unitary operator}
\label{topic:unitary-operator}

\subsection*{Mechanism}
Unitary operator is the reversible-map constructor: it changes the state while preserving inner products and probabilities.

Unitary operators express reversible quantum change and changes of representation with one mathematical condition: preservation of the inner product. That condition preserves normalization, orthogonality, and transition probabilities.

Unitary maps are the admissible reversible transformations of closed-system quantum theory. They preserve normalization and carry projective geometry of the state space through time or controlled operations.

\subsection*{Physical Construction}
The state carrier is a state vector, density operator, wavefunction, field state, or register on a specified domain. The calculation involves hamiltonian, unitary map, channel generator, action, constraint, or differential operator that transports the state. Self-adjointness, complete positivity, trace preservation, gauge constraints, and boundary/domain conditions decide whether the evolution is legal. The calculated quantities include time-dependent probabilities, spectra, transition amplitudes, conserved quantities, or response functions.

\subsection*{Topic Equations}
\begin{centeredalign}
U^\dagger U=UU^\dagger=I\\
\ket{\psi'}=U\ket{\psi}\\
\langle\psi'|\phi'\rangle=\langle\psi|\phi\rangle
\end{centeredalign}

\subsection*{Physical Meaning}
Applying a unitary map to a state rotates it within Hilbert space without losing information. Conjugating an observable by the same map gives the equivalent description in the transformed representation. Continuous unitary evolution is generated by a self-adjoint operator.

A spin rotation changes the amplitudes assigned to vertical and horizontal outcomes, but a second inverse rotation restores the original state exactly. Decoherence cannot be represented by such a unitary map on the subsystem alone.

The Hamiltonian page identifies the generator that produces unitary time evolution.

\subsection*{Invariance And Realization}
The rule connecting prepared states, observables, and spectral probability measures across wave, matrix, path-integral, circuit, or field notation. The operator-to-spectrum relation: admissible observations are represented through eigenvalues, projections, modes, or outcome channels. The dependence of admissible observable on measurement context or boundary condition. The non-commuting compatibility structure, which survives changes of representation.

The name of the carrier: particle, wave, field, qubit, or excitation. Where time dependence is represented: on the state, on the operator, or in a path weight. The coordinate system, basis, or geometric picture used to display the same relation. The physical implementation of detector, boundary, preparation, or observable.

\subsection*{Discriminating Consequences}
Lawful closed-system evolution preserves norm or trace; open-system evolution must preserve positivity and trace under the stated approximation. The short-time and classical limits identify whether the generator has the correct physical regime. Hamiltonian, propagator, and path-integral forms are equivalent only when they yield the same transition amplitudes or correlation functions.




\clearpage
\section{Path integral formulation}
\label{topic:path-integral-formulation}

\subsection*{Mechanism}
Path integral formulation is the formulation in which the generator is represented by action-weighted histories.

Kinematics and dynamics are separate. The state space lists what can exist, whereas a Hamiltonian, action, Liouvillian, or channel generator specifies which changes are allowed and on what timescale.

This page belongs with lawful change because it changes how the transport step is calculated. The observable predictions remain probabilities after amplitudes are composed and squared or traced.

\subsection*{Physical Construction}
The state carrier is a state vector, density operator, wavefunction, field state, or register on a specified domain. The calculation involves hamiltonian, unitary map, channel generator, action, constraint, or differential operator that transports the state. Self-adjointness, complete positivity, trace preservation, gauge constraints, and boundary/domain conditions decide whether the evolution is legal. The calculated quantities include time-dependent probabilities, spectra, transition amplitudes, conserved quantities, or response functions.

\subsection*{Topic Equations}
\begin{centeredalign}
\langle q_f,t_f|q_i,t_i\rangle=\int\mathcal Dq\,e^{iS[q]/\hbar}\\
Z[J]=\int\mathcal D\phi\,\exp\!\left(\frac{i}{\hbar}(S[\phi]+\int J\phi)\right)
\end{centeredalign}

\subsection*{Physical Meaning}
The evolution law transports a state without redefining it. Closed-system evolution is unitary; effective open-system evolution must preserve trace and positivity. Equivalent Hamiltonian, propagator, and path-integral descriptions agree on transition amplitudes.

The evolved state becomes experimentally meaningful through an observable whose spectrum and expectation values expose the consequences of the dynamics.

\subsection*{Invariance And Realization}
The rule connecting prepared states, observables, and spectral probability measures across wave, matrix, path-integral, circuit, or field notation. The operator-to-spectrum relation: admissible observations are represented through eigenvalues, projections, modes, or outcome channels. The dependence of admissible observable on measurement context or boundary condition. The non-commuting compatibility structure, which survives changes of representation.

The name of the carrier: particle, wave, field, qubit, or excitation. Where time dependence is represented: on the state, on the operator, or in a path weight. The coordinate system, basis, or geometric picture used to display the same relation. The physical implementation of detector, boundary, preparation, or observable.

\subsection*{Discriminating Consequences}
Lawful closed-system evolution preserves norm or trace; open-system evolution must preserve positivity and trace under the stated approximation. The short-time and classical limits identify whether the generator has the correct physical regime. Hamiltonian, propagator, and path-integral forms are equivalent only when they yield the same transition amplitudes or correlation functions.




\clearpage
\section{Path integral}
\label{topic:path-integral}

\subsection*{Mechanism}
Path integral is an alternate generator constructor: transition amplitudes are obtained by summing phase weights over histories.

Kinematics and dynamics are separate. The state space lists what can exist, whereas a Hamiltonian, action, Liouvillian, or channel generator specifies which changes are allowed and on what timescale.

The path integral does not replace the operator constructor. It repackages the generator step as a weighted sum over histories between boundary conditions. It is especially useful when action, symmetry, and field degrees of freedom are more natural than state-vector evolution.

\subsection*{Physical Construction}
The state carrier is a state vector, density operator, wavefunction, field state, or register on a specified domain. The calculation involves hamiltonian, unitary map, channel generator, action, constraint, or differential operator that transports the state. Self-adjointness, complete positivity, trace preservation, gauge constraints, and boundary/domain conditions decide whether the evolution is legal. The calculated quantities include time-dependent probabilities, spectra, transition amplitudes, conserved quantities, or response functions.

\subsection*{Topic Equations}
\begin{centeredalign}
K(x_f,t_f;x_i,t_i)=\int_{x_i}^{x_f}\mathcal D x(t)\,\exp\!\left(\frac{i}{\hbar}S[x]\right)\\
\psi(x_f,t_f)=\int K(x_f,t_f;x_i,t_i)\psi(x_i,t_i)\,dx_i
\end{centeredalign}

\subsection*{Physical Meaning}
The evolution law transports a state without redefining it. Closed-system evolution is unitary; effective open-system evolution must preserve trace and positivity. Equivalent Hamiltonian, propagator, and path-integral descriptions agree on transition amplitudes.

The evolved state becomes experimentally meaningful through an observable whose spectrum and expectation values expose the consequences of the dynamics.

\subsection*{Invariance And Realization}
The rule connecting prepared states, observables, and spectral probability measures across wave, matrix, path-integral, circuit, or field notation. The operator-to-spectrum relation: admissible observations are represented through eigenvalues, projections, modes, or outcome channels. The dependence of admissible observable on measurement context or boundary condition. The non-commuting compatibility structure, which survives changes of representation.

The name of the carrier: particle, wave, field, qubit, or excitation. Where time dependence is represented: on the state, on the operator, or in a path weight. The coordinate system, basis, or geometric picture used to display the same relation. The physical implementation of detector, boundary, preparation, or observable.

\subsection*{Discriminating Consequences}
Lawful closed-system evolution preserves norm or trace; open-system evolution must preserve positivity and trace under the stated approximation. The short-time and classical limits identify whether the generator has the correct physical regime. Hamiltonian, propagator, and path-integral forms are equivalent only when they yield the same transition amplitudes or correlation functions.




\clearpage
\section{Hamiltonian (quantum mechanics)}
\label{topic:hamiltonian-quantum-mechanics}

\subsection*{Mechanism}
Hamiltonian (quantum mechanics) is the generator observable: it both transports states and supplies the energy spectrum.

The Hamiltonian specifies both the energy observable and, for a closed system, the generator of time translation. These roles coincide but should not be confused: one concerns possible energy values, the other the path followed by every prepared state.

The Hamiltonian has a dual role. Dynamically, it generates unitary time evolution. Spectrally, its eigenvalues are admissible energy observables. This dual role is one reason the operator/spectrum branch is central.

\subsection*{Physical Construction}
The state carrier is a state vector, density operator, wavefunction, field state, or register on a specified domain. The calculation involves hamiltonian, unitary map, channel generator, action, constraint, or differential operator that transports the state. Self-adjointness, complete positivity, trace preservation, gauge constraints, and boundary/domain conditions decide whether the evolution is legal. The calculated quantities include time-dependent probabilities, spectra, transition amplitudes, conserved quantities, or response functions.

\subsection*{Topic Equations}
\begin{centeredalign}
H\ket{E_n}=E_n\ket{E_n}\\
U(t)=e^{-iHt/\hbar}\\
\rho(t)=U(t)\rho(0)U(t)^\dagger
\end{centeredalign}

\subsection*{Physical Meaning}
Exponentiating a self-adjoint Hamiltonian produces the unitary propagator. Its eigenvectors acquire phases at rates set by their energies; relative phases then produce interference and motion. Time dependence or external driving requires a time-ordered propagator.

For a particle in a static potential, kinetic and potential terms determine both stationary energy levels and the evolution of a wave packet assembled from those levels.

The Schrödinger equation gives the differential form of this evolution, while path integrals and the Heisenberg picture reorganize the same predictions.

\subsection*{Invariance And Realization}
The rule connecting prepared states, observables, and spectral probability measures across wave, matrix, path-integral, circuit, or field notation. The operator-to-spectrum relation: admissible observations are represented through eigenvalues, projections, modes, or outcome channels. The dependence of admissible observable on measurement context or boundary condition. The non-commuting compatibility structure, which survives changes of representation.

The name of the carrier: particle, wave, field, qubit, or excitation. Where time dependence is represented: on the state, on the operator, or in a path weight. The coordinate system, basis, or geometric picture used to display the same relation. The physical implementation of detector, boundary, preparation, or observable.

\subsection*{Discriminating Consequences}
Lawful closed-system evolution preserves norm or trace; open-system evolution must preserve positivity and trace under the stated approximation. The short-time and classical limits identify whether the generator has the correct physical regime. Hamiltonian, propagator, and path-integral forms are equivalent only when they yield the same transition amplitudes or correlation functions.




\clearpage
\section{Schrödinger equation}
\label{topic:schr-dinger-equation}

\subsection*{Mechanism}
Schrödinger equation is the state-transport constructor: the Hamiltonian generates lawful change of the state before measurement.

The Schrödinger equation turns a Hamiltonian into a local rule for the time dependence of a state. It is the point at which the chosen state space, boundary conditions, and interaction model become a calculable prediction.

The Schrödinger equation is not a measurement rule. It is the generator step of the quantum constructor. It evolves the predictive carrier while preserving normalization when the Hamiltonian is self-adjoint.

\subsection*{Physical Construction}
The state carrier is a state vector, density operator, wavefunction, field state, or register on a specified domain. The calculation involves hamiltonian, unitary map, channel generator, action, constraint, or differential operator that transports the state. Self-adjointness, complete positivity, trace preservation, gauge constraints, and boundary/domain conditions decide whether the evolution is legal. The calculated quantities include time-dependent probabilities, spectra, transition amplitudes, conserved quantities, or response functions.

\subsection*{Topic Equations}
\begin{centeredalign}
i\hbar\,\partial_t\ket{\psi(t)}=H\ket{\psi(t)}\\
\ket{\psi(t)}=U(t)\ket{\psi(0)},\qquad U(t)=e^{-iHt/\hbar}\\
\frac{d}{dt}\langle\psi(t)|\psi(t)\rangle=0\quad (H=H^\dagger)
\end{centeredalign}

\subsection*{Physical Meaning}
The time derivative of the state is fixed by the Hamiltonian acting on that state. For a self-adjoint Hamiltonian, the equation preserves norm. Stationary solutions separate time dependence from the spatial eigenvalue problem, but a general preparation is their superposition.

A wave packet in a potential well spreads and interferes according to the same equation whose stationary solutions define the well's energy levels.

Observables determine which consequences of the evolved state are compared with experiment.

\subsection*{Invariance And Realization}
The rule connecting prepared states, observables, and spectral probability measures across wave, matrix, path-integral, circuit, or field notation. The operator-to-spectrum relation: admissible observations are represented through eigenvalues, projections, modes, or outcome channels. The dependence of admissible observable on measurement context or boundary condition. The non-commuting compatibility structure, which survives changes of representation.

The name of the carrier: particle, wave, field, qubit, or excitation. Where time dependence is represented: on the state, on the operator, or in a path weight. The coordinate system, basis, or geometric picture used to display the same relation. The physical implementation of detector, boundary, preparation, or observable.

\subsection*{Discriminating Consequences}
Lawful closed-system evolution preserves norm or trace; open-system evolution must preserve positivity and trace under the stated approximation. The short-time and classical limits identify whether the generator has the correct physical regime. Hamiltonian, propagator, and path-integral forms are equivalent only when they yield the same transition amplitudes or correlation functions.




\clearpage

\section{Symmetry in quantum mechanics}

\label{topic:symmetry-in-quantum-mechanics}

A rotation of an isolated atom changes the description of its orientation while leaving transition probabilities unchanged. Such transformations act on quantum rays and are represented, under the hypotheses of Wigner's theorem, by unitary or antiunitary maps. Continuous transformations connected to the identity are represented unitarily; their generators turn an invariance of probabilities into a relation between operators.

Let a one-parameter symmetry be \(U(\epsilon)=\exp(-i\epsilon G/\hbar)\), with self-adjoint generator \(G\) and a parameter whose units make the exponent dimensionless. For a time-independent Hamiltonian, invariance under this group implies

\begin{centeredalign}
U(\epsilon)HU(\epsilon)^\dagger=H,
\qquad [G,H]=0,\qquad \frac{d}{dt}\langle G\rangle=0.
\end{centeredalign}

The last equality follows from the Heisenberg equation when \(G\) has no explicit time dependence and the operator domains permit the calculation. Spatial translations give momentum conservation; rotations give angular momentum conservation. A boundary or an external field can break the symmetry even if a local bulk term remains invariant. Conservation therefore belongs to the full physical problem.

Symmetry also controls which transitions can occur. Suppose parity \(\Pi\) commutes with \(H\) and two nondegenerate states have definite parities \(\pi_i,\pi_f=\pm1\). Position is odd under parity. Inserting \(\Pi^\dagger\Pi\) into its matrix element yields

\begin{centeredalign}
\langle f|X|i\rangle=-\pi_f\pi_i\langle f|X|i\rangle.
\end{centeredalign}

The electric-dipole matrix element vanishes for states of the same parity. This is a selection rule for the pair consisting of the dynamics and the coupling to the probe. It can suppress a spectral line without creating an additional conserved quantity. Conversely, an accidental equality of two energy differences can merge lines without imposing any selection rule. These mechanisms must be distinguished when a spectral pattern simplifies.

For a proposed transfer between physical systems, a shared symmetry is useful because it restricts the admissible target interactions and observables. It does not determine their coupling constants or all matrix elements. The construction becomes predictive only after the target representation, Hamiltonian and measurement are fixed. Symmetry then supplies exact relations, such as forbidden transitions, that can test the proposed realization.

\subsection*{References For The Physical Derivation}
\begin{itemize}
\item \href{https://arxiv.org/abs/1508.06951}{V. Moretti, Mathematical Foundations of Quantum Mechanics: An Advanced Short Course; states, operators and symmetry.}
\end{itemize}



\clearpage

\section{Heisenberg picture}

\label{topic:heisenberg-picture}

The Schrödinger and Heisenberg pictures assign time dependence to different parts of the same prediction. With unitary evolution \(U(t)\), an initial density operator \(\rho_0\) and an observable \(O\), the expectation can be evaluated either by evolving the state or by evolving the operator:

\begin{centeredalign}
\rho_S(t)=U(t)\rho_0U(t)^\dagger,\qquad
O_H(t)=U(t)^\dagger OU(t),\qquad
\operatorname{Tr}[\rho_S(t)O]=\operatorname{Tr}[\rho_0O_H(t)].
\end{centeredalign}

Differentiating \(O_H\) for a time-independent Hamiltonian \(H\) gives the Heisenberg equation:

\begin{centeredalign}
\dot O_H=\frac{i}{\hbar}[H,O_H].
\end{centeredalign}

The resulting commutators identify which other observables enter the prediction. This is particularly useful when an experiment observes only a small part of a many-body system.

Consider two spins with \(H=\hbar g\,Z_1Z_2\), where \(X_j,Y_j,Z_j\) denote dimensionless Pauli operators and \(g\) is a frequency. If the measured quantity is \(x=\langle X_1\rangle\), its derivative contains the correlation \(c=\langle Y_1Z_2\rangle\). A second commutator closes the pair:

\begin{centeredalign}
\dot x=-2gc,\qquad \dot c=2gx,\qquad
x(t)=x(0)\cos(2gt)-c(0)\sin(2gt).
\end{centeredalign}

The correlation is physically necessary: preparations with the same initial transverse magnetization but different \(c(0)\) give different later magnetizations. Its appearance is fixed by the interaction algebra. Retaining the two expectations supplies a complete closed prediction for this observable even though it does not reconstruct every entry of the two-spin density operator.

Eliminating the correlation replaces the pair of first-order equations by a history-dependent equation for \(x\). The initial correlation survives as a separate term:

\begin{centeredalign}
\dot x(t)=-2g c(0)-4g^2\int_0^t x(s)\,ds.
\end{centeredalign}

Both descriptions give the same signal when they use the same preparation. Discarding the integral or the initial-correlation term gives a different physical prediction. Repeated commutators thus provide a concrete way to construct the set of observables needed by a measurement, and to identify when a smaller description requires memory. For more complicated interactions the sequence may generate a much larger space rather than closing after two steps.

\subsection*{References For The Physical Derivation}
\begin{itemize}
\item \href{https://arxiv.org/abs/1508.06951}{V. Moretti, Mathematical Foundations of Quantum Mechanics: An Advanced Short Course; states, operators and symmetry.}
\end{itemize}

\subsection*{Relations In The Original Papers}

\href{https://arxiv.org/html/1508.06951#S3.Ex177}{arXiv:1508.06951, S3.Ex177}. State evolution and observable evolution give identical spectral probabilities. The source chooses its unitary convention explicitly. The two-spin calculation uses U = exp(-iHt/hbar).

\clearpage
\section{Creation and annihilation operators}
\label{topic:creation-and-annihilation-operators}

\subsection*{Mechanism}
Creation and annihilation operators are sector-changing operators: they add or remove one quantum from a mode and make many-body or field descriptions executable.

Kinematics and dynamics are separate. The state space lists what can exist, whereas a Hamiltonian, action, Liouvillian, or channel generator specifies which changes are allowed and on what timescale.

The page is about the algebraic move that changes occupation number. Creation raises the population of a mode, annihilation lowers it, and the commutation or anticommutation rule determines the statistics. The number operator gives the spectral prediction.

\subsection*{Physical Construction}
The state carrier is a state vector, density operator, wave function, field state, or register on a specified domain. The calculation involves a Hamiltonian, action, Liouvillian, channel generator, or differential operator that transports the state. Self-adjointness, complete positivity, trace preservation, gauge constraints, and boundary conditions determine whether the evolution is legal. The calculated quantities include time-dependent probabilities, transition amplitudes, response functions, conserved quantities, or spectra implied by the dynamics.

\subsection*{Topic Equations}
\begin{centeredalign}
a_i^\dagger\ket{\ldots,n_i,\ldots}=\sqrt{n_i+1}\ket{\ldots,n_i+1,\ldots}\\
a_i\ket{\ldots,n_i,\ldots}=\sqrt{n_i}\ket{\ldots,n_i-1,\ldots}\\
N_i=a_i^\dagger a_i
\end{centeredalign}

\subsection*{Physical Meaning}
The evolution law transports a state without redefining it. Closed-system evolution is unitary; effective open-system evolution must preserve trace and positivity. Equivalent Hamiltonian, propagator, and path-integral descriptions agree on transition amplitudes.

The evolved state becomes experimentally meaningful through an observable whose spectrum and expectation values expose the consequences of the dynamics.

\subsection*{Invariance And Realization}
The relation between prepared states, observables, and spectral probability measures. The use of eigenvalues, projectors, modes, or outcome channels to represent admissible observations. The dependence of the observable on basis, domain, potential, preparation, or measurement context. The commutator structure that limits which observables can be jointly diagonalized.

The physical carrier: particle, wave, field mode, spin, qubit, detector, or excitation. The representation: wave mechanics, matrix mechanics, density matrices, path integrals, circuits, or fields. Where time dependence is placed: on the state, on the operator, in a propagator, or in a path weight. The implementation of preparation, boundary condition, detector, or outcome channel.

\subsection*{Discriminating Consequences}
A transfer target provides a state space, a transformation law, and a spectral or categorical observable, with one compatibility relation experimentally unresolved. A useful validation varies the basis, domain, or measurement context and measures whether the allowed observable changes while the underlying transformation law remains identifiable. A stronger validation contains two candidate observables whose predicted commutator controls joint resolvability.




\part{Observables, Compatibility, And Probability}

\chapter{Observables and spectra}

\begin{claimbox}

An observable is a permitted question whose operator form determines the possible numerical answers.

\end{claimbox}

An observable determines possible outcomes through its spectral measure. The state then fixes their probabilities.

\section*{Topics}

\begin{longtable}{p{0.92\linewidth}}
\toprule
Page \\
\midrule
\endhead
\phantomsection\label{page:spectral-theory}Spectral theory \\
\phantomsection\label{page:pauli-matrices}Pauli matrices \\
\phantomsection\label{page:angular-momentum-operator}Angular momentum operator \\
\phantomsection\label{page:observable}Observable \\
\phantomsection\label{page:self-adjoint-operator}Self-adjoint operator \\
\phantomsection\label{page:heisenberg-group}Heisenberg group \\
\bottomrule
\end{longtable}

\clearpage

\section{Spectral theory}

\label{topic:spectral-theory}

An observable specifies both possible outcomes and how a prepared state distributes probability among them. For a self-adjoint operator \(A\), the spectral measure \(E_A\) assigns an orthogonal projector to each measurable set of real outcomes. The operator and its measurement probabilities are reconstructed from this measure:

\begin{centeredalign}
A=\int_{\mathbb R}\lambda\,dE_A(\lambda),\qquad
p(\lambda\in B)=\operatorname{Tr}[\rho E_A(B)].
\end{centeredalign}

Here \(B\) is an interval or other measurable outcome set and \(\rho\) the density operator. For a discrete spectrum the integral becomes a sum over eigenspace projectors. For position or free-particle momentum it is genuinely continuous; their ideal eigenvectors are generalized distributions rather than normalizable prepared states. A detector with finite resolution measures an interval probability, not a normalizable state at one exact continuum value.

The same spectral measure determines functions of the operator. If \(H\) is the time-independent Hamiltonian, the function \(\exp(-iEt/\hbar)\) gives its unitary evolution. An energy decomposition therefore connects a spectroscopic question to dynamics:

\begin{centeredalign}
U(t)=\int e^{-iEt/\hbar}\,dE_H(E),\qquad
\langle\psi|U(t)|\psi\rangle=\int e^{-iEt/\hbar}\,d\mu_\psi(E),
\end{centeredalign}

where \(\mu_\psi(B)=\langle\psi|E_H(B)|\psi\rangle\) is the energy distribution of the preparation. The survival amplitude is its Fourier transform. A narrow energy distribution changes phase slowly relative to itself; a broad distribution can dephase rapidly. This connects spectral width to temporal evolution without assuming irreversible decay.

The spectrum belongs to the operator with its domain. On an interval of length \(L\), the same differential expression \(-\hbar^2\partial_x^2/(2m)\) has a zero-energy constant mode with Neumann conditions, while Dirichlet conditions exclude it and begin at \(\pi^2\hbar^2/(2mL^2)\). A change of boundary therefore changes both the energy outcomes and the time evolution, even though the bulk equation is unchanged.

Two Hamiltonians with identical energy values can still have different observable matrix elements. Transporting a mechanism requires a correspondence between states and observables, not just a matched list of eigenvalues. Spectral theory supplies the predictions once that correspondence and the operator domains are specified; it also identifies precisely what an isospectral comparison leaves undecided.

\subsection*{References For The Physical Derivation}
\begin{itemize}
\item \href{https://arxiv.org/abs/1508.06951}{V. Moretti, Mathematical Foundations of Quantum Mechanics: An Advanced Short Course; states, operators and symmetry.}
\end{itemize}

\subsection*{Relations In The Original Papers}

\href{https://arxiv.org/html/1508.06951#S2.EGx25}{arXiv:1508.06951, S2.EGx25}. Functions of a self-adjoint observable are defined by integration against its spectral measure. The domain of an unbounded function of the operator must be retained.

\clearpage

\section{Pauli matrices}

\label{topic:pauli-matrices}

A two-level system has three independent traceless Hermitian observables. Choosing an orthonormal basis of states identifies them with the Pauli matrices. Their multiplication law determines rotations, transition amplitudes and the closure of equations for spin expectations:

\begin{centeredalign}
\sigma_x=\begin{pmatrix}0&1\\1&0\end{pmatrix},\quad
\sigma_y=\begin{pmatrix}0&-i\\i&0\end{pmatrix},\quad
\sigma_z=\begin{pmatrix}1&0\\0&-1\end{pmatrix},
\end{centeredalign}

\begin{centeredalign}
\sigma_i\sigma_j=\delta_{ij}I+i\sum_k\epsilon_{ijk}\sigma_k.
\end{centeredalign}

For a unit vector \(\mathbf n\), this identity gives \((\mathbf n\cdot\boldsymbol\sigma)^2=I\). Exponentiating a Hamiltonian proportional to that operator therefore separates into a cosine and a sine. The rotation of a spinor through angle \(\theta\) is

\begin{centeredalign}
U(\theta)=e^{-i\theta\mathbf n\cdot\boldsymbol\sigma/2}
=\cos(\theta/2)I-i\sin(\theta/2)\mathbf n\cdot\boldsymbol\sigma.
\end{centeredalign}

A full \(2\pi\) rotation gives \(-I\). This phase leaves an isolated spin probability unchanged, but becomes observable relative to a coherent reference branch that was not rotated. The Bloch vector itself completes one ordinary spatial rotation. The distinction between spinor amplitude and observable vector is the reason the half-angle appears.

Writing a two-level Hamiltonian as \(H=E_0I+(\hbar/2)\boldsymbol\omega\cdot\boldsymbol\sigma\) and a state as \(\rho=(I+\mathbf r\cdot\boldsymbol\sigma)/2\) gives \(\dot{\mathbf r}=\boldsymbol\omega\times\mathbf r\). The identity term contributes only an overall phase, while the traceless part produces precession. The vector's length is conserved by this isolated unitary dynamics; relaxation requires an additional physical coupling.

For multiple spins, tensor products of Pauli matrices form an operator basis. Commuting a Hamiltonian with a one-spin observable can produce a two-spin product, showing how a local signal becomes coupled to a correlation. Counting the independent operators generated by those commutators measures the size of that particular observable dynamics. A reduction in this size can result from symmetry, selection rules or equal transition frequencies; the Pauli algebra permits these possibilities to be tested separately.

\subsection*{References For The Physical Derivation}
\begin{itemize}
\item \href{https://arxiv.org/abs/1508.06951}{V. Moretti, Mathematical Foundations of Quantum Mechanics: An Advanced Short Course; states, operators and symmetry.}
\end{itemize}

\subsection*{Relations In The Original Papers}

\href{https://arxiv.org/html/quant-ph/9705052#Ch1.E3}{arXiv:quant-ph/9705052, Ch1.E3}. The Pauli matrices represent the three spin components used to calculate rotations and measurement probabilities. One spin-1/2 degree of freedom; these finite matrices do not realize a canonical position-momentum pair.

\clearpage

\section{Angular momentum operator}

\label{topic:angular-momentum-operator}

Rotations about different axes change orientation in an order-dependent way. In quantum mechanics their generators \(J_x,J_y,J_z\) obey an algebra whose representations determine the allowed angular-momentum values. With \(\hbar\) the reduced Planck constant,

\begin{centeredalign}
[J_i,J_j]=i\hbar\sum_k\epsilon_{ijk}J_k,\qquad
[J^2,J_i]=0,\qquad J^2=J_x^2+J_y^2+J_z^2.
\end{centeredalign}

A simultaneous eigenstate of \(J^2\) and \(J_z\) can be labeled \(|j,m\rangle\). Define \(J_\pm=J_x\pm iJ_y\). Their commutators with \(J_z\) show that they raise or lower its eigenvalue by \(\hbar\), while preserving \(j\). Positivity of the squared norms fixes the endpoints of the ladder:

\begin{centeredalign}
J^2|j,m\rangle=\hbar^2j(j+1)|j,m\rangle,\qquad
J_z|j,m\rangle=\hbar m|j,m\rangle,
\end{centeredalign}

\begin{centeredalign}
J_\pm|j,m\rangle=\hbar\sqrt{j(j+1)-m(m\pm1)}\,|j,m\pm1\rangle.
\end{centeredalign}

The ladder terminates at \(m=\pm j\), giving \(2j+1\) states. Single-valued orbital wave functions on ordinary three-dimensional space carry integer orbital angular momentum. Intrinsic spin uses representations of the covering group SU(2), allowing half-integer \(j\). The same local commutator algebra is therefore realized on carriers with different global transformation properties.

For two spins, total angular momentum is \(\mathbf J=\mathbf J_1+\mathbf J_2\). An isotropic exchange interaction can be rewritten using the total Casimir:

\begin{centeredalign}
H=\frac{K}{\hbar^2}\mathbf J_1\cdot\mathbf J_2
=\frac{K}{2\hbar^2}(J^2-J_1^2-J_2^2).
\end{centeredalign}

Here \(K\) is an energy. Two spin-one-half particles then have a singlet energy \(-3K/4\) and a triplet energy \(K/4\). The gap is fixed by the coupling and the representation, without diagonalizing a general four-dimensional matrix. A magnetic field separates magnetic sublevels; an anisotropic interaction can also mix sectors that the isotropic model kept separate.

This calculation connects symmetry to a measurable excitation energy. Transferring the exchange formula to another object requires identifying which angular-momentum representation its degrees of freedom carry and whether anisotropy, orbital coupling or the environment changes the generator. A common algebra supplies the ladder relations, while the Hamiltonian selects which of those states a physical system occupies.

\subsection*{References For The Physical Derivation}
\begin{itemize}
\item \href{https://arxiv.org/abs/1508.06951}{V. Moretti, Mathematical Foundations of Quantum Mechanics: An Advanced Short Course; states, operators and symmetry.}
\end{itemize}



\clearpage
\section{Observable}
\label{topic:observable}

\subsection*{Mechanism}
Observable is the legal-question constructor: it turns a physical question into an operator with spectral outcome channels.

An observable converts a broad physical question such as position, energy, or spin into an operator with a defined domain and spectrum. Without that operator, the state alone does not specify which distribution is being predicted.

An observable is the mathematical form of a question that can be asked of a state. Its spectral decomposition defines the possible answers.

\subsection*{Physical Construction}
The state carrier is an admissible quantum state space on which the physical quantity is represented. The calculation involves a self-adjoint operator, operator-valued measure, or algebra element representing the physical question. Domain, self-adjointness, gauge invariance, and spectral conditions determine whether the quantity is a physical observable. The calculated quantities include eigenvalues, spectral measures, expectation values, moments, and response functions associated with the observable.

\subsection*{Topic Equations}
\begin{centeredalign}
A=A^\dagger\\
A=\sum_i a_iP_i\\
p(a_i)=\operatorname{Tr}(\rho P_i)
\end{centeredalign}

\subsection*{Physical Meaning}
The spectral measure decomposes the observable into possible outcome sectors. Pairing those sectors with the state gives probabilities, and weighting them by their spectral values gives expectation values. Degenerate eigenspaces correspond to outcomes that do not resolve every state component.

A spin state has different probability distributions for measurements along different axes. The preparation is unchanged; the observable selects the question.

The Born rule supplies the probability assigned to each spectral sector, and measurement theory describes its physical registration.

\subsection*{Invariance And Realization}
The rule connecting prepared states, observables, and spectral probability measures across wave, matrix, path-integral, circuit, or field notation. The operator-to-spectrum relation: admissible observations are represented through eigenvalues, projections, modes, or outcome channels. The dependence of admissible observable on measurement context or boundary condition. The non-commuting compatibility structure, which survives changes of representation.

The name of the carrier: particle, wave, field, qubit, or excitation. Where time dependence is represented: on the state, on the operator, or in a path weight. The coordinate system, basis, or geometric picture used to display the same relation. The physical implementation of detector, boundary, preparation, or observable.

\subsection*{Discriminating Consequences}
Self-adjointness, or the appropriate POVM positivity condition, is what makes the question a legal observable. A complete spectral resolution supplies all outcome channels for the question being asked. Equivalent representations preserve expectation values and probability distributions.




\clearpage
\section{Self-adjoint operator}
\label{topic:self-adjoint-operator}

\subsection*{Mechanism}
Self-adjoint operator is the admissible-observable condition: it gives real spectra and well-defined spectral measures.

A prediction concerns a stated physical question. The same state supports many incompatible questions, so a prediction requires an operator, spectral measure, or effect family in addition to the state itself.

Self-adjointness is not a technical decoration. It is the condition that makes an operator a legitimate spectral question in ordinary quantum mechanics.

\subsection*{Physical Construction}
The state carrier is an admissible quantum state space on which the physical quantity is represented. The calculation involves a self-adjoint operator, operator-valued measure, or algebra element representing the physical question. Domain, self-adjointness, gauge invariance, and spectral conditions determine whether the quantity is a physical observable. The calculated quantities include eigenvalues, spectral measures, expectation values, moments, and response functions associated with the observable.

\subsection*{Topic Equations}
\begin{centeredalign}
A=A^\dagger\\
A=\int_{\sigma(A)}\lambda\,dE_A(\lambda)\\
\Pr(\Delta)=\operatorname{Tr}(\rho E_A(\Delta))
\end{centeredalign}

\subsection*{Physical Meaning}
The operator's spectrum lists possible sharp values, while the state determines their weights. Matrix entries depend on basis, but the spectrum, expectation values, and probability distribution are unchanged by an equivalent representation. Domain and self-adjointness conditions are part of the physical definition.

An observable defines possible outcomes. The measurement chapter adds the probability rule and, when needed, the physical interaction that records one of those outcomes.

\subsection*{Invariance And Realization}
The rule connecting prepared states, observables, and spectral probability measures across wave, matrix, path-integral, circuit, or field notation. The operator-to-spectrum relation: admissible observations are represented through eigenvalues, projections, modes, or outcome channels. The dependence of admissible observable on measurement context or boundary condition. The non-commuting compatibility structure, which survives changes of representation.

The name of the carrier: particle, wave, field, qubit, or excitation. Where time dependence is represented: on the state, on the operator, or in a path weight. The coordinate system, basis, or geometric picture used to display the same relation. The physical implementation of detector, boundary, preparation, or observable.

\subsection*{Discriminating Consequences}
Self-adjointness, or the appropriate POVM positivity condition, is what makes the question a legal observable. A complete spectral resolution supplies all outcome channels for the question being asked. Equivalent representations preserve expectation values and probability distributions.




\clearpage

\section{Heisenberg group}

\label{topic:heisenberg-group}

A displacement of a particle's position followed by a momentum kick traces a rectangle in phase space when both operations are undone. Although the final position and momentum shifts cancel, the quantum amplitude acquires a phase. The Heisenberg group includes this phase as part of the composition law.

For canonical operators satisfying \([X,P]=i\hbar I\), define the displacement with position increment \(a\) and momentum increment \(b\) by \(D(a,b)=\exp[i(bX-aP)/\hbar]\). The commutator of the exponents is a scalar, so the Baker--Campbell--Hausdorff series stops after that commutator:

\begin{centeredalign}
D(a,b)D(a',b')=
e^{i(ba'-ab')/(2\hbar)}D(a+a',b+b').
\end{centeredalign}

The antisymmetric combination \(ba'-ab'\) is the oriented symplectic area. Its appearance explains why a description retaining only the two displacement coordinates has a projective composition law, whereas adding the phase gives an ordinary group representation. The central coordinate records how operations compose, rather than another position of the particle.

\begin{centeredalign}
D(a,0)D(0,b)D(-a,0)D(0,-b)=e^{-iab/\hbar}I.
\end{centeredalign}

A common overall phase has no effect on the probability of a single prepared state. To measure the loop phase, the displaced sequence must interfere with a reference sequence, for example through two coherent branches of a controlled operation. The observable is then a relative phase between the branches. A phase generated by a chosen path is compatible with fully specified dynamics.

The construction also relates particle motion to oscillator modes. Appropriate dimensionless quadratures express the same displacement algebra through creation and annihilation operators. Acting on the oscillator vacuum produces a coherent state; displacement changes its mean quadratures while preserving the vacuum covariance. The carrier and physical units differ, but the central commutation relation is the same mathematical relation being realized.

The uniqueness of the usual irreducible representation requires finitely many canonical pairs, regularity and a fixed nonzero central parameter. Quantum fields contain infinitely many modes and can admit inequivalent representations. Consequently the finite-dimensional symplectic picture cannot be promoted without qualification to an equivalence of arbitrary field-theory vacua. The group provides a precise reusable mechanism of composition within its stated representation assumptions.

\subsection*{References For The Physical Derivation}
\begin{itemize}
\item \href{https://arxiv.org/abs/1508.06951}{V. Moretti, Mathematical Foundations of Quantum Mechanics: An Advanced Short Course; states, operators and symmetry.}
\end{itemize}

\subsection*{Relations In The Original Papers}

\href{https://arxiv.org/html/1508.06951#S1.E21}{arXiv:1508.06951, S1.E21}. The central commutator supplies the infinitesimal relation used in the displacement composition law. This display establishes the infinitesimal relation, not by itself the global Weyl representation; the chapter states that representation explicitly.

\chapter{Noncommuting observables}

\begin{claimbox}

Two observables need not possess a common sharp spectral resolution. Their commutator and uncertainty relations quantify this algebraic incompatibility. Entanglement belongs instead to the structure of composite states, while Bell experiments join that state structure to local measurements.

\end{claimbox}

Noncommuting observables lack a common sharp spectral resolution. Commutators and uncertainty relations quantify this algebraic fact.

\section*{Topics}

\begin{longtable}{p{0.92\linewidth}}
\toprule
Page \\
\midrule
\endhead
\phantomsection\label{page:commutator}Commutator \\
\phantomsection\label{page:uncertainty-principle}Uncertainty principle \\
\phantomsection\label{page:canonical-commutation-relation}Canonical commutation relation \\
\bottomrule
\end{longtable}

\clearpage
\section{Commutator}
\label{topic:commutator}

\subsection*{Mechanism}
Commutator measures the physical consequence of exchanging two operations: it governs joint measurability, generated motion, and the leading difference between reversed protocols.

The commutator measures whether the order of two transformations matters. For observables it also tests whether a common eigenbasis, and hence a joint sharp description, can exist.

For two operators on a common domain, [A,B]=AB-BA compares the two possible orders. For observables, a nonzero commutator obstructs a common sharp spectral resolution and enters uncertainty bounds. For a Hamiltonian H, [H,A] gives the dynamical change of A. For short control pulses, [A,B] is the first term that distinguishes the two reversed sequences. One algebraic object therefore connects compatibility, dynamics, and order-sensitive response.

\subsection*{Physical Construction}
The state carrier is a common state space on which two transformations, observables, or questions are both defined. The calculation involves the ordered products AB and BA, compared through the obstruction [A,B]=AB-BA. A nonzero commutator marks an order-dependence or compatibility limit; a zero commutator permits a common sharp refinement only when the remaining spectral conditions hold. The calculated quantities include compatibility tests, uncertainty bounds, common eigenspaces, or canonical commutation relations.

\subsection*{Topic Equations}
\begin{centeredalign}
[A,B]=AB-BA\\
\frac{dA}{dt}=\frac{i}{\hbar}[H,A]+\left(\frac{\partial A}{\partial t}\right)\\
\Delta A\,\Delta B\geq\frac12\left|\langle[A,B]\rangle\right|\\
e^{\epsilon A}e^{\epsilon B}e^{-\epsilon A}e^{-\epsilon B}=I+\epsilon^2[A,B]+O(\epsilon^3)
\end{centeredalign}

\subsection*{Physical Meaning}
A zero commutator is a compatibility condition under the stated domains; a nonzero commutator identifies an algebraic obstruction. The value can also generate symmetry transformations, so the same operation links compatibility, dynamics, and conservation.

Position followed by momentum is not the same operation as momentum followed by position. Their canonical commutator fixes the lower bound in the corresponding uncertainty relation.

The uncertainty-principle page converts this algebraic obstruction into a state-dependent bound on simultaneous dispersion.

\subsection*{Consequences Forced By The Relation}
A nonzero commutator prevents a generic common sharp eigenbasis and imposes state-dependent uncertainty bounds. The commutator with the Hamiltonian determines whether an observable is conserved and how it changes in time. Reversing two short operations produces a difference proportional to their commutator at leading nontrivial order.
\subsection*{Domain Of The Construction}
For unbounded operators, the common domain belongs to the physical statement; a formal commutator without a domain need not define an observable relation. A nonzero commutator is not by itself a new interaction. It becomes physical through a state and an observable consequence.

\subsection*{Invariance And Realization}
The commutator transforms covariantly under a simultaneous unitary change of representation. Joint spectral compatibility is unchanged by a consistent representation change. The leading order-sensitive response survives when different physical controls realize the same operator algebra.

The matrix entries, basis, carrier, and physical implementation of the two operations may change. Domains can change the meaning of formal commutation relations for unbounded operators. The observable consequence depends on the prepared state and on how the operator difference is read out.

\subsection*{Discriminating Consequences}
Reversing two calibrated operations isolates the response generated by their commutator. A commuting control pair removes the order-sensitive contribution without changing the individual operations. Agreement requires the same operator domain and measured observable, not merely similar notation.

\subsection*{Equation Sources}
\begin{itemize}
\item \href{https://arxiv.org/abs/quant-ph/0005019}{arXiv:quant-ph/0005019}
\item \href{https://arxiv.org/abs/quant-ph/0001056}{arXiv:quant-ph/0001056}
\item \href{https://arxiv.org/abs/hep-ph/0104068}{arXiv:hep-ph/0104068}
\end{itemize}



\clearpage
\section{Uncertainty principle}
\label{topic:uncertainty-principle}

\subsection*{Mechanism}
Uncertainty principle is a compatibility-limit theorem: non-commuting observables impose lower bounds on joint sharpness.

The uncertainty principle states a quantitative consequence of incompatibility. It is not a claim about poor apparatus: it constrains the spreads of outcome distributions for every state in the relevant operator domains.

Uncertainty is not detector imperfection. It is a structural consequence of state variance and non-commuting observables.

\subsection*{Physical Construction}
The state carrier is one state space or a multipartite state space on which several questions can be asked. The calculation involves two or more observables, contexts, correlation operators, or hidden-variable assignments being compared. Commutators, uncertainty bounds, contextuality constraints, or Bell-type inequalities decide which joint assignments are possible. The calculated quantities include joint spectra, correlations, inequality violations, uncertainty products, or incompatible outcome statistics.

\subsection*{Topic Equations}
\begin{centeredalign}
\Delta_\rho A\,\Delta_\rho B\ge \frac12\left|\operatorname{Tr}(\rho[A,B])\right|\\
\Delta x\,\Delta p\ge\frac{\hbar}{2}
\end{centeredalign}

\subsection*{Physical Meaning}
The lower bound depends on the expectation value of the commutator and can be strengthened by covariance terms. A small spread in one observable therefore requires a compensating spread in its incompatible partner for the same prepared state.

Narrowing a particle's position wave packet broadens its momentum distribution because the two amplitudes are Fourier transforms, not because the momentum detector disturbs a pre-existing sharp value.

Multipartite correlations generalize incompatibility beyond pairs of observables and lead to entanglement and Bell tests.

\subsection*{Invariance And Realization}
The rule connecting prepared states, observables, and spectral probability measures across wave, matrix, path-integral, circuit, or field notation. The operator-to-spectrum relation: admissible observations are represented through eigenvalues, projections, modes, or outcome channels. The dependence of admissible observable on measurement context or boundary condition. The non-commuting compatibility structure, which survives changes of representation.

The name of the carrier: particle, wave, field, qubit, or excitation. Where time dependence is represented: on the state, on the operator, or in a path weight. The coordinate system, basis, or geometric picture used to display the same relation. The physical implementation of detector, boundary, preparation, or observable.

\subsection*{Discriminating Consequences}
A common eigenbasis or joint probability model exists only when the relevant compatibility conditions are satisfied. Commutators, uncertainty bounds, Bell inequalities, and contextuality tests are different forms of the same joint-observable obstruction. Classical joint-assignment failure is meaningful only under the explicitly stated locality, realism, or measurement-independence assumptions.




\clearpage

\section{Canonical commutation relation}

\label{topic:canonical-commutation-relation}

Translations in position and changes in momentum act differently when their order is reversed. For a particle on the real line, let \(X\) multiply a wave function by \(x\) and let \(P=-i\hbar\partial_x\). On smooth rapidly decreasing wave functions, direct differentiation gives

\begin{centeredalign}
[X,P]\psi=X(-i\hbar\partial_x\psi)+i\hbar\partial_x(x\psi)
=i\hbar\psi.
\end{centeredalign}

The additional term comes from differentiating the coordinate itself. It fixes the relation between the two operations, independently of the Hamiltonian subsequently chosen. With standard deviations defined in the prepared state and finite relevant moments, the Cauchy--Schwarz inequality applied to \((X-\langle X\rangle)\psi\) and \((P-\langle P\rangle)\psi\) then gives

\begin{centeredalign}
\Delta X\,\Delta P\geq\frac{\hbar}{2}.
\end{centeredalign}

This bound concerns the spread of outcomes in identically prepared ensembles. A measurement-disturbance experiment introduces an apparatus and a sequence of measurements and requires its own dynamical model. The preparation bound should not be substituted for that separate calculation.

The exponentiated relation makes domain issues more transparent. Let \(T(a)=\exp(-iaP/\hbar)\) translate a wave function by length \(a\), and \(B(b)=\exp(ibX/\hbar)\) shift its momentum by \(b\). Acting on a wave function in either order gives

\begin{centeredalign}
[T(a)\psi](x)=\psi(x-a),\qquad
T(a)B(b)=e^{-iab/\hbar}B(b)T(a).
\end{centeredalign}

The phase is proportional to the phase-space area enclosed by the two moves. It is common to the state for a single path, but can become a relative phase when the two operation sequences form coherent alternatives. This is the physical content of the central extension represented by the Heisenberg group.

No finite-dimensional matrices obey \([X,P]=i\hbar I\) exactly: taking a trace makes the left side zero and the right side nonzero. A truncated oscillator basis or a finite numerical grid must therefore depart from the relation somewhere. Such departures can be approximation effects rather than additional interactions. An exact construction records the representation and tests the states on which the commutator identity is being used.

\subsection*{References For The Physical Derivation}
\begin{itemize}
\item \href{https://arxiv.org/abs/1508.06951}{V. Moretti, Mathematical Foundations of Quantum Mechanics: An Advanced Short Course; states, operators and symmetry.}
\end{itemize}

\subsection*{Relations In The Original Papers}

\href{https://arxiv.org/html/1508.06951#S1.E21}{arXiv:1508.06951, S1.E21}. Position and momentum obey the canonical commutation relations on a common invariant domain. The surrounding derivation specifies Schwartz functions; the finite-matrix obstruction is calculated separately in the chapter.

\chapter{Measurement, instruments, and probabilities}

\begin{claimbox}

Measurement connects a state and an observable to recorded frequencies.  Projection, POVMs, Born weights, and collapse language are alternative ways of presenting this state-to-spectrum probability assignment.

\end{claimbox}

A quantum instrument joins an outcome probability to the corresponding state change. Bell experiments enter here because local measurements interrogate correlations in a composite state.

\section*{Topics}

\begin{longtable}{p{0.92\linewidth}}
\toprule
Page \\
\midrule
\endhead
\phantomsection\label{page:povm}POVM \\
\phantomsection\label{page:measurement-in-quantum-mechanics}Measurement in quantum mechanics \\
\phantomsection\label{page:quantum-jump}Quantum jump \\
\phantomsection\label{page:born-rule}Born rule \\
\phantomsection\label{page:measurement-problem}Measurement problem \\
\phantomsection\label{page:projection-valued-measure}Projection-valued measure \\
\phantomsection\label{page:quantum-amplifier}Quantum amplifier \\
\phantomsection\label{page:bell-s-theorem}Bell's theorem \\
\bottomrule
\end{longtable}

\clearpage
\section{POVM}
\label{topic:povm}

\subsection*{Mechanism}
POVM is the generalized-observable constructor: outcome effects need not be orthogonal projectors.

The measurement rule connects the formal state and observable to experimental frequencies. It distinguishes the probability assigned to an outcome from the conditional state change that may follow a recorded event.

POVMs separate the probability assignment from the idealized projection assumption. They are the natural mechanism for noisy, coarse-grained, indirect, or open-system measurements.

\subsection*{Physical Construction}
The state carrier is a state vector or density operator together with the measurement context in which outcome channels are defined. The calculation involves a projection-valued measure, POVM, update map, or instrument map connecting state to record. Outcome probabilities must be positive, normalized, and tied to a specified measurement map rather than to informal observer language. The calculated quantities include Born probabilities, detector records, post-measurement states, ensemble frequencies, or decision probabilities.

\subsection*{Topic Equations}
\begin{centeredalign}
E_i\ge0,\qquad \sum_iE_i=I\\
p(i)=\operatorname{Tr}(\rho E_i)\\
E_i=\sum_\alpha K_{i\alpha}^\dagger K_{i\alpha}
\end{centeredalign}

\subsection*{Physical Meaning}
Each positive effect represents an outcome channel and the effects sum to the identity, which enforces normalized probabilities. A projective measurement is a special case. A complete detector model may further specify a quantum instrument, whose maps describe both the outcome probability and the corresponding post-measurement state.

Once the probability rule is explicit, incompatibility can be tested rather than asserted. Commutators, uncertainty relations, and Bell-type constraints identify when several measurement questions cannot share one sharp assignment.

\subsection*{Invariance And Realization}
The rule connecting prepared states, observables, and spectral probability measures across wave, matrix, path-integral, circuit, or field notation. The operator-to-spectrum relation: admissible observations are represented through eigenvalues, projections, modes, or outcome channels. The dependence of admissible observable on measurement context or boundary condition. The non-commuting compatibility structure, which survives changes of representation.

The name of the carrier: particle, wave, field, qubit, or excitation. Where time dependence is represented: on the state, on the operator, or in a path weight. The coordinate system, basis, or geometric picture used to display the same relation. The physical implementation of detector, boundary, preparation, or observable.

\subsection*{Discriminating Consequences}
Outcome probabilities are non-negative and normalized because the observable acts on a valid state with a complete effect family. Projective measurement is the sharp limit of the same probability rule when effects become orthogonal projectors. The probability assignment is distinct from any optional post-measurement update convention.




\clearpage
\section{Measurement in quantum mechanics}
\label{topic:measurement-in-quantum-mechanics}

\subsection*{Mechanism}
Measurement in quantum mechanics is the complete measurement junction: it combines a state, a measurement model, probabilities, and sometimes an update rule.

Measurement theory must describe both an outcome distribution and the physical operation that produces a record. Keeping these roles separate prevents an interpretive account of state change from being mistaken for the probability law itself.

Measurement is not the root of quantum theory in this book. It is the junction where a prepared state and an observable or POVM are converted into probabilities and recorded outcomes.

\subsection*{Physical Construction}
The state carrier is a state vector or density operator together with the measurement context in which outcome channels are defined. The calculation involves a projection-valued measure, POVM, update map, or instrument map connecting state to record. Outcome probabilities must be positive, normalized, and tied to a specified measurement map rather than to informal observer language. The calculated quantities include Born probabilities, detector records, post-measurement states, ensemble frequencies, or decision probabilities.

\subsection*{Topic Equations}
\begin{centeredalign}
p(i)=\operatorname{Tr}(\rho E_i)\\
\rho\mapsto \rho_i=\frac{K_i\rho K_i^\dagger}{\operatorname{Tr}(K_i\rho K_i^\dagger)}\\
E_i=K_i^\dagger K_i
\end{centeredalign}

\subsection*{Physical Meaning}
A POVM gives outcome probabilities, whereas a quantum instrument gives the corresponding conditional transformations. Projective measurement is an ideal sharp limit. Real detectors are calibrated by showing that their effects are positive, complete, and consistent with observed frequencies.

A photon counter may report click or no click with non-unit efficiency. A two-effect POVM models those probabilities; the associated instrument is needed only when the state after the event matters.

The incompatibility chapter asks which families of such measurements can be jointly realized or assigned simultaneous sharp values.

\subsection*{Invariance And Realization}
The rule connecting prepared states, observables, and spectral probability measures across wave, matrix, path-integral, circuit, or field notation. The operator-to-spectrum relation: admissible observations are represented through eigenvalues, projections, modes, or outcome channels. The dependence of admissible observable on measurement context or boundary condition. The non-commuting compatibility structure, which survives changes of representation.

The name of the carrier: particle, wave, field, qubit, or excitation. Where time dependence is represented: on the state, on the operator, or in a path weight. The coordinate system, basis, or geometric picture used to display the same relation. The physical implementation of detector, boundary, preparation, or observable.

\subsection*{Discriminating Consequences}
Outcome probabilities are non-negative and normalized because the observable acts on a valid state with a complete effect family. Projective measurement is the sharp limit of the same probability rule when effects become orthogonal projectors. The probability assignment is distinct from any optional post-measurement update convention.




\clearpage

\section{Quantum jump}

\label{topic:quantum-jump}

A fluorescence detector records discrete photons, whereas an unobserved excited atom is often described by a smoothly decaying density operator. These descriptions refer to different information about the same atom and its radiation field. Conditioning on the detection record gives quantum trajectories; averaging over records recovers the master equation.

For Hamiltonian \(H\) and monitored decay operators \(L_k\), assume a Markovian reservoir and time-independent detection channels. In a short interval \(dt\), the no-count operator \(M_0\) and the count operators \(M_k\) are

\begin{centeredalign}
M_0=I-\frac{iH\,dt}{\hbar}-\frac{dt}{2}\sum_kL_k^\dagger L_k,
\qquad M_k=\sqrt{dt}\,L_k.
\end{centeredalign}

They satisfy \(\sum_{k\geq0}M_k^\dagger M_k=I+O(dt^2)\). The probability of a count in channel \(k\) is \(dt\,\operatorname{Tr}(L_k^\dagger L_k\rho)\); the state after that count is \(L_k\rho L_k^\dagger\) divided by its trace. Even the absence of a count changes the conditioned state, because it supplies information about its excited-state population.

\begin{centeredalign}
\dot\rho=-\frac{i}{\hbar}[H,\rho]+\sum_k
\left(L_k\rho L_k^\dagger-\frac12\{L_k^\dagger L_k,\rho\}\right).
\end{centeredalign}

This equation follows by summing the unnormalized states for all outcomes, expanding to first order in \(dt\) and subtracting the previous density operator. It is the ensemble prediction, not one particular observed sequence. For a two-level atom with \(L=\sqrt{\gamma}|g\rangle\langle e|\), no drive and an initially excited state, the survival probability is \(e^{-\gamma t}\) and the first-photon waiting-time density is \(\gamma e^{-\gamma t}\).

Changing the measurement of the outgoing field can change the conditioned trajectories while leaving the ensemble master equation unchanged. Direct photon counting gives jumps; quadrature detection gives a continuous noisy record. Consequently a trajectory belongs to the atom-field dynamics together with the measurement, not to the reduced Hamiltonian alone.

Finite efficiency separates the decay operator into observed and unobserved channels. A missed photon still affects the atom, but does not produce an observed jump. Feedback adds a further dependence: the detector record controls a subsequent Hamiltonian or operation. A mechanism assembled from decay, observation and feedback must retain these distinctions to predict both the mean population and the distribution of individual records.

\subsection*{References For The Physical Derivation}
\begin{itemize}
\item \href{https://arxiv.org/abs/quant-ph/9702007}{M. B. Plenio and P. L. Knight, The Quantum Jump Approach to Dissipative Dynamics in Quantum Optics; conditional and ensemble evolution.}
\end{itemize}

\subsection*{Relations In The Original Papers}

\href{https://arxiv.org/html/quant-ph/9702007#S4.E66}{arXiv:quant-ph/9702007, S4.E66}. Evolution conditioned on no detected emission is generated by an effective non-Hermitian Hamiltonian. The neighbouring equation gives the no-emission probability as the squared norm; detector assumptions determine the conditional evolution.

\clearpage
\section{Born rule}
\label{topic:born-rule}

\subsection*{Mechanism}
Born rule is the probability-observable constructor: it maps a state and a spectral channel to an observed probability.

The Born rule is the bridge from complex amplitudes to experimentally testable probabilities. It is an additional postulate: linear state evolution alone does not say how often a detector outcome should occur.

The Born rule is the point where the constructor becomes predictive. It does not name an object; it connects state preparation and a legal question to frequencies over outcome channels.

\subsection*{Physical Construction}
The state carrier is a state vector or density operator together with the measurement context in which outcome channels are defined. The calculation involves a projection-valued measure, POVM, update map, or instrument map connecting state to record. Outcome probabilities must be positive, normalized, and tied to a specified measurement map rather than to informal observer language. The calculated quantities include Born probabilities, detector records, post-measurement states, ensemble frequencies, or decision probabilities.

\subsection*{Topic Equations}
\begin{centeredalign}
p(i|\rho,\{P_i\})=\operatorname{Tr}(\rho P_i)\\
\Pr(X\in\Delta|\psi)=\int_\Delta|\psi(x)|^2\,d\mu(x)\\
\sum_i p(i)=1
\end{centeredalign}

\subsection*{Physical Meaning}
For a sharp measurement, the probability of an outcome is the state weight in the corresponding eigenspace. The density-operator form extends the same rule to mixtures and generalized measurements. Completeness of the outcome operators makes the probabilities sum to one.

A spin prepared equally between two vertical outcomes gives one-half for each vertical detector channel, even though each individual run records only one result.

Measurement theory adds detector coupling and conditional state change without altering this probability assignment.

\subsection*{Invariance And Realization}
The rule connecting prepared states, observables, and spectral probability measures across wave, matrix, path-integral, circuit, or field notation. The operator-to-spectrum relation: admissible observations are represented through eigenvalues, projections, modes, or outcome channels. The dependence of admissible observable on measurement context or boundary condition. The non-commuting compatibility structure, which survives changes of representation.

The name of the carrier: particle, wave, field, qubit, or excitation. Where time dependence is represented: on the state, on the operator, or in a path weight. The coordinate system, basis, or geometric picture used to display the same relation. The physical implementation of detector, boundary, preparation, or observable.

\subsection*{Discriminating Consequences}
Outcome probabilities are non-negative and normalized because the observable acts on a valid state with a complete effect family. Projective measurement is the sharp limit of the same probability rule when effects become orthogonal projectors. The probability assignment is distinct from any optional post-measurement update convention.




\clearpage
\section{Measurement problem}
\label{topic:measurement-problem}

\subsection*{Mechanism}
Measurement problem is the junction between unitary system--apparatus coupling, probability assignment, and conditional state update; these are distinct maps and need not be identified.

The measurement rule connects the formal state and observable to experimental frequencies. It distinguishes the probability assigned to an outcome from the conditional state change that may follow a recorded event.

A measurement model first couples the system to an apparatus or environment. A POVM or instrument then assigns outcome probabilities, and a conditional map specifies the post-record state. The foundational problem concerns the relation between these operations and a definite record, not the absence of a probability formula.

\subsection*{Physical Construction}
The state carrier is a joint system--apparatus state, possibly enlarged by environmental degrees of freedom. The calculation involves a premeasurement interaction followed by a measurement instrument whose components label possible records. The instrument maps are completely positive and their sum is trace preserving; the outcome effects sum to the identity. Outcome probabilities and conditional post-record states must be stated separately.

\subsection*{Topic Equations}
\begin{centeredalign}
\rho_{SA}'=U_{SA}(\rho_S\otimes\rho_A)U_{SA}^{\dagger}\\
p(i)=\operatorname{Tr}[\mathcal I_i(\rho_S)]=\operatorname{Tr}(\rho_SE_i)\\
\rho_{S|i}=\frac{\mathcal I_i(\rho_S)}{p(i)},\qquad \rho_S'=\sum_i\mathcal I_i(\rho_S)
\end{centeredalign}

\subsection*{Physical Meaning}
Each positive effect represents an outcome channel and the effects sum to the identity, which enforces normalized probabilities. A projective measurement is a special case. A complete detector model may further specify a quantum instrument, whose maps describe both the outcome probability and the corresponding post-measurement state.

Once the probability rule is explicit, incompatibility can be tested rather than asserted. Commutators, uncertainty relations, and Bell-type constraints identify when several measurement questions cannot share one sharp assignment.

\subsection*{Invariance And Realization}
The rule connecting prepared states, observables, and spectral probability measures across wave, matrix, path-integral, circuit, or field notation. The operator-to-spectrum relation: admissible observations are represented through eigenvalues, projections, modes, or outcome channels. The dependence of admissible observable on measurement context or boundary condition. The non-commuting compatibility structure, which survives changes of representation.

The name of the carrier: particle, wave, field, qubit, or excitation. Where time dependence is represented: on the state, on the operator, or in a path weight. The coordinate system, basis, or geometric picture used to display the same relation. The physical implementation of detector, boundary, preparation, or observable.

\subsection*{Discriminating Consequences}
Outcome probabilities are non-negative and normalized because the observable acts on a valid state with a complete effect family. Projective measurement is the sharp limit of the same probability rule when effects become orthogonal projectors. The probability assignment is distinct from any optional post-measurement update convention.




\clearpage
\section{Projection-valued measure}
\label{topic:projection-valued-measure}

\subsection*{Mechanism}
Projection-valued measure is the sharp-observable constructor: mutually exclusive outcome projectors partition the identity.

The measurement rule connects the formal state and observable to experimental frequencies. It distinguishes the probability assigned to an outcome from the conditional state change that may follow a recorded event.

A projection-valued measure encodes an ideal sharp measurement. It defines outcome channels that are orthogonal and exhaustive.

\subsection*{Physical Construction}
The state carrier is a state vector or density operator together with the measurement context in which outcome channels are defined. The calculation involves a projection-valued measure, POVM, update map, or instrument map connecting state to record. Outcome probabilities must be positive, normalized, and tied to a specified measurement map rather than to informal observer language. The calculated quantities include Born probabilities, detector records, post-measurement states, ensemble frequencies, or decision probabilities.

\subsection*{Topic Equations}
\begin{centeredalign}
P_iP_j=\delta_{ij}P_i,\qquad \sum_iP_i=I\\
p(i)=\operatorname{Tr}(\rho P_i)\\
\rho\mapsto \frac{P_i\rho P_i}{\operatorname{Tr}(\rho P_i)}
\end{centeredalign}

\subsection*{Physical Meaning}
Each positive effect represents an outcome channel and the effects sum to the identity, which enforces normalized probabilities. A projective measurement is a special case. A complete detector model may further specify a quantum instrument, whose maps describe both the outcome probability and the corresponding post-measurement state.

Once the probability rule is explicit, incompatibility can be tested rather than asserted. Commutators, uncertainty relations, and Bell-type constraints identify when several measurement questions cannot share one sharp assignment.

\subsection*{Invariance And Realization}
The rule connecting prepared states, observables, and spectral probability measures across wave, matrix, path-integral, circuit, or field notation. The operator-to-spectrum relation: admissible observations are represented through eigenvalues, projections, modes, or outcome channels. The dependence of admissible observable on measurement context or boundary condition. The non-commuting compatibility structure, which survives changes of representation.

The name of the carrier: particle, wave, field, qubit, or excitation. Where time dependence is represented: on the state, on the operator, or in a path weight. The coordinate system, basis, or geometric picture used to display the same relation. The physical implementation of detector, boundary, preparation, or observable.

\subsection*{Discriminating Consequences}
Outcome probabilities are non-negative and normalized because the observable acts on a valid state with a complete effect family. Projective measurement is the sharp limit of the same probability rule when effects become orthogonal projectors. The probability assignment is distinct from any optional post-measurement update convention.




\clearpage

\section{Quantum amplifier}

\label{topic:quantum-amplifier}

An amplifier that increases both quadratures of a bosonic signal must also introduce fluctuations. Let \(a\) be the input annihilation operator, with \([a,a^\dagger]=1\), and let \(G>1\) be the intensity gain. Multiplying \(a\) by \(\sqrt G\) alone would make the output commutator equal to \(G\), so it cannot describe an output mode with the original normalization.

A second, independent mode \(b\) supplies the required additional degree of freedom. The phase-insensitive amplification relation is

\begin{centeredalign}
a_{\rm out}=\sqrt G\,a_{\rm in}+\sqrt{G-1}\,b_{\rm in}^\dagger,
\qquad [a_{\rm out},a_{\rm out}^\dagger]=G-(G-1)=1.
\end{centeredalign}

The creation operator has the opposite commutator sign, exactly compensating the excess from the amplified signal. A parametric interaction realizes this transformation by producing correlated excitations in a signal and an idler. Energy is supplied by a pump; the idler is a physical channel rather than an adjustable numerical correction.

Define \(X=(a+a^\dagger)/\sqrt2\) and \(P=(a-a^\dagger)/(i\sqrt2)\), so vacuum has variance \(1/2\) in each quadrature. For uncorrelated inputs and a vacuum idler,

\begin{centeredalign}
\operatorname{Var}X_{\rm out}=G\operatorname{Var}X_{\rm in}+\frac{G-1}{2},
\qquad \langle n_{\rm out}\rangle=G\langle n_{\rm in}\rangle+G-1.
\end{centeredalign}

The second expression includes spontaneous output even for vacuum signal input. Referring the quadrature noise back to the input gives an added variance \((G-1)/(2G)\), approaching half a quantum at large gain. A thermally occupied idler adds still more noise. The commutator determines the necessary channel, while its state determines how much fluctuation that channel contributes.

A phase-sensitive amplifier follows a different relation: one quadrature is stretched and the conjugate one is compressed, preserving their commutator without amplifying both equally. It can amplify a known quadrature without the same added-noise bound. Thus the bound is attached to a specific physical task, not to every process called amplification.

This is an explicit construction from a failed relation. The proposed gain-only map has a nonzero commutator defect; adding an independent idler with the appropriate coefficient restores the algebra and predicts the noise. The correction is already established quantum optics. Its value here is that every additional ingredient changes a calculable experimental consequence.

\subsection*{References For The Physical Derivation}
\begin{itemize}
\item \href{https://arxiv.org/abs/1110.3234}{C. Weedbrook and colleagues, Gaussian Quantum Information; Gaussian channels and amplification.}
\end{itemize}

\subsection*{Relations In The Original Papers}

\href{https://arxiv.org/html/1110.3234v1#S5.E110}{arXiv:1110.3234v1, S5.E110}. Complete positivity ties the noise covariance to the gain of a one-mode Gaussian channel. The source uses vacuum quadrature variance one; the chapter uses one half. Its minimum noise follows after this normalization change.

\clearpage
\section{Bell's theorem}
\label{topic:bell-s-theorem}

\subsection*{Mechanism}
Bell's theorem is a compatibility/locality stress test: quantum correlations violate bounds satisfied by local hidden-variable assignments.

The measurement rule connects the formal state and observable to experimental frequencies. It distinguishes the probability assigned to an outcome from the conditional state change that may follow a recorded event.

Bell's theorem is not a page about a mysterious object. It is a falsifier for a classical joint-assignment model of measurement outcomes.

\subsection*{Physical Construction}
The state carrier is one state space or a multipartite state space on which several questions can be asked. The calculation involves two or more observables, contexts, correlation operators, or hidden-variable assignments being compared. Commutators, uncertainty bounds, contextuality constraints, or Bell-type inequalities decide which joint assignments are possible. The calculated quantities include joint spectra, correlations, inequality violations, uncertainty products, or incompatible outcome statistics.

\subsection*{Topic Equations}
\begin{centeredalign}
|E(a,b)+E(a,b')+E(a',b)-E(a',b')|\le 2\\
\text{quantum prediction can reach }2\sqrt2
\end{centeredalign}

\subsection*{Physical Meaning}
Each positive effect represents an outcome channel and the effects sum to the identity, which enforces normalized probabilities. A projective measurement is a special case. A complete detector model may further specify a quantum instrument, whose maps describe both the outcome probability and the corresponding post-measurement state.

Once the probability rule is explicit, incompatibility can be tested rather than asserted. Commutators, uncertainty relations, and Bell-type constraints identify when several measurement questions cannot share one sharp assignment.

\subsection*{Invariance And Realization}
The rule connecting prepared states, observables, and spectral probability measures across wave, matrix, path-integral, circuit, or field notation. The operator-to-spectrum relation: admissible observations are represented through eigenvalues, projections, modes, or outcome channels. The dependence of admissible observable on measurement context or boundary condition. The non-commuting compatibility structure, which survives changes of representation.

The name of the carrier: particle, wave, field, qubit, or excitation. Where time dependence is represented: on the state, on the operator, or in a path weight. The coordinate system, basis, or geometric picture used to display the same relation. The physical implementation of detector, boundary, preparation, or observable.

\subsection*{Discriminating Consequences}
Outcome probabilities are non-negative and normalized because the observable acts on a valid state with a complete effect family. Projective measurement is the sharp limit of the same probability rule when effects become orthogonal projectors. The probability assignment is distinct from any optional post-measurement update convention.




\part{Controlled And Spatial Realizations}

\chapter{Control sequences and quantum channels}

\begin{claimbox}

Quantum computing, channels, circuits, algorithms, networks, sensors, and error correction turn the same formal machinery into controlled sequences of operations.

\end{claimbox}

An ordered sequence becomes dynamics when its composition defines the implemented channel. Quantum control and information theory exploit this promotion directly.

\section*{Topics}

\begin{longtable}{p{0.92\linewidth}}
\toprule
Page \\
\midrule
\endhead
\phantomsection\label{page:quantum-error-correction}Quantum error correction \\
\phantomsection\label{page:quantum-channel}Quantum channel \\
\phantomsection\label{page:quantum-circuit}Quantum circuit \\
\phantomsection\label{page:quantum-key-distribution}Quantum key distribution \\
\phantomsection\label{page:quantum-teleportation}Quantum teleportation \\
\phantomsection\label{page:quantum-simulator}Quantum simulator \\
\phantomsection\label{page:quantum-metrology}Quantum metrology \\
\bottomrule
\end{longtable}

\clearpage

\section{Quantum error correction}

\label{topic:quantum-error-correction}

A quantum memory must distinguish errors without measuring the unknown logical amplitudes it is intended to preserve. The state is encoded into a subspace in which different correctable errors leave distinguishable information in auxiliary degrees of freedom, while their probabilities reveal no logical-state information.

Let \(P\) project onto the code subspace and let \(E_a\) be the error operators in a specified noise model. Exact correction of their linear span is possible when there is a matrix \(c\) such that

\begin{centeredalign}
P E_a^\dagger E_b P=c_{ab}P.
\end{centeredalign}

The right side is proportional to the identity within the code. Thus overlaps of the error-affected states depend on the error labels, but not on which logical superposition was stored. Diagonalizing \(c\) organizes the error spaces so that a measurement can identify the required correction without resolving the encoded state.

For a simple illustration, encode \(\alpha|0\rangle+\beta|1\rangle\) as \(\alpha|000\rangle+\beta|111\rangle\) and consider at most one bit flip. The commuting observables \(Z_1Z_2\) and \(Z_2Z_3\) give four distinct syndromes:

\begin{centeredalign}
\begin{array}{c|rrrr}
\text{error}&I&X_1&X_2&X_3\\\hline
(Z_1Z_2,Z_2Z_3)&(+,+)&(-,+)&(-,-)&(+,-)
\end{array}
\end{centeredalign}

Both basis codewords have the same syndrome before an error. The syndrome measurement therefore preserves their coherence. Applying the indicated bit flip restores the encoded state for any \(\alpha\) and \(\beta\). Measuring the individual \(Z_j\) instead would reveal which codeword was present and destroy an unknown superposition.

The noise specification matters. A phase flip acts within this code as a logical error and is not identified by those two syndromes. Protecting arbitrary single-qubit errors requires additional structure, such as a code that corrects both bit and phase errors. Increasing the number of spins without changing the encoded subspace and syndrome observables does not automatically supply that protection.

Error correction connects interaction, preparation, measurement and conditional dynamics in one mechanism. The environment determines the error operators; encoding determines the carrier subspace; syndrome extraction determines what information becomes observable; and the conditional operation restores the logical state. The algebraic correction condition is an exact test of this assembly. Fault tolerance adds another problem: the physical operations used to extract a syndrome can themselves propagate errors, so those operations must be analyzed as part of the noise model.

\subsection*{References For The Physical Derivation}
\begin{itemize}
\item \href{https://arxiv.org/abs/quant-ph/9705052}{D. Gottesman, Stabilizer Codes and Quantum Error Correction; error correction conditions and syndrome measurements.}
\end{itemize}

\subsection*{Relations In The Original Papers}

\href{https://arxiv.org/html/quant-ph/9705052#Ch2.E10}{arXiv:quant-ph/9705052, Ch2.E10}. Correctable errors have overlaps independent of the logical state inside the code space. A specified code and error set; the three-qubit example corrects single bit flips, not arbitrary single-qubit errors.

\href{https://arxiv.org/html/quant-ph/0302006#S3.E22}{arXiv:quant-ph/0302006, S3.E22}. A detected error can be reversed on the code when its norm does not distinguish logical codewords. Detected-emission feedback with a specified error record; it is narrower than correction of an arbitrary unobserved error set.

\clearpage
\section{Quantum channel}
\label{topic:quantum-channel}

\subsection*{Mechanism}
Quantum channel is the open-system protocol constructor: it maps input states to output states while preserving complete positivity and trace.

A quantum channel is the general transformation available to a state when an environment, uncontrolled degree of freedom, or measurement outcome is not retained. It extends unitary dynamics without abandoning positivity or probability conservation.

A channel is the mechanism for noisy transformations, measurements with forgotten outcomes, and subsystem evolution.

\subsection*{Physical Construction}
The state carrier is input and output density operators, possibly on different Hilbert spaces or subsystem carriers. The calculation involves a completely positive trace-preserving map, often represented by Kraus operators or by a Stinespring dilation. Complete positivity and trace preservation are the legal conditions; non-trace-preserving maps require an explicitly conditioned outcome. The calculated quantities include output state, final POVM probabilities, fidelity, capacity, error rate, or recovered subsystem statistics.

\subsection*{Topic Equations}
\begin{centeredalign}
\mathcal E(\rho)=\sum_a K_a\rho K_a^\dagger\\
\sum_aK_a^\dagger K_a=I\\
p(y)=\operatorname{Tr}(M_y\mathcal E(\rho))
\end{centeredalign}

\subsection*{Physical Meaning}
Complete positivity guarantees that the map remains physical when the input is entangled with an untouched system, and trace preservation guarantees total probability. A Kraus representation displays one realization, but different Kraus sets can describe the same channel.

Loss of phase coherence can be represented by a dephasing channel. Its action suppresses off-diagonal density-matrix elements while leaving the corresponding populations unchanged.

Quantum error correction asks whether information can be encoded so that a specified family of channels is detectable and reversible on the code space.

\subsection*{Invariance And Realization}
The rule connecting prepared states, observables, and spectral probability measures across wave, matrix, path-integral, circuit, or field notation. The operator-to-spectrum relation: admissible observations are represented through eigenvalues, projections, modes, or outcome channels. The dependence of admissible observable on measurement context or boundary condition. The non-commuting compatibility structure, which survives changes of representation.

The name of the carrier: particle, wave, field, qubit, or excitation. Where time dependence is represented: on the state, on the operator, or in a path weight. The coordinate system, basis, or geometric picture used to display the same relation. The physical implementation of detector, boundary, preparation, or observable.

\subsection*{Discriminating Consequences}
Each operation in the sequence is constrained by the map class it claims: unitary, completely positive, trace preserving, measurement, correction, or conditional update. The composed protocol is defined by its output state and outcome probabilities, not only by the names of the gates. Changing operation order or replacing a quantum channel with a classical control identifies which part of the protocol carries the effect.




\clearpage
\section{Quantum circuit}
\label{topic:quantum-circuit}

\subsection*{Mechanism}
Quantum circuit is the engineered-composition constructor: a finite sequence of admissible maps prepares, transforms, and measures a register.

A protocol is an ordered sequence of operations. Order is physical whenever the maps do not commute, so a list of available gates or channels is insufficient to define an algorithm, sensor, communication scheme, or correction cycle.

A circuit is the protocol layer of the same state-operator-observable machinery. Gates are controlled unitary or channel maps; measurement converts final states into output probabilities.

\subsection*{Physical Construction}
The state carrier is an input state, register, channel state, error syndrome, key, or controlled experimental configuration. The calculation involves an ordered sequence of gates, channels, measurements, corrections, encodings, or conditional maps. Each step must belong to the claimed map class: unitary, completely positive, trace-preserving, projective, conditional, or corrective. The calculated quantities include output state, key, error rate, fidelity, channel capacity, algorithmic success probability, or sensor estimate.

\subsection*{Topic Equations}
\begin{centeredalign}
\rho_{\mathrm{out}}=U_m\cdots U_2U_1\,\rho_{\mathrm{in}}\,U_1^\dagger U_2^\dagger\cdots U_m^\dagger\\
p(y)=\operatorname{Tr}(M_y\rho_{\mathrm{out}})
\end{centeredalign}

\subsection*{Physical Meaning}
The ordered composition carries a prepared input to a final state. Every intermediate map must preserve its stated physical conditions, and conditional operations are tied to explicit measurement outcomes. Performance is quantified through fidelity, error rate, capacity, precision, or success probability.

This is the executable end of the mechanism tree. It also closes the loop: failed predictions can be traced backward to the operation order, the generator, the state preparation, or the mathematical domain rather than attributed to the topic as a whole.

\subsection*{Invariance And Realization}
The rule connecting prepared states, observables, and spectral probability measures across wave, matrix, path-integral, circuit, or field notation. The operator-to-spectrum relation: admissible observations are represented through eigenvalues, projections, modes, or outcome channels. The dependence of admissible observable on measurement context or boundary condition. The non-commuting compatibility structure, which survives changes of representation.

The name of the carrier: particle, wave, field, qubit, or excitation. Where time dependence is represented: on the state, on the operator, or in a path weight. The coordinate system, basis, or geometric picture used to display the same relation. The physical implementation of detector, boundary, preparation, or observable.

\subsection*{Discriminating Consequences}
Each operation in the sequence is constrained by the map class it claims: unitary, completely positive, trace preserving, measurement, correction, or conditional update. The composed protocol is defined by its output state and outcome probabilities, not only by the names of the gates. Changing operation order or replacing a quantum channel with a classical control identifies which part of the protocol carries the effect.




\clearpage

\section{Quantum key distribution}

\label{topic:quantum-key-distribution}

Quantum key distribution creates correlated classical data whose secrecy is inferred from a specified quantum communication model. Its physical basis is that information about nonorthogonal signal states cannot be acquired perfectly while leaving those states unchanged. The task is key generation, not the transmission of an encrypted message or the authentication of an unknown partner.

In ideal BB84 the sender chooses between the eigenstates of \(Z\) and \(X\). Write \(|\pm\rangle=(|0\rangle\pm|1\rangle)/\sqrt2\). Suppose an eavesdropper's unitary interaction preserves both \(|0\rangle\) and \(|1\rangle\) while recording them in probe states \(|e_0\rangle\) and \(|e_1\rangle\). Linearity then requires

\begin{centeredalign}
|+\rangle|e\rangle\longmapsto
\frac{|0\rangle|e_0\rangle+|1\rangle|e_1\rangle}{\sqrt2}.
\end{centeredalign}

The signal remains a pure \(|+\rangle\) only when the two probe states coincide up to the phase compatible with the signal. Distinguishable probe records suppress its off-diagonal coherence. A check in the complementary basis can therefore detect the disturbance associated with that information gain.

After transmission, the legitimate parties disclose their basis choices, retain the compatible outcomes and estimate an error rate from a random sample. Classical error correction reconciles their strings and privacy amplification reduces the information that may remain with an adversary. For the ideal asymptotic single-photon BB84 setting with the relevant bit and phase error rates both bounded by \(Q\), the one-way secret fraction per sifted bit has the familiar form

\begin{centeredalign}
r\geq1-2h_2(Q),\qquad
h_2(Q)=-Q\log_2Q-(1-Q)\log_2(1-Q).
\end{centeredalign}

This expression is tied to its security model. Finite data require statistical confidence terms; multiphoton emission, detector imperfections and information disclosed during reconciliation require explicit treatment. An observed low bit-error rate alone does not establish secrecy for an arbitrary device. Loss and source statistics can carry information that a two-level idealization omits.

The classical channel must be authenticated. Otherwise an adversary can impersonate each party in a separate quantum exchange. The quantum mechanism supplies a relation between information and disturbance under stated source and measurement assumptions; authentication supplies the identity of the parties using that relation. Keeping these contributions distinct is necessary to transfer a security argument from an ideal qubit model to an optical implementation.

\subsection*{References For The Physical Derivation}
\begin{itemize}
\item \href{https://arxiv.org/abs/0802.4155}{V. Scarani and colleagues, The Security of Practical Quantum Key Distribution; ideal protocols, device assumptions and security bounds.}
\end{itemize}

\subsection*{Relations In The Original Papers}

\href{https://arxiv.org/html/0802.4155#S4.EGx43}{arXiv:0802.4155, S4.EGx43}. For the specified entanglement-based protocol, the asymptotic key rate subtracts phase-error information and error-correction leakage. Ideal single-photon assumptions and asymptotic rates; setting leakage to binary entropy gives the bound discussed in the chapter.

\clearpage

\section{Quantum teleportation}

\label{topic:quantum-teleportation}

An unknown qubit can be transferred using a shared entangled pair and two classical bits. The entangled resource alone does not transmit a usable signal: a joint measurement at the sender and a correction conditioned on its outcome are essential parts of the mechanism.

Let the input be \(|\psi\rangle=a|0\rangle+b|1\rangle\) and let the shared pair be \(|\Phi^+\rangle=(|00\rangle+|11\rangle)/\sqrt2\), with the second member held by the receiver. Define the Bell states \(|B_{mn}\rangle=(I\otimes X^nZ^m)|\Phi^+\rangle\), where \(m,n\in\{0,1\}\). Expansion in this joint basis gives

\begin{centeredalign}
|\psi\rangle_1|\Phi^+\rangle_{23}
=\frac12\sum_{m,n=0}^{1}|B_{mn}\rangle_{12}\,X^nZ^m|\psi\rangle_3.
\end{centeredalign}

Each Bell outcome has probability one quarter, independent of the input amplitudes. Once the sender communicates \(m,n\), the receiver applies \(Z^mX^n\) and obtains the original state. The measurement determines a known transformation of the state, rather than its unknown amplitudes. This is why transmitting two classical bits suffices when the entangled pair has already been supplied.

Without the outcome record the receiver averages the four possible transformed states. For any input density operator \(\rho\), that average is

\begin{centeredalign}
\frac14\sum_{m,n}X^nZ^m\rho Z^mX^n=\frac I2.
\end{centeredalign}

The receiver's local statistics then contain no dependence on the input. The same calculation establishes that the protocol cannot be used for faster-than-light signalling. The original input has participated in a destructive joint measurement, so the transfer does not create two independently available copies.

A nonideal entangled resource changes the resulting channel. For a Bell-diagonal resource, its Bell weights become probabilities of Pauli errors in the corrected output. Resource quality, Bell-measurement fidelity and the conditional correction can therefore be distinguished experimentally; all three can reduce the final state fidelity, but they enter the calculation at different steps.

Teleportation illustrates how a mechanism can be transferred across physical carriers. Photonic polarization, internal atomic levels and superconducting circuits may realize the same qubit relation, provided their entangling resource, joint measurement and correction implement the required maps. Changing only the name of the carrier leaves those requirements unresolved; demonstrating the maps establishes the actual correspondence.

\subsection*{References For The Physical Derivation}
\begin{itemize}
\item \href{https://arxiv.org/abs/quant-ph/9705052}{D. Gottesman, Stabilizer Codes and Quantum Error Correction; entanglement, measurements and teleportation.}
\item \href{https://arxiv.org/abs/1110.3234}{C. Weedbrook and colleagues, Gaussian Quantum Information; continuous-variable realizations and finite-resource effects.}
\end{itemize}

\subsection*{Relations In The Original Papers}

\href{https://arxiv.org/html/1110.3234v1#S4.E101}{arXiv:1110.3234v1, S4.E101}. Teleportation of coherent states with a finite two-mode squeezed resource has fidelity below one. A continuous-variable comparison to the exact ideal qubit protocol; it is not the derivation of the qubit Bell-state identity.

\clearpage
\section{Quantum simulator}
\label{topic:quantum-simulator}

\subsection*{Mechanism}
Quantum simulator is a target--carrier--validation construction: a controllable physical system encodes another model, and selected observables test whether the encoded dynamics is faithful.

A protocol is an ordered sequence of operations. Order is physical whenever the maps do not commute, so a list of available gates or channels is insufficient to define an algorithm, sensor, communication scheme, or correction cycle.

The simulator Hamiltonian is not by itself the target theory. The claim also needs an encoding between target and device states, a correspondence between their generators or channels, and validation observables with an error budget over the stated time and parameter range.

\subsection*{Physical Construction}
The state carrier is a controllable device state space together with an explicit encoding of the target state space. The calculation involves device Hamiltonians, channels, or gate sequences intended to reproduce target dynamics under the encoding. Control errors, leakage, finite size, noise, and approximation order define the regime in which the correspondence is claimed. The calculated quantities include encoded target observables compared with independently predicted or calibrated device measurements.

\subsection*{Topic Equations}
\begin{centeredalign}
V:\mathcal H_{\mathrm{target}}\hookrightarrow\mathcal H_{\mathrm{device}}\\
\left\|U_{\mathrm{device}}(t)V-VU_{\mathrm{target}}(t)\right\|\le\varepsilon(t)\\
\left|\langle O\rangle_{\mathrm{target}}-\langle VO V^{\dagger}\rangle_{\mathrm{device}}\right|\le\delta_O
\end{centeredalign}

\subsection*{Physical Meaning}
The ordered composition carries a prepared input to a final state. Every intermediate map must preserve its stated physical conditions, and conditional operations are tied to explicit measurement outcomes. Performance is quantified through fidelity, error rate, capacity, precision, or success probability.

This is the executable end of the mechanism tree. It also closes the loop: failed predictions can be traced backward to the operation order, the generator, the state preparation, or the mathematical domain rather than attributed to the topic as a whole.

\subsection*{Invariance And Realization}
The rule connecting prepared states, observables, and spectral probability measures across wave, matrix, path-integral, circuit, or field notation. The operator-to-spectrum relation: admissible observations are represented through eigenvalues, projections, modes, or outcome channels. The dependence of admissible observable on measurement context or boundary condition. The non-commuting compatibility structure, which survives changes of representation.

The name of the carrier: particle, wave, field, qubit, or excitation. Where time dependence is represented: on the state, on the operator, or in a path weight. The coordinate system, basis, or geometric picture used to display the same relation. The physical implementation of detector, boundary, preparation, or observable.

\subsection*{Discriminating Consequences}
Each operation in the sequence is constrained by the map class it claims: unitary, completely positive, trace preserving, measurement, correction, or conditional update. The composed protocol is defined by its output state and outcome probabilities, not only by the names of the gates. Changing operation order or replacing a quantum channel with a classical control identifies which part of the protocol carries the effect.




\clearpage

\section{Quantum metrology}

\label{topic:quantum-metrology}

A sensor estimates a physical parameter from how that parameter changes a prepared state. Sensitivity depends on the interaction that encodes the parameter, on the fluctuations of its generator in the preparation, and on the measurement used to distinguish nearby states. Entanglement is useful only when it improves this complete estimation problem.

Let a dimensionless parameter \(\theta\) be encoded by \(|\psi_\theta\rangle=e^{-i\theta G}|\psi_0\rangle\), with Hermitian dimensionless generator \(G\). For this pure-state unitary family, the quantum Fisher information is

\begin{centeredalign}
F_Q=4(\langle G^2\rangle-\langle G\rangle^2),\qquad
\operatorname{Var}\widehat\theta\geq\frac{1}{\nu F_Q}.
\end{centeredalign}

Here \(\nu\) counts independent repetitions and the bound applies to locally unbiased estimation, with attainability requiring a suitable measurement and statistical regime. Large generator variance makes neighboring parameter-dependent states more distinguishable. A state that is an eigenstate of \(G\) acquires only an overall phase and has no sensitivity to \(\theta\) in this task.

For \(N\) spins exposed to the same phase, choose \(G=\tfrac12\sum_j Z_j\). Independent spins prepared along \(x\) have \(F_Q=N\). The coherent superposition \((|0\rangle^{\otimes N}+|1\rangle^{\otimes N})/\sqrt2\) has \(F_Q=N^2\), because the two components acquire phases separated by \(N\theta\).

\begin{centeredalign}
\Delta\theta_{\rm product}\geq\frac{1}{\sqrt{\nu N}},\qquad
\Delta\theta_{\rm GHZ}\geq\frac{1}{N\sqrt\nu}.
\end{centeredalign}

These expressions compare the same single-spin coupling with a counted number of uses. They are not universal bounds for arbitrary many-body Hamiltonians or uncounted preparation resources. The enhanced oscillation also creates phase ambiguities over a broad prior interval, so an estimation scheme must establish which fringe contains the parameter.

Independent dephasing exposes the cost of the collective coherence. If a single-spin off-diagonal element decays as \(e^{-\gamma t}\), the coherence between the two GHZ branches decays as \(e^{-N\gamma t}\). For phase encoding over a fixed duration this gives \(F_Q=N^2e^{-2N\gamma t}\). Increasing \(N\) can then reduce rather than improve the usable information. Optimizing frequency estimation further requires counting the interrogation time and available repetitions.

A proposed sensing mechanism must therefore specify what is being estimated and which resources are fixed. The generator identifies the useful preparation, while environmental coupling and measurement determine how much of its distinguishability is available. This makes sensitivity a calculable property of the assembled experiment rather than a consequence of entanglement alone.

\subsection*{References For The Physical Derivation}
\begin{itemize}
\item \href{https://arxiv.org/abs/quant-ph/0509179}{V. Giovannetti, S. Lloyd and L. Maccone, Quantum metrology; generator fluctuations and counted resources.}
\end{itemize}

\subsection*{Relations In The Original Papers}

\href{https://arxiv.org/html/quant-ph/0509179#S0.EGx4}{arXiv:quant-ph/0509179, S0.EGx4}. An entangled probe gives an uncertainty bound scaling inversely with the number of uses of the parameter generator. Unitary encoding, counted resources and the estimator assumptions of the source; noise changes the attainable scaling.

\chapter{Boundaries and operator domains}

\begin{claimbox}

Walls, interfaces, asymptotic conditions, and media become part of the mechanism when they select the operator domain. The resulting self-adjoint extension fixes spectra, scattering channels, tunnelling amplitudes, and edge states.

\end{claimbox}

A boundary becomes physical closure when it fixes an operator domain. Tunnelling, confinement, scattering, and edge spectra follow from that domain.

\section*{Topics}

\begin{longtable}{p{0.92\linewidth}}
\toprule
Page \\
\midrule
\endhead
\phantomsection\label{page:particle-in-a-box}Particle in a box \\
\phantomsection\label{page:quantum-tunnelling}Quantum tunnelling \\
\phantomsection\label{page:quantum-harmonic-oscillator}Quantum harmonic oscillator \\
\bottomrule
\end{longtable}

\clearpage
\section{Particle in a box}
\label{topic:particle-in-a-box}

\subsection*{Mechanism}
Particle in a box is a boundary-spectrum constructor: a spatial domain and boundary condition discretize the allowed energy spectrum.

The particle in a box shows in the simplest form that a boundary condition is part of the Hamiltonian's physical definition. Confinement turns a continuous free-particle spectrum into discrete allowed energies.

The page shows how a boundary condition changes the domain of the Hamiltonian and therefore the allowed spectra.

\subsection*{Physical Construction}
The state carrier is Fock space, field configuration space, or a sector selected by charge, spin, momentum, statistics, or gauge data. The calculation involves creation, annihilation, field, charge, spin, Hamiltonian, or scattering operators acting on the admissible sector. Statistics, gauge constraints, commutation or anticommutation rules, domain conditions, and sector labels decide which states are legal. The calculated quantities include occupation number, charge, spin, momentum, energy, correlation function, cross-section, or scattering amplitude.

\subsection*{Topic Equations}
\begin{centeredalign}
\psi(0)=\psi(L)=0\\
\psi_n(x)=\sqrt{\frac{2}{L}}\sin\frac{n\pi x}{L}\\
E_n=\frac{\hbar^2\pi^2n^2}{2mL^2}
\end{centeredalign}

\subsection*{Physical Meaning}
The wave function satisfies the same differential equation inside the box as a free particle, but vanishing boundary values select standing waves. Only wavelengths fitting the interval are admissible, and their curvature fixes the quantized energies.

Doubling the box length reduces every level spacing by a factor of four. This follows from the boundary-selected wavelength and provides a direct check of the construction.

Replacing an impenetrable wall by a finite barrier leads to tunnelling, resonances, and scattering channels.

\subsection*{Invariance And Realization}
The rule connecting prepared states, observables, and spectral probability measures across wave, matrix, path-integral, circuit, or field notation. The operator-to-spectrum relation: admissible observations are represented through eigenvalues, projections, modes, or outcome channels. The dependence of admissible observable on measurement context or boundary condition. The non-commuting compatibility structure, which survives changes of representation.

The name of the carrier: particle, wave, field, qubit, or excitation. Where time dependence is represented: on the state, on the operator, or in a path weight. The coordinate system, basis, or geometric picture used to display the same relation. The physical implementation of detector, boundary, preparation, or observable.

\subsection*{Discriminating Consequences}
The boundary changes the operator domain, and therefore the allowed modes, transmission amplitudes, or scattering channels. The same operator can have different spectra when the admissible domain changes. Removing the boundary recovers the appropriate free, infinite-domain, or asymptotic limit.




\clearpage
\section{Quantum tunnelling}
\label{topic:quantum-tunnelling}

\subsection*{Mechanism}
Quantum tunnelling is a boundary-realization constructor: a state has nonzero transmission through a classically forbidden region.

Quantum tunnelling tests the consequence of wave evolution across a classically forbidden region. The effect depends jointly on the generator, barrier geometry, matching conditions, and incident state.

Tunnelling shows that the realization layer matters. A potential barrier changes the admissible wave solutions and produces transmission even where classical kinetic energy would be negative.

\subsection*{Physical Construction}
The state carrier is a Hilbert space with a selected domain, potential, interface, asymptotic channel, cavity, well, or boundary condition. The calculation involves a Hamiltonian, wave operator, transfer operator, or scattering map whose domain is changed by the boundary. Boundary conditions and matching conditions determine allowed states, resonances, transmission amplitudes, and spectra. The calculated quantities include eigenvalues, resonances, tunnelling probabilities, phase shifts, reflection/transmission amplitudes, or scattering data.

\subsection*{Topic Equations}
\begin{centeredalign}
T\sim \exp\!\left(-2\int_{x_1}^{x_2}\sqrt{\frac{2m(V(x)-E)}{\hbar^2}}\,dx\right)\\
-\frac{\hbar^2}{2m}\psi''(x)+V(x)\psi(x)=E\psi(x)
\end{centeredalign}

\subsection*{Physical Meaning}
Inside the barrier the stationary wave function is evanescent rather than oscillatory. Continuity and flux conditions connect it to incoming and outgoing waves, producing a nonzero transmission amplitude that falls approximately exponentially with barrier width in the semiclassical regime.

Changing the barrier width while holding its height and the incident energy fixed isolates the predicted exponential dependence and separates tunnelling from an over-barrier contribution.

Scattering theory generalizes the same matching construction to many incoming, outgoing, and resonant channels.

\subsection*{Invariance And Realization}
The rule connecting prepared states, observables, and spectral probability measures across wave, matrix, path-integral, circuit, or field notation. The operator-to-spectrum relation: admissible observations are represented through eigenvalues, projections, modes, or outcome channels. The dependence of admissible observable on measurement context or boundary condition. The non-commuting compatibility structure, which survives changes of representation.

The name of the carrier: particle, wave, field, qubit, or excitation. Where time dependence is represented: on the state, on the operator, or in a path weight. The coordinate system, basis, or geometric picture used to display the same relation. The physical implementation of detector, boundary, preparation, or observable.

\subsection*{Discriminating Consequences}
The boundary changes the operator domain, and therefore the allowed modes, transmission amplitudes, or scattering channels. The same operator can have different spectra when the admissible domain changes. Removing the boundary recovers the appropriate free, infinite-domain, or asymptotic limit.




\clearpage

\section{Quantum harmonic oscillator}

\label{topic:quantum-harmonic-oscillator}

Small oscillations about a stable equilibrium have a quadratic energy. Quantizing that local approximation produces equally spaced excitations and a nonzero ground-state variance. For mass \(m\), angular frequency \(\omega>0\), position \(X\) and momentum \(P\), take

\begin{centeredalign}
H=\frac{P^2}{2m}+\frac12m\omega^2X^2,\qquad
a=\sqrt{\frac{m\omega}{2\hbar}}X+\frac{iP}{\sqrt{2m\hbar\omega}}.
\end{centeredalign}

The canonical commutator gives \([a,a^\dagger]=1\). Substituting \(a\) and \(a^\dagger\) into the energy leaves an additive term from their noncommutation:

\begin{centeredalign}
H=\hbar\omega\left(a^\dagger a+\frac12\right),\qquad
E_n=\hbar\omega\left(n+\frac12\right).
\end{centeredalign}

Positivity of \(a^\dagger a\) bounds the energy below. Repeated lowering reaches a state annihilated by \(a\); raising generates the nonnegative integer occupations. The half-quantum is consequently fixed by the operator algebra. It is not a fitted offset introduced to match a spectrum.

The ground state has position variance \(\hbar/(2m\omega)\) and momentum variance \(m\hbar\omega/2\). Their product saturates the preparation uncertainty relation. Displacing that state changes the mean position and momentum while leaving the variances fixed. Under the harmonic Hamiltonian, those means obey the classical oscillator equation, even though the state retains quantum fluctuations.

\begin{centeredalign}
X(t)=X(0)\cos\omega t+\frac{P(0)}{m\omega}\sin\omega t.
\end{centeredalign}

The equation is an operator identity. Taking expectations gives the classical motion of the center; taking products gives correlations and variances. A thermal mixture has zero mean displacement but fluctuating energy and a larger position variance. A squeezed state instead redistributes variance between the two quadratures, and that anisotropy rotates in phase space.

The same construction describes a normal mode of a vibrating solid or electromagnetic cavity after its generalized coordinate, conjugate momentum and normalization have been identified. The material and boundary conditions determine the frequency and mode shape. Anharmonic interactions couple modes or shift their spacing, while damping adds an environment and its fluctuations. The oscillator is therefore a transferable quadratic mechanism, not an assertion that all its realizations have identical lifetimes, couplings or measurement units.

For a finite matrix truncation, the highest retained state breaks the exact ladder algebra. This can be checked directly rather than interpreted as a physical correction. Predictions for low occupations converge as the truncation is enlarged; strong driving that reaches the cutoff requires a larger state space.

\subsection*{References For The Physical Derivation}
\begin{itemize}
\item \href{https://arxiv.org/abs/1110.3234}{C. Weedbrook and colleagues, Gaussian Quantum Information; canonical quadratures, Gaussian states and transformations.}
\end{itemize}

\subsection*{Relations In The Original Papers}

\href{https://arxiv.org/html/1110.3234v1#S2.E3}{arXiv:1110.3234v1, S2.E3}. The annihilation operator lowers an oscillator number state with amplitude equal to the square root of its occupation. An infinite bosonic Fock space; finite truncations change the commutator at the upper boundary.

\part{Fields, Constraints, And Scale}

\chapter{Fields, constraints, and scale}

\begin{claimbox}

Quantum field theory promotes fields to operator-valued degrees of freedom. Gauge constraints select the physical Hilbert space, particle statistics select admissible sectors, and renormalization makes the effective law depend on scale.

\end{claimbox}

Fields enlarge the state space to variable occupation and local operator algebras. Gauge constraints, statistics, and scale then select the physical sectors and effective law.

\section*{Topics}

\begin{longtable}{p{0.92\linewidth}}
\toprule
Page \\
\midrule
\endhead
\phantomsection\label{page:gauge-theory}Gauge theory \\
\phantomsection\label{page:photon}Photon \\
\phantomsection\label{page:renormalization}Renormalization \\
\phantomsection\label{page:quantum-field-theory}Quantum field theory \\
\phantomsection\label{page:fermi-dirac-statistics}Fermi–Dirac statistics \\
\phantomsection\label{page:quantum-electrodynamics}Quantum electrodynamics \\
\phantomsection\label{page:boson}Boson \\
\phantomsection\label{page:quantum-chromodynamics}Quantum chromodynamics \\
\phantomsection\label{page:fermion}Fermion \\
\phantomsection\label{page:electron}Electron \\
\phantomsection\label{page:ads-cft-correspondence}AdS/CFT correspondence \\
\phantomsection\label{page:bose-einstein-statistics}Bose–Einstein statistics \\
\phantomsection\label{page:quantum-statistical-mechanics}Quantum statistical mechanics \\
\phantomsection\label{page:quantum-gravity}Quantum gravity \\
\phantomsection\label{page:quantum-geometry}Quantum geometry \\
\phantomsection\label{page:fock-space}Fock space \\
\bottomrule
\end{longtable}

\clearpage
\section{Gauge theory}
\label{topic:gauge-theory}

\subsection*{Mechanism}
Gauge theory defines how internal states are compared at different spacetime points: a connection relates neighboring frames, and curvature records the path dependence that no single gauge choice can remove.

Gauge theory separates physical states from multiple mathematical descriptions related by local transformations. The redundancy is useful because it makes locality and interaction structure explicit, but constraints are required to remove unphysical degrees of freedom.

A local gauge transformation changes the field coordinates used at each point without changing the physical state. Ordinary derivatives compare fields in different local frames and therefore cease to transform covariantly. The gauge connection repairs that comparison. Its commutator gives the field strength, while Wilson loops measure the accumulated transport around a closed path. Gauss constraints select physical states and charges. The connection is representation dependent; curvature, loop observables, and gauge-invariant amplitudes carry the physical content.

\subsection*{Physical Construction}
The state carrier is matter and gauge fields modulo local gauge equivalence, restricted to the physical constraint sector. The calculation involves a covariant derivative and connection whose commutator gives the field strength. Gauss constraints, gauge covariance, operator domains, and boundary conditions select the physical states and charges. The calculated quantities include gauge-invariant contractions of field strengths, traced Wilson loops, physical charges, and scattering amplitudes.

\subsection*{Topic Equations}
\begin{centeredalign}
D_\mu=\partial_\mu+igA_\mu\\
[D_\mu,D_\nu]=igF_{\mu\nu}\\
F_{\mu\nu}'=U F_{\mu\nu}U^{-1}\\
W(\gamma)=\operatorname{Tr}\,\mathcal P\exp\!\left(-ig\oint_\gamma A_\mu dx^\mu\right)\\
G^a\ket{\Psi_{\mathrm{phys}}}=0
\end{centeredalign}

\subsection*{Physical Meaning}
Gauge-related field configurations represent the same physical state. Observable quantities must therefore be gauge invariant, or transform covariantly within a construction whose final predictions are invariant. Gauge fixing chooses one representative without changing the physical equivalence class.

Changing the electromagnetic scalar and vector potentials by a gauge transformation changes their formulas but leaves electric and magnetic fields, phase-consistent amplitudes, and measured forces unchanged.

Renormalization adds a second kind of transformation: changing the scale at which the same theory is parametrized.

\subsection*{Consequences Forced By The Relation}
In a non-Abelian theory the field-strength matrix changes by conjugation. Traces of its products and closed Wilson loops are invariant. For electromagnetism the group is Abelian, so the field strength itself is unchanged. A flat connection can still have nontrivial holonomy on a multiply connected domain. The connection and the topology together determine this global information. The Gauss constraint ties matter charge to the divergence of the electric field. Transformations that act nontrivially at a boundary can carry charges and must be distinguished from gauge redundancies that vanish there.
\subsection*{Domain Of The Construction}
The gauge group, representation, matter content, dimension, and boundary conditions distinguish different physical theories. A measured interference phase includes the connection along the paths and the phases of the charged states. The traced holonomy is gauge invariant; an untraced non-Abelian holonomy transforms by conjugation at its base point.

\subsection*{Invariance And Realization}
Gauge-equivalent potentials give the same physical probabilities, invariant contractions of the field strength, and traced closed-loop observables. Curvature records the infinitesimal holonomy of the connection and cannot be removed by a local gauge choice. The physical state belongs to the constraint sector rather than to an arbitrary field-coordinate representation.

The gauge potential, local basis, gauge-fixing condition, and coordinate description may change. The gauge group, representation, matter content, dimension, and boundary conditions specify different physical theories. Topological sectors and boundary charges can survive even where the local field strength vanishes.

\subsection*{Discriminating Consequences}
A closed-loop phase or Wilson observable distinguishes nontrivial holonomy from a removable local gauge choice. Gauge-related descriptions give identical probabilities for the same physical preparation and observable. An additional field or interaction changes a gauge-invariant observable rather than only the gauge potential.

\subsection*{Equation Sources}
\begin{itemize}
\item \href{https://arxiv.org/abs/hep-th/0106242}{arXiv:hep-th/0106242}
\item \href{https://arxiv.org/abs/hep-th/0010277}{arXiv:hep-th/0010277}
\end{itemize}
\subsection*{References For The Physical Derivation}
\begin{itemize}
\item \href{https://arxiv.org/abs/hep-ph/0001312}{G. S. Bali, QCD forces and heavy quark bound states, Appendix B: covariant derivative, curvature and gauge transformations.}
\end{itemize}


\clearpage
\section{Photon}
\label{topic:photon}

\subsection*{Mechanism}
Photon is a field-mode constructor: a one-quantum excitation of the electromagnetic field, constrained by massless dispersion and transverse polarization.

Relativistic, many-body, field, gauge, geometric, and scale-dependent settings each need their own state space and operator domain, specified for that setting rather than inferred from a single-particle model.

The native photon mechanism is field quantization. The electromagnetic field is decomposed into modes, creation and annihilation operators act on those modes, and a one-photon state is created from the vacuum. The relevant observables are occupation number, energy, momentum, polarization, and detector clicks; the constraints are dispersion and gauge-compatible transversality.

\subsection*{Physical Construction}
The state carrier is Fock space, field configuration space, or a sector selected by charge, spin, momentum, statistics, or gauge data. The calculation involves creation, annihilation, field, charge, spin, Hamiltonian, or scattering operators acting on the admissible sector. Statistics, gauge constraints, commutation or anticommutation rules, domain conditions, and sector labels decide which states are legal. The calculated quantities include occupation number, charge, spin, momentum, energy, correlation function, cross-section, or scattering amplitude.

\subsection*{Topic Equations}
\begin{centeredalign}
E=\hbar\omega,\qquad \mathbf p=\hbar\mathbf k,\qquad \omega=c|\mathbf k|\\
\ket{1_{\mathbf k,\lambda}}=a_{\mathbf k,\lambda}^{\dagger}\ket{0}\\
\hat N_{\mathbf k,\lambda}=a_{\mathbf k,\lambda}^{\dagger}a_{\mathbf k,\lambda}\\
\mathbf k\cdot\boldsymbol\epsilon_{\mathbf k,\lambda}=0
\end{centeredalign}

\subsection*{Physical Meaning}
Different topics in this branch use different carriers: spinor wave functions, Fock spaces, many-body states, gauge sectors, geometric states, or effective low-energy sectors. Their physical content is fixed by the associated field equation or Hamiltonian, its constraints and domain, and the amplitudes, spectra, charges, or correlation functions it predicts.

Field and many-body mechanisms become experimentally useful when assembled into an ordered intervention. The protocol chapter shows how preparation, controlled evolution, measurement, and correction compose into one executable map.

\subsection*{Invariance And Realization}
The relation between prepared states, observables, and spectral probability measures. The use of eigenvalues, projectors, modes, or outcome channels to represent admissible observations. The dependence of the observable on basis, domain, potential, preparation, or measurement context. The commutator structure that limits which observables can be jointly diagonalized.

The physical carrier: particle, wave, field mode, spin, qubit, detector, or excitation. The representation: wave mechanics, matrix mechanics, density matrices, path integrals, circuits, or fields. Where time dependence is placed: on the state, on the operator, in a propagator, or in a path weight. The implementation of preparation, boundary condition, detector, or outcome channel.

\subsection*{Discriminating Consequences}
A transfer target provides a state space, a transformation law, and a spectral or categorical observable, with one compatibility relation experimentally unresolved. A useful validation varies the basis, domain, or measurement context and measures whether the allowed observable changes while the underlying transformation law remains identifiable. A stronger validation contains two candidate observables whose predicted commutator controls joint resolvability.




\clearpage
\section{Renormalization}
\label{topic:renormalization}

\subsection*{Mechanism}
Renormalization relates physical descriptions at different resolutions by changing couplings and operator weights while preserving long-distance predictions.

Renormalization explains how a theory preserves predictions while its effective parameters and degrees of freedom change with scale. It is therefore a transformation between descriptions, not merely a device for removing infinities.

Coarse graining removes short-distance variables and generates every operator allowed by the remaining symmetries. Their coefficients flow with scale. Relevant directions grow, irrelevant directions decay, and fixed points organize scale-invariant behaviour. Microscopically different systems can therefore share critical exponents and scaling functions when their flows approach the same fixed point. Universality is the invariance class of this scale transformation, not a visual similarity between equations.

\subsection*{Physical Construction}
The state carrier is effective fields and degrees of freedom defined at a stated resolution or cutoff. The calculation involves a coarse-graining transformation and beta functions acting on the coefficients of an effective operator expansion. Symmetry, dimensionality, locality assumptions, relevant directions, and renormalization conditions determine the allowed flow. The calculated quantities include running couplings, correlation functions, scaling dimensions, critical exponents, and corrections to scaling.

\subsection*{Topic Equations}
\begin{centeredalign}
\mu\frac{dg_i}{d\mu}=\beta_i(\{g\})\\
\left(\mu\partial_\mu+\beta_i\partial_{g_i}+n\gamma\right)G^{(n)}=0\\
\beta_i(g_*)=0,\qquad \delta g_i(b)=b^{y_i}\delta g_i\\
\mathcal L_{\mathrm{eff}}(\mu)=\sum_i c_i(\mu)\mathcal O_i
\end{centeredalign}

\subsection*{Physical Meaning}
A change of scale moves the couplings along a flow. Predictions remain fixed when explicit scale dependence and coupling dependence compensate. Fixed points and relevant directions then organize universality across microscopically different systems.

Different lattice models can approach the same critical exponents because coarse-graining removes microscopic details while preserving the long-distance transformation structure.

Effective field theory uses this scale organization to decide which operators must be retained for a specified accuracy.

\subsection*{Consequences Forced By The Relation}
Observable predictions are independent of the arbitrary renormalization scale when explicit scale dependence and coupling flow compensate. Systems approaching the same fixed point with the same relevant directions share critical exponents and scaling functions. Relevant perturbations drive the flow away from a fixed point and select the accessible phases or crossover regimes.
\subsection*{Domain Of The Construction}
Universality requires matching dimensionality, symmetries, conserved quantities, interaction range, and relevant directions; a shared dimensionless group alone is insufficient. An effective theory is tied to a scale and accuracy, because omitted operators return through controlled corrections.

\subsection*{Invariance And Realization}
Observable predictions remain independent of the arbitrary renormalization scale when couplings and fields flow consistently. Critical exponents and scaling functions are shared by systems approaching the same fixed point with the same relevant directions. Symmetry and dimensionality constrain the effective operators generated by coarse graining.

Couplings, field normalizations, effective degrees of freedom, and operator coefficients depend on scale. Microscopic Hamiltonians can differ while their long-distance flows enter the same universality class. A relevant perturbation can drive the system away from one fixed point toward another phase or crossover regime.

\subsection*{Discriminating Consequences}
Measurements at several scales determine whether running couplings follow one beta function. Microscopically different realizations determine whether critical exponents and scaling functions coincide. Corrections to scaling distinguish approach to a fixed point from an exact scale-invariant law.

\subsection*{Equation Sources}
\begin{itemize}
\item \href{https://arxiv.org/abs/hep-ph/0106305}{arXiv:hep-ph/0106305}
\item \href{https://arxiv.org/abs/hep-th/0104197}{arXiv:hep-th/0104197}
\end{itemize}



\clearpage
\section{Quantum field theory}
\label{topic:quantum-field-theory}

\subsection*{Mechanism}
Quantum field theory is the many-mode local-field extension of the quantum constructor.

Quantum field theory is needed when relativity, locality, and particle creation must be described together. Fields become the primary operator-valued objects, while particles appear as excitations of their modes.

Quantum field theory lifts the state-operator-spectrum construction to fields, local operators, creation and annihilation modes, and scattering amplitudes. Particles become stable excitation/observable roles of fields.

\subsection*{Physical Construction}
The state carrier is Fock space, field configuration space, or a sector selected by charge, spin, momentum, statistics, or gauge data. The calculation involves creation, annihilation, field, charge, spin, Hamiltonian, or scattering operators acting on the admissible sector. Statistics, gauge constraints, commutation or anticommutation rules, domain conditions, and sector labels decide which states are legal. The calculated quantities include occupation number, charge, spin, momentum, energy, correlation function, cross-section, or scattering amplitude.

\subsection*{Topic Equations}
\begin{centeredalign}
\Phi(x)=\sum_k\left(a_k u_k(x)+a_k^\dagger u_k^*(x)\right)\\
[a_k,a_l^\dagger]_{\mp}=\delta_{kl}\\
\langle 0|T\{\Phi(x)\Phi(y)\}|0\rangle
\end{centeredalign}

\subsection*{Physical Meaning}
The state lives in a many-particle or field space, local operators create correlations, and the action or Hamiltonian generates their evolution. Perturbative amplitudes are expansions of these operator relations, not independent pictorial rules.

A photon is represented as an excitation of the electromagnetic field. Processes that change photon number are natural in Fock space but cannot be represented in a fixed one-particle Hilbert space.

Gauge theory identifies redundant field descriptions, and renormalization connects parameters measured at different scales.

\subsection*{Invariance And Realization}
The rule connecting prepared states, observables, and spectral probability measures across wave, matrix, path-integral, circuit, or field notation. The operator-to-spectrum relation: admissible observations are represented through eigenvalues, projections, modes, or outcome channels. The dependence of admissible observable on measurement context or boundary condition. The non-commuting compatibility structure, which survives changes of representation.

The name of the carrier: particle, wave, field, qubit, or excitation. Where time dependence is represented: on the state, on the operator, or in a path weight. The coordinate system, basis, or geometric picture used to display the same relation. The physical implementation of detector, boundary, preparation, or observable.

\subsection*{Discriminating Consequences}
Commutation, anticommutation, gauge, and occupation rules define which many-mode states are admissible. The field or many-mode construction must reduce to the appropriate single-particle, quasiparticle, or low-energy limit when those limits exist. Dual presentations are credible only when observables, spectra, or correlation functions are preserved across the translation.




\clearpage
\section{Fermi–Dirac statistics}
\label{topic:fermi-dirac-statistics}

\subsection*{Mechanism}
Fermi–Dirac statistics is an antisymmetric-state constructor: exchange statistics restricts which many-fermion occupation patterns are admissible.

Fermi-Dirac statistics is needed because identical fermions do not occupy many-particle states independently. Antisymmetry under particle exchange, and the exclusion principle that follows from it, changes the allowed occupation patterns before any Hamiltonian dynamics is considered.

Fermi-Dirac statistics is an admissibility rule, not a generator of time evolution. Anticommutation and antisymmetry imply single occupation of each one-particle mode, while the thermal distribution states the mean occupation of those modes.

\subsection*{Physical Construction}
The state carrier is a fermionic Fock space assembled from antisymmetric many-particle sectors or occupation-number modes. The calculation involves creation, annihilation, and number operators obeying canonical anticommutation relations. Exchange antisymmetry restricts each one-particle mode to occupation zero or one for each internal state. The calculated quantities include mode occupations, Fermi energy, particle density, pressure, heat capacity, and other equilibrium response functions.

\subsection*{Topic Equations}
\begin{centeredalign}
\{a_i,a_j^\dagger\}=\delta_{ij},\qquad n_i\in\{0,1\}\\
\bar n_i=\frac{1}{e^{\beta(\varepsilon_i-\mu)}+1}
\end{centeredalign}

\subsection*{Physical Meaning}
Single-particle energy levels provide the available modes, while each mode can be occupied at most once per internal state. The Fermi-Dirac distribution gives the mean occupation at thermal equilibrium. At zero temperature it fills modes up to the Fermi energy.

The degeneracy pressure of an electron gas follows from filling successively higher momentum states even without a repulsive force between the electrons. The effect is a consequence of antisymmetric state construction.

Bose-Einstein statistics changes the exchange rule from antisymmetric to symmetric and therefore permits unlimited occupation of one mode.

\subsection*{Invariance And Realization}
Exchange antisymmetry and canonical anticommutation define the fermionic many-particle sectors. Each one-particle mode has occupation zero or one for each internal state. The Fermi-Dirac function gives the equilibrium mean occupation once energy, temperature, and chemical potential are specified.

The dispersion relation, dimensionality, degeneracy, density of states, and interaction approximation depend on the physical system. Electrons, atoms, nucleons, and fermionic quasiparticles realize the same exchange rule with different Hamiltonians and observables. Finite temperature smooths the occupation edge that is sharp at the Fermi energy in the ideal zero-temperature limit.

\subsection*{Discriminating Consequences}
The theory recovers occupations restricted to zero or one and the ideal zero-temperature Fermi sea. The theory recovers the Maxwell-Boltzmann distribution in the dilute low-fugacity limit. Integrate the mode occupations against the density of states and verify the specified particle number.




\clearpage
\section{Quantum electrodynamics}
\label{topic:quantum-electrodynamics}

\subsection*{Mechanism}
An electromagnetic field can exchange energy and momentum with charged particles. Quantum electrodynamics describes both the charged Dirac field and the electromagnetic field as dynamical quantum degrees of freedom.

Treating the potential as prescribed describes motion in an external field. Allowing the field to fluctuate permits emission, absorption and virtual photon exchange. The additional dynamical variables change which processes the theory can predict.

In natural units \(\hbar=c=1\), let \(\psi\) be the charged Dirac field, \(m\) its mass, \(q\) its signed coupling, and \(A_\mu\) the electromagnetic potential. Local phase covariance replaces the ordinary derivative in the Dirac equation by \(D_\mu\). Expanding that derivative gives the current-potential interaction. The field-strength term supplies the electromagnetic dynamics, so radiation and recoil are calculated within one action rather than prescribed independently.

\subsection*{Physical Construction}
The state carrier is Dirac and electromagnetic field states satisfying the Gauss constraint. The calculation involves the minimally coupled Dirac-Maxwell action and its quantized evolution. Charge conservation, fermionic statistics and gauge constraints restrict amplitudes. The calculated quantities include emission and scattering probabilities, energy shifts and electromagnetic response.

\subsection*{Topic Equations}
\begin{centeredalign}
D_\mu=\partial_\mu+iqA_\mu,\qquad F_{\mu\nu}=\partial_\mu A_\nu-\partial_\nu A_\mu\\
\mathcal L_{\mathrm{QED}}=-\frac14F_{\mu\nu}F^{\mu\nu}+\bar\psi(i\gamma^\mu D_\mu-m)\psi\\
j^\mu=q\bar\psi\gamma^\mu\psi,\qquad\partial_\mu F^{\mu\nu}=j^\nu\\
\psi'=e^{-iq\chi}\psi,\quad A_\mu'=A_\mu+\partial_\mu\chi,\quad F_{\mu\nu}'=F_{\mu\nu}\\
k_\mu\Gamma^\mu(p+k,p)=S^{-1}(p+k)-S^{-1}(p)
\end{centeredalign}

\subsection*{Physical Meaning}
The U(1) connection acts by multiplication, so its components commute and the field strength contains no potential commutator. The classical gauge-field action is quadratic. Photon-photon scattering nevertheless arises through charged-particle loops; absence of a classical photon self-coupling does not mean absence of that quantum process.

A calculation of photon emission can replace a polarization vector by the photon momentum. For a complete on-shell amplitude the resulting longitudinal contribution vanishes. Failure of that cancellation exposes a missing diagram or an inconsistent approximation.

Replacing U(1) by a non-Abelian colour group changes this calculation at the level of the field strength and adds direct gauge-field interactions. The distinction is an algebraic change with physical consequences, not a change of particle names.

\subsection*{Consequences Forced By The Relation}
The antisymmetry of F makes the divergence of Maxwell's equation vanish on its left-hand side, giving conservation of electric current. The Dirac equation and its adjoint give the same continuity relation. Agreement of these two calculations connects matter evolution to the electromagnetic source. For the final identity, \(S\) is the full electron propagator and \(\Gamma\) is the proper vertex with the overall charge removed. The Ward--Takahashi relation ties a longitudinal photon insertion to the change in inverse propagator. Radiative corrections to the vertex and the electron propagator must respect it together.
\subsection*{Domain Of The Construction}
Gauge fixing is needed to invert the photon kinetic operator in perturbation theory. Observable probabilities are independent of that choice when the state preparation and calculation are treated consistently. A prescribed classical potential is an approximation in which the corresponding field degrees of freedom are not evolved quantum mechanically.

\subsection*{Invariance And Realization}
A simultaneous local phase change of the charged field and gauge transformation of the potential preserves current conservation and physical probabilities.

Replacing a classical field by a quantized field permits its fluctuations and recoil to affect charged matter. Coupling and mass values determine numerical predictions within the theory.

\subsection*{Discriminating Consequences}
The Ward--Takahashi identity compares independently calculated vertex and propagator corrections; an inconsistent truncation can violate it even when both calculations are finite.

\subsection*{References For The Physical Derivation}
\begin{itemize}
\item \href{https://arxiv.org/abs/hep-ph/0508242}{A. Grozin, Lectures on QED and QCD, sections 3.1-3.7: Lagrangian, Ward identity and charge renormalization.}
\end{itemize}


\clearpage
\section{Boson}
\label{topic:boson}

\subsection*{Mechanism}
Boson is an exchange-symmetry constructor: identical bosons live in symmetric sectors and can share the same mode.

Relativistic, many-body, field, gauge, geometric, and scale-dependent settings each need their own state space and operator domain, specified for that setting rather than inferred from a single-particle model.

The mechanism is symmetric exchange. Commuting creation and annihilation operators allow arbitrary nonnegative occupation of a mode, which supports field modes, coherent states, condensates, and photon-like observables.

\subsection*{Physical Construction}
The state carrier is Fock space, field configuration space, or a sector selected by charge, spin, momentum, statistics, or gauge data. The calculation involves creation, annihilation, field, charge, spin, Hamiltonian, or scattering operators acting on the admissible sector. Statistics, gauge constraints, commutation or anticommutation rules, domain conditions, and sector labels decide which states are legal. The calculated quantities include occupation number, charge, spin, momentum, energy, correlation function, cross-section, or scattering amplitude.

\subsection*{Topic Equations}
\begin{centeredalign}
\mathcal F_{+}(\mathcal H)=\bigoplus_{n=0}^{\infty}\operatorname{Sym}^n\mathcal H\\
[a_i,a_j^\dagger]=\delta_{ij},\qquad [a_i,a_j]=0\\
n_i=a_i^\dagger a_i\in\{0,1,2,\ldots\}
\end{centeredalign}

\subsection*{Physical Meaning}
Different topics in this branch use different carriers: spinor wave functions, Fock spaces, many-body states, gauge sectors, geometric states, or effective low-energy sectors. Their physical content is fixed by the associated field equation or Hamiltonian, its constraints and domain, and the amplitudes, spectra, charges, or correlation functions it predicts.

Field and many-body mechanisms become experimentally useful when assembled into an ordered intervention. The protocol chapter shows how preparation, controlled evolution, measurement, and correction compose into one executable map.

\subsection*{Invariance And Realization}
The relation between prepared states, observables, and spectral probability measures. The use of eigenvalues, projectors, modes, or outcome channels to represent admissible observations. The dependence of the observable on basis, domain, potential, preparation, or measurement context. The commutator structure that limits which observables can be jointly diagonalized.

The physical carrier: particle, wave, field mode, spin, qubit, detector, or excitation. The representation: wave mechanics, matrix mechanics, density matrices, path integrals, circuits, or fields. Where time dependence is placed: on the state, on the operator, in a propagator, or in a path weight. The implementation of preparation, boundary condition, detector, or outcome channel.

\subsection*{Discriminating Consequences}
A transfer target provides a state space, a transformation law, and a spectral or categorical observable, with one compatibility relation experimentally unresolved. A useful validation varies the basis, domain, or measurement context and measures whether the allowed observable changes while the underlying transformation law remains identifiable. A stronger validation contains two candidate observables whose predicted commutator controls joint resolvability.




\clearpage
\section{Quantum chromodynamics}
\label{topic:quantum-chromodynamics}

\subsection*{Mechanism}
Quarks carry three colour components. Quantum chromodynamics couples them to eight gluon fields through a local SU(3) connection.

The colour connection specifies how quark amplitudes at neighbouring spacetime points are compared. Its action is inseparable from the colour representation of the quark: replacing the three-component carrier by an uncharged scalar would remove the very interaction being described.

The colour matrices do not commute. Their commutator enters the field strength, so the gluon field contributes to its own interactions as well as mediating interactions between quarks. With natural units \(\hbar=c=1\), the quark field of flavour \(f\) is \(q_f\), its mass is \(m_f\), the coupling is \(g_s\), and \(A_\mu\) is a matrix in colour space. The action joins their covariant Dirac motion to the energy of the colour field.

\subsection*{Physical Construction}
The state carrier is quark and gluon field states in the physical SU(3) constraint sector. The calculation involves the QCD action, with covariant Dirac motion and non-Abelian field strength. Gauss constraints, quark statistics and specified boundary conditions select physical states. The calculated quantities include colour-singlet spectra, correlation functions and scattering cross sections.

\subsection*{Topic Equations}
\begin{centeredalign}
[T^a,T^b]=if^{abc}T^c,\qquad\operatorname{Tr}(T^aT^b)=\frac12\delta^{ab}\\
A_\mu=A_\mu^aT^a,\qquad D_\mu=\partial_\mu+ig_s A_\mu\\
F_{\mu\nu}=\partial_\mu A_\nu-\partial_\nu A_\mu+ig_s[A_\mu,A_\nu]\\
\mathcal L_{\mathrm{QCD}}=-\frac12\operatorname{Tr}(F_{\mu\nu}F^{\mu\nu})+\sum_f\bar q_f(i\gamma^\mu D_\mu-m_f)q_f\\
\mu\frac{d g_s}{d\mu}=-\frac{11-2n_f/3}{16\pi^2}g_s^3+O(g_s^5)
\end{centeredalign}

\subsection*{Physical Meaning}
Physical observables include colour-singlet correlation functions, hadron energies and scattering cross sections. A Wilson loop follows a heavy test colour charge around a closed contour and traces over its colour index. Its large-time decay defines the static quark-antiquark potential after the self-energy convention is fixed.

In pure gauge theory a long confining flux tube gives an approximately linear static potential. With dynamical quarks the tube can break through pair creation. The observable must therefore be specified along with the matter content; the same loop notation does not impose the same long-distance force in both theories.

Changing resolution replaces short-distance fluctuations by effective couplings and operators. Heavy-quark effective descriptions preserve selected QCD amplitudes through matching, which supplies a concrete example of transferring a physical relation between different sets of variables.

\subsection*{Consequences Forced By The Relation}
Squaring the field strength produces terms cubic and quartic in the gauge potential. These are the three-gluon and four-gluon interactions. They disappear when the gauge algebra is replaced by an Abelian one; the covariant derivative alone is therefore not the whole distinction between QCD and electrodynamics. For \(n_f\) active flavours with \(11-2n_f/3\) positive, the weak-coupling beta function makes the coupling decrease as the energy scale increases. This permits short-distance perturbation theory. The same expansion loses accuracy when the coupling grows at hadronic scales; confinement is investigated through nonperturbative calculations rather than inferred by extrapolating the one-loop formula.
\subsection*{Domain Of The Construction}
The beta function is the one-loop result for SU(3) with fundamental quarks. Active flavour thresholds require matching. Hadron masses and long-distance forces additionally depend on quark masses and the nonperturbative state of the gauge field.

\subsection*{Invariance And Realization}
Gauge changes rotate colour coordinates while leaving colour-singlet probabilities unchanged. Matching to an effective theory preserves the amplitudes included at the stated order.

Quark masses, the coupling, temperature and volume change spectra and correlations. A different gauge group changes the interaction algebra itself.

\subsection*{Discriminating Consequences}
Static potentials, hadron spectra and short-distance cross sections probe different regimes of the same action and require the corresponding controlled approximations.

\subsection*{References For The Physical Derivation}
\begin{itemize}
\item \href{https://arxiv.org/abs/hep-ph/0001312}{G. S. Bali, QCD forces and heavy quark bound states, sections 3-4 and Appendices B-C: colour fields, static potential and running coupling.}
\end{itemize}


\clearpage
\section{Fermion}
\label{topic:fermion}

\subsection*{Mechanism}
Fermion is an exchange-antisymmetry construction: exchanging two identical fermions reverses the many-body amplitude, so coincident one-particle states are removed from the admissible state space.

The tensor product of one-particle spaces contains both symmetric and antisymmetric states. Identical fermions occupy its antisymmetric sector: exchanging the complete position and spin coordinates changes the sign of the amplitude. This restriction determines which many-particle occupations are allowed for a specified Hamiltonian.

Let \(z=(x,\sigma)\) denote the complete one-particle coordinate, including spin. Antisymmetry makes the amplitude vanish when two \(z\) coordinates coincide. Equal positions with opposite spins do not satisfy that condition: two electrons may occupy one spatial orbital in a spin singlet. The anticommutation relations impose the same exclusion when particles are added to individual spin-orbitals.

\subsection*{Physical Construction}
The state carrier is the antisymmetric many-particle sector of the one-particle Hilbert space, including spin. The calculation involves the many-body Hamiltonian and fermionic creation and annihilation operators. Exchange antisymmetry restricts states and mode occupations; the Hamiltonian specifies interactions and evolution. The calculated quantities include mode occupations, spin-resolved correlations, energies and thermodynamic response.

\subsection*{Topic Equations}
\begin{centeredalign}
z_i=(x_i,\sigma_i),\qquad \Psi(\ldots,z_i,\ldots,z_j,\ldots)=-\Psi(\ldots,z_j,\ldots,z_i,\ldots)\\
\Psi(\ldots,z,\ldots,z,\ldots)=0\\
\mathcal F_{-}(\mathcal H)=\bigoplus_{n=0}^{\infty}\wedge^n\mathcal H\\
\{a_i,a_j^\dagger\}=\delta_{ij},\qquad \{a_i,a_j\}=0\\
n_i=a_i^\dagger a_i,\qquad n_i^2=n_i,\qquad\operatorname{spec}(n_i)=\{0,1\}
\end{centeredalign}

\subsection*{Physical Meaning}
The minus sign acquired under exchange makes every Slater determinant vanish when two columns coincide. In occupation language the same restriction is encoded by anticommutation and by eigenvalues zero or one of each mode-number operator. These are equivalent descriptions of one state-space constraint.

A zero-temperature ideal Fermi gas fills all momentum modes up to the Fermi momentum. Compressing the gas therefore forces particles into higher-momentum states and raises its pressure even when the interaction potential is set to zero.

Fermi--Dirac statistics adds thermal occupation to this exchange-constrained state space. Pairing, bosonization and Jordan--Wigner transformations then test which consequences of fermionic algebra survive on a different carrier.

\subsection*{Consequences Forced By The Relation}
Antisymmetry forces the wave function to vanish when two identical fermions occupy the same one-particle state. The resulting exchange hole is present before a dynamical interaction is specified. For a homogeneous ideal gas at zero temperature, filling distinct momentum spin-orbitals up to the Fermi energy produces degeneracy pressure without pairwise repulsion. Pairing preserves antisymmetry of the constituent fermions. A pair has even fermion parity and may support a bosonic collective order parameter, as in superconductors and superfluid helium-3; it need not be a tightly bound elementary boson.
\subsection*{Transformations To Other Physical Realizations}
In one dimension, hard-core bosons and free fermions can share the same density spectrum. The map changes exchange phases and off-diagonal correlations, so equality of energies does not imply identity of all observables. A Jordan--Wigner map carries local fermionic occupation into spins by adding a parity string. The exchange algebra survives, but locality is transferred into an ordered, generally nonlocal operator. In two dimensions, exchanges are classified by braids rather than only by permutations. Anyonic statistics therefore extends, rather than merely interpolates between, the boson and fermion constructions.
\subsection*{Domain Of The Construction}
The relativistic spin--statistics theorem additionally requires locality, positive energy and the relativistic field framework. Antisymmetric lattice or effective quasiparticle models do not by themselves establish that theorem. Fermi surfaces and degeneracy pressure require a many-mode spectrum and a specified density or particle-number constraint; they do not follow from the word fermion alone.

\subsection*{Invariance And Realization}
Exchange antisymmetry, exterior-product state space, canonical anticommutation and zero-or-one mode occupation are equivalent forms of the fermionic restriction. The node at equal complete coordinates and the same-spin exchange hole survive changes between first-quantized wave functions and second-quantized fields. Antisymmetric correlated states can require a superposition of Slater determinants. Fermion parity remains meaningful when particle number changes, including in paired and superconducting states.

Mass, charge, dispersion, dimensionality, interaction law and gauge representation belong to the physical realization and are not fixed by exchange statistics. A change of carrier can turn local fermion operators into nonlocal strings, as in the Jordan--Wigner transformation. Two-dimensional braid statistics and composite-particle structure alter the exchange construction beyond the elementary boson--fermion dichotomy.

\subsection*{Discriminating Consequences}
The defining relation is antisymmetry or canonical anticommutation and the resulting zero-or-one occupation spectrum. At fixed one-particle spectrum, compare with distinguishable-particle and sign-erased controls to isolate exchange. For a claimed transfer, compare correlations and locality as well as energies; spectral agreement alone is insufficient.

\subsection*{Equation Sources}
\begin{itemize}
\item \href{https://arxiv.org/abs/cond-mat/0005069}{arXiv:cond-mat/0005069}
\item \href{https://arxiv.org/abs/quant-ph/0305150}{arXiv:quant-ph/0305150}
\item \href{https://arxiv.org/abs/hep-ph/0007343}{arXiv:hep-ph/0007343}
\end{itemize}
\subsection*{References For The Physical Derivation}
\begin{itemize}
\item \href{https://arxiv.org/abs/0706.3360}{S. Giorgini, L. P. Pitaevskii and S. Stringari, Theory of ultracold atomic Fermi gases, sections II and V.2.2: occupation and spin-resolved correlations.}
\end{itemize}


\clearpage
\section{Electron}
\label{topic:electron}

\subsection*{Mechanism}
Electron is a charged spinor constructor: its identity is fixed by mass, charge, spin-1/2 representation, fermionic statistics, and electromagnetic coupling.

Relativistic, many-body, field, gauge, geometric, and scale-dependent settings each need their own state space and operator domain, specified for that setting rather than inferred from a single-particle model.

The electron is not a generic object label in this tree. Its native mechanism combines a spinor state, a Schrödinger/Pauli/Dirac generator depending on regime, conserved charge, and fermionic anticommutation. The observables are charge, spin, momentum, energy, and scattering response.

\subsection*{Physical Construction}
The state carrier is Fock space, field configuration space, or a sector selected by charge, spin, momentum, statistics, or gauge data. The calculation involves creation, annihilation, field, charge, spin, Hamiltonian, or scattering operators acting on the admissible sector. Statistics, gauge constraints, commutation or anticommutation rules, domain conditions, and sector labels decide which states are legal. The calculated quantities include occupation number, charge, spin, momentum, energy, correlation function, cross-section, or scattering amplitude.

\subsection*{Topic Equations}
\begin{centeredalign}
(i\hbar\gamma^{\mu}D_{\mu}-mc)\psi=0\\
D_{\mu}=\partial_{\mu}+\frac{ie}{\hbar c}A_{\mu}\\
\{\psi_{\alpha}(x),\psi_{\beta}^{\dagger}(y)\}=\delta_{\alpha\beta}\delta(x-y)
\end{centeredalign}

\subsection*{Physical Meaning}
Different topics in this branch use different carriers: spinor wave functions, Fock spaces, many-body states, gauge sectors, geometric states, or effective low-energy sectors. Their physical content is fixed by the associated field equation or Hamiltonian, its constraints and domain, and the amplitudes, spectra, charges, or correlation functions it predicts.

Field and many-body mechanisms become experimentally useful when assembled into an ordered intervention. The protocol chapter shows how preparation, controlled evolution, measurement, and correction compose into one executable map.

\subsection*{Invariance And Realization}
The relation between prepared states, observables, and spectral probability measures. The use of eigenvalues, projectors, modes, or outcome channels to represent admissible observations. The dependence of the observable on basis, domain, potential, preparation, or measurement context. The commutator structure that limits which observables can be jointly diagonalized.

The physical carrier: particle, wave, field mode, spin, qubit, detector, or excitation. The representation: wave mechanics, matrix mechanics, density matrices, path integrals, circuits, or fields. Where time dependence is placed: on the state, on the operator, in a propagator, or in a path weight. The implementation of preparation, boundary condition, detector, or outcome channel.

\subsection*{Discriminating Consequences}
A transfer target provides a state space, a transformation law, and a spectral or categorical observable, with one compatibility relation experimentally unresolved. A useful validation varies the basis, domain, or measurement context and measures whether the allowed observable changes while the underlying transformation law remains identifiable. A stronger validation contains two candidate observables whose predicted commutator controls joint resolvability.




\clearpage
\section{AdS/CFT correspondence}
\label{topic:ads-cft-correspondence}

\subsection*{Mechanism}
AdS/CFT correspondence is a geometry-translation constructor: bulk gravitational data and boundary field data are treated as dual presentations of one operator structure.

Relativistic, many-body, field, gauge, geometric, and scale-dependent settings each need their own state space and operator domain, specified for that setting rather than inferred from a single-particle model.

AdS/CFT belongs at the interface where geometry becomes a representation of the quantum construction rather than the invariant root. The practical content is a dictionary between bulk fields and boundary operators.

\subsection*{Physical Construction}
The state carrier is a spacetime, boundary algebra, gauge orbit, spin network, bulk/boundary pair, or geometric representation of a quantum state space. The calculation involves hamiltonian, action, constraint, boundary operator, correlation map, or dictionary between two representations. Gauge, boundary, metric, covariance, and constraint conditions decide which geometric descriptions represent the same physical content. The calculated quantities include boundary correlators, spectra, entropies, scattering data, geometric invariants, or reconstructed bulk quantities.

\subsection*{Topic Equations}
\begin{centeredalign}
Z_{\mathrm{grav}}[\phi_0]\simeq Z_{\mathrm{CFT}}[J=\phi_0]\\
\left\langle \exp\!\int J\mathcal O\right\rangle_{\mathrm{CFT}}=Z_{\mathrm{bulk}}[\phi|_{\partial}=J]
\end{centeredalign}

\subsection*{Physical Meaning}
Different topics in this branch use different carriers: spinor wave functions, Fock spaces, many-body states, gauge sectors, geometric states, or effective low-energy sectors. Their physical content is fixed by the associated field equation or Hamiltonian, its constraints and domain, and the amplitudes, spectra, charges, or correlation functions it predicts.

Field and many-body mechanisms become experimentally useful when assembled into an ordered intervention. The protocol chapter shows how preparation, controlled evolution, measurement, and correction compose into one executable map.

\subsection*{Invariance And Realization}
The rule connecting prepared states, observables, and spectral probability measures across wave, matrix, path-integral, circuit, or field notation. The operator-to-spectrum relation: admissible observations are represented through eigenvalues, projections, modes, or outcome channels. The dependence of admissible observable on measurement context or boundary condition. The non-commuting compatibility structure, which survives changes of representation.

The name of the carrier: particle, wave, field, qubit, or excitation. Where time dependence is represented: on the state, on the operator, or in a path weight. The coordinate system, basis, or geometric picture used to display the same relation. The physical implementation of detector, boundary, preparation, or observable.

\subsection*{Discriminating Consequences}
Commutation, anticommutation, gauge, and occupation rules define which many-mode states are admissible. The field or many-mode construction must reduce to the appropriate single-particle, quasiparticle, or low-energy limit when those limits exist. Dual presentations are credible only when observables, spectra, or correlation functions are preserved across the translation.




\clearpage
\section{Bose–Einstein statistics}
\label{topic:bose-einstein-statistics}

\subsection*{Mechanism}
Bose–Einstein statistics is a symmetric-state constructor: exchange statistics permits collective occupation of one quantum mode.

Bose-Einstein statistics is needed because identical bosons occupy symmetric many-particle states. Multiple particles may share one mode, which produces collective occupation effects absent from distinguishable particles and fermions.

Bose-Einstein statistics is an admissibility rule for symmetric many-particle states. Commutation permits unrestricted mode occupation, while the thermal distribution determines the mean population and the conditions for macroscopic ground-state occupation.

\subsection*{Physical Construction}
The state carrier is a bosonic Fock space assembled from symmetric many-particle sectors or occupation-number modes. The calculation involves creation, annihilation, and number operators obeying canonical commutation relations. Exchange symmetry permits any non-negative integer occupation of a one-particle mode. The calculated quantities include mode occupations, condensate fraction, particle density, pressure, heat capacity, and coherence observables.

\subsection*{Topic Equations}
\begin{centeredalign}
[a_i,a_j^\dagger]=\delta_{ij},\qquad n_i\in\{0,1,2,\ldots\}\\
\bar n_i=\frac{1}{e^{\beta(\varepsilon_i-\mu)}-1}
\end{centeredalign}

\subsection*{Physical Meaning}
The Bose-Einstein distribution gives the mean occupation of each energy mode at thermal equilibrium. The chemical potential and density determine whether a macroscopic population accumulates in the lowest mode, as in Bose-Einstein condensation.

Cooling a dilute bosonic gas below its critical temperature transfers a macroscopic fraction of the atoms into one quantum mode. The condensate is not caused by an attractive binding force; it follows from the symmetric occupation rule and thermodynamic constraints.

Fock space provides one common language for bosonic and fermionic occupations, while field operators create and remove individual mode excitations.

\subsection*{Invariance And Realization}
Exchange symmetry and canonical commutation define the bosonic many-particle sectors. Each one-particle mode admits any non-negative integer occupation. The Bose-Einstein function gives the equilibrium mean occupation once energy, temperature, and chemical potential are specified.

The dispersion relation, dimensionality, density of states, conserved particle number, and interaction approximation depend on the physical system. Photons, cold atoms, phonons, and bosonic quasiparticles realize the same exchange rule with different chemical-potential constraints. Macroscopic occupation of the lowest mode occurs only when the density of states and thermodynamic limit support condensation.

\subsection*{Discriminating Consequences}
The theory recovers unrestricted mode occupation and Bose enhancement relative to distinguishable particles. The theory recovers the Maxwell-Boltzmann distribution in the dilute low-fugacity limit. The defining relation is particle-number or energy constraints and test whether the inferred condensate fraction follows the correct dimensional dependence.




\clearpage

\section{Quantum statistical mechanics}

\label{topic:quantum-statistical-mechanics}

Thermal equilibrium assigns probabilities to quantum energy states while respecting the quantities that can be exchanged with the surroundings. For a system that exchanges energy with a large reservoir at temperature \(T\) but has fixed particle number, the canonical state is determined by the Hamiltonian \(H\):

\begin{centeredalign}
\rho_\beta=\frac{e^{-\beta H}}{Z},\qquad
Z=\operatorname{Tr}e^{-\beta H},\qquad \beta=(k_BT)^{-1}.
\end{centeredalign}

The trace sums over the actual many-particle state space, including exchange statistics and boundary conditions. It is not an independent classical distribution imposed on particle positions. Derivatives of the same partition function connect average energy and fluctuations:

\begin{centeredalign}
\langle H\rangle=-\partial_\beta\log Z,\qquad
\operatorname{Var}H=\partial_\beta^2\log Z,\qquad
C=\frac{\operatorname{Var}H}{k_BT^2}.
\end{centeredalign}

The last equality holds for a temperature-independent Hamiltonian in the canonical ensemble. It explains why equilibrium energy fluctuations measure the heat capacity: both are responses of the same probability weights to temperature. A different ensemble constrains different fluctuations and must be analyzed separately.

For independent modes of energy \(\epsilon\) exchanging particles with a reservoir at chemical potential \(\mu\), set \(z=\exp[-\beta(\epsilon-\mu)]\). A fermionic mode contributes \(1+z\) to the grand partition function because only occupations zero and one are allowed. A bosonic mode contributes the geometric sum \((1-z)^{-1}\) when \(z<1\). Differentiation gives

\begin{centeredalign}
\overline n_F=\frac{1}{e^{\beta(\epsilon-\mu)}+1},\qquad
\overline n_B=\frac{1}{e^{\beta(\epsilon-\mu)}-1}.
\end{centeredalign}

The sign difference follows from the allowed state occupations. It produces Fermi filling and Bose enhancement before a specific interaction is added. In the dilute limit both expressions approach the Maxwell--Boltzmann weight. Interactions generally prevent factorization into independent mode contributions, although an effective quasiparticle description may recover it approximately.

For a single harmonic mode, the same sum gives \(Z=[2\sinh(\beta\hbar\omega/2)]^{-1}\) and mean energy \((\hbar\omega/2)\coth(\beta\hbar\omega/2)\). The mean approaches the zero-point energy as temperature tends to zero; quadrature fluctuations persist while the energy variance vanishes. At high temperature the mean energy approaches \(k_BT\). This links the oscillator algebra to thermodynamic response through the choice of preparation.

Equilibrium statistics do not specify the rate of approach to equilibrium. Relaxation requires dynamics, reservoir coupling and sometimes additional conserved quantities. A Gibbs state can describe stationary observables while giving no information about a transport coefficient or memory time. Constructing those predictions requires extending the equilibrium description by the corresponding physical evolution.

\subsection*{References For The Physical Derivation}
\begin{itemize}
\item \href{https://arxiv.org/abs/0706.3360}{S. Giorgini, L. P. Pitaevskii and S. Stringari, Theory of ultracold atomic Fermi gases; ideal and interacting quantum gases.}
\end{itemize}

\subsection*{Relations In The Original Papers}

\href{https://arxiv.org/html/1110.3234v1#S2.E27}{arXiv:1110.3234v1, S2.E27}. A thermal harmonic mode has a geometric distribution of number-state populations. The single-mode bosonic example, not the interacting Fermi-gas discussion.

\clearpage
\section{Quantum gravity}
\label{topic:quantum-gravity}

\subsection*{Mechanism}
Quantum gravity asks whether geometry is a background realization, a constrained quantum carrier, or an emergent observable of a deeper quantum state.

Relativistic, many-body, field, gauge, geometric, and scale-dependent settings each need their own state space and operator domain, specified for that setting rather than inferred from a single-particle model.

The constructor cannot treat quantum gravity as an ordinary wave function on a fixed domain. A candidate theory must specify the state space of geometric and matter degrees of freedom, the constraint or evolution operators, the gauge-invariant or relational observables, and the limit in which classical spacetime is recovered.

\subsection*{Physical Construction}
The state carrier is a state space for geometric and matter degrees of freedom, or a deeper carrier from which geometry is reconstructed. The calculation involves constraint, evolution, or amplitude operators that do not presuppose an unexamined fixed spacetime background. Gauge and diffeomorphism constraints determine the physical state space and which operators are observable. The calculated quantities include relational observables, boundary amplitudes, geometric spectra, or semiclassical spacetime data.

\subsection*{Topic Equations}
\begin{centeredalign}
\Psi\in\mathcal H_{\mathrm{geom}\otimes\mathrm{matter}},\qquad \widehat{\mathcal C}_a\Psi=0\\
[\widehat O_{\mathrm{phys}},\widehat{\mathcal C}_a]\Psi=0\\
\langle\Psi|\widehat g_{\mu\nu}|\Psi\rangle\longrightarrow g^{\mathrm{cl}}_{\mu\nu}\quad\text{in a controlled semiclassical regime}
\end{centeredalign}

\subsection*{Physical Meaning}
Different topics in this branch use different carriers: spinor wave functions, Fock spaces, many-body states, gauge sectors, geometric states, or effective low-energy sectors. Their physical content is fixed by the associated field equation or Hamiltonian, its constraints and domain, and the amplitudes, spectra, charges, or correlation functions it predicts.

Field and many-body mechanisms become experimentally useful when assembled into an ordered intervention. The protocol chapter shows how preparation, controlled evolution, measurement, and correction compose into one executable map.

\subsection*{Invariance And Realization}
The rule connecting prepared states, observables, and spectral probability measures across wave, matrix, path-integral, circuit, or field notation. The operator-to-spectrum relation: admissible observations are represented through eigenvalues, projections, modes, or outcome channels. The dependence of admissible observable on measurement context or boundary condition. The non-commuting compatibility structure, which survives changes of representation.

The name of the carrier: particle, wave, field, qubit, or excitation. Where time dependence is represented: on the state, on the operator, or in a path weight. The coordinate system, basis, or geometric picture used to display the same relation. The physical implementation of detector, boundary, preparation, or observable.

\subsection*{Discriminating Consequences}
Commutation, anticommutation, gauge, and occupation rules define which many-mode states are admissible. The field or many-mode construction must reduce to the appropriate single-particle, quasiparticle, or low-energy limit when those limits exist. Dual presentations are credible only when observables, spectra, or correlation functions are preserved across the translation.




\clearpage
\section{Quantum geometry}
\label{topic:quantum-geometry}

\subsection*{Mechanism}
Quantum geometry is a geometry-realization page: geometric quantities are promoted to quantum observables rather than assumed as a smooth background.

Relativistic, many-body, field, gauge, geometric, and scale-dependent settings each need their own state space and operator domain, specified for that setting rather than inferred from a single-particle model.

Quantum geometry uses a quantum state of geometry, often represented by graph or spin-network data. The operator-to-spectrum step asks for eigenvalues of geometric observables such as area or volume. This places the page near the geometry/boundary interface rather than inside a generic many-mode field layer.

\subsection*{Physical Construction}
The state carrier is a spacetime, boundary algebra, gauge orbit, spin network, bulk/boundary pair, or geometric representation of a quantum state space. The calculation involves hamiltonian, action, constraint, boundary operator, correlation map, or dictionary between two representations. Gauge, boundary, metric, covariance, and constraint conditions decide which geometric descriptions represent the same physical content. The calculated quantities include boundary correlators, spectra, entropies, scattering data, geometric invariants, or reconstructed bulk quantities.

\subsection*{Topic Equations}
\begin{centeredalign}
\mathcal H_{\Gamma}=L^2\!\left(SU(2)^E/SU(2)^V\right),\qquad \ket{\Gamma,j_e,\iota_v}\\
\hat A(S)\ket{\Gamma,j,\iota}=8\pi\gamma\ell_P^2\sum_{e\cap S}\sqrt{j_e(j_e+1)}\,\ket{\Gamma,j,\iota}\\
\text{geometry observable}:\quad \hat G\,\ket{g_i}=g_i\ket{g_i}
\end{centeredalign}

\subsection*{Physical Meaning}
Different topics in this branch use different carriers: spinor wave functions, Fock spaces, many-body states, gauge sectors, geometric states, or effective low-energy sectors. Their physical content is fixed by the associated field equation or Hamiltonian, its constraints and domain, and the amplitudes, spectra, charges, or correlation functions it predicts.

Field and many-body mechanisms become experimentally useful when assembled into an ordered intervention. The protocol chapter shows how preparation, controlled evolution, measurement, and correction compose into one executable map.

\subsection*{Invariance And Realization}
The relation between prepared states, observables, and spectral probability measures. The use of eigenvalues, projectors, modes, or outcome channels to represent admissible observations. The dependence of the observable on basis, domain, potential, preparation, or measurement context. The commutator structure that limits which observables can be jointly diagonalized.

The physical carrier: particle, wave, field mode, spin, qubit, detector, or excitation. The representation: wave mechanics, matrix mechanics, density matrices, path integrals, circuits, or fields. Where time dependence is placed: on the state, on the operator, in a propagator, or in a path weight. The implementation of preparation, boundary condition, detector, or outcome channel.

\subsection*{Discriminating Consequences}
A transfer target provides a state space, a transformation law, and a spectral or categorical observable, with one compatibility relation experimentally unresolved. A useful validation varies the basis, domain, or measurement context and measures whether the allowed observable changes while the underlying transformation law remains identifiable. A stronger validation contains two candidate observables whose predicted commutator controls joint resolvability.




\clearpage
\section{Fock space}
\label{topic:fock-space}

\subsection*{Mechanism}
Fock space is the occupation-number state space: the construction that replaces a fixed-particle Hilbert space by a direct sum over particle number.

Relativistic, many-body, field, gauge, geometric, and scale-dependent settings each need their own state space and operator domain, specified for that setting rather than inferred from a single-particle model.

Fock space changes the carrier of the quantum state. Instead of describing one system in one Hilbert space, it builds sectors with zero, one, two, and more identical quanta, then imposes the bosonic or fermionic exchange rule. Creation and annihilation operators are the native coordinates of this page because they move the state between occupation sectors.

\subsection*{Physical Construction}
The state carrier is Fock space, field configuration space, or a sector selected by charge, spin, momentum, statistics, or gauge data. The calculation involves creation, annihilation, field, charge, spin, Hamiltonian, or scattering operators acting on the admissible sector. Statistics, gauge constraints, commutation or anticommutation rules, domain conditions, and sector labels decide which states are legal. The calculated quantities include occupation number, charge, spin, momentum, energy, correlation function, cross-section, or scattering amplitude.

\subsection*{Topic Equations}
\begin{centeredalign}
\mathcal F_{\pm}(\mathcal H)=\bigoplus_{n=0}^{\infty} \mathcal S_{\pm}\mathcal H^{\otimes n}\\
[a_i,a_j^\dagger]_{\mp}=\delta_{ij},\qquad [a_i,a_j]_{\mp}=0\\
N=\sum_i a_i^\dagger a_i,\qquad N\ket{n_1,n_2,\ldots}=\left(\sum_i n_i\right)\ket{n_1,n_2,\ldots}
\end{centeredalign}

\subsection*{Physical Meaning}
Different topics in this branch use different carriers: spinor wave functions, Fock spaces, many-body states, gauge sectors, geometric states, or effective low-energy sectors. Their physical content is fixed by the associated field equation or Hamiltonian, its constraints and domain, and the amplitudes, spectra, charges, or correlation functions it predicts.

Field and many-body mechanisms become experimentally useful when assembled into an ordered intervention. The protocol chapter shows how preparation, controlled evolution, measurement, and correction compose into one executable map.

\subsection*{Invariance And Realization}
The relation between prepared states, observables, and spectral probability measures. The use of eigenvalues, projectors, modes, or outcome channels to represent admissible observations. The dependence of the observable on basis, domain, potential, preparation, or measurement context. The commutator structure that limits which observables can be jointly diagonalized.

The physical carrier: particle, wave, field mode, spin, qubit, detector, or excitation. The representation: wave mechanics, matrix mechanics, density matrices, path integrals, circuits, or fields. Where time dependence is placed: on the state, on the operator, in a propagator, or in a path weight. The implementation of preparation, boundary condition, detector, or outcome channel.

\subsection*{Discriminating Consequences}
A transfer target provides a state space, a transformation law, and a spectral or categorical observable, with one compatibility relation experimentally unresolved. A useful validation varies the basis, domain, or measurement context and measures whether the allowed observable changes while the underlying transformation law remains identifiable. A stronger validation contains two candidate observables whose predicted commutator controls joint resolvability.




\chapter{Topic Reference Index}
This index retains every subject in the field map. A worked treatment has its own section in the preceding chapters. An overview entry records an editorial placement only; no repeated branch equation is offered as a derivation of that subject. Historical entries and alternative names are distinguished from physical treatments.
\section*{State space, domain, and representation}
\begin{longtable}{p{0.57\linewidth}p{0.35\linewidth}}
\toprule Subject & Treatment in this edition \\
\midrule\endhead
\phantomsection\label{index:transformation-theory-quantum-mechanics}Transformation theory (quantum mechanics) & Overview entry; no individual derivation \\
\phantomsection\label{index:mathematical-formulation-of-quantum-mechanics}Mathematical formulation of quantum mechanics & Overview entry; no individual derivation \\
\phantomsection\label{index:hilbert-space}Hilbert space & \hyperref[topic:hilbert-space]{Worked treatment, p.~\pageref*{topic:hilbert-space}} \\
\phantomsection\label{index:quantum-differential-calculus}Quantum differential calculus & Overview entry; no individual derivation \\
\phantomsection\label{index:quantum-mechanics}Quantum mechanics & \hyperref[topic:quantum-mechanics]{Worked treatment, p.~\pageref*{topic:quantum-mechanics}} \\
\phantomsection\label{index:old-quantum-theory}Old quantum theory & \hyperref[topic:old-quantum-theory]{Worked treatment, p.~\pageref*{topic:old-quantum-theory}} \\
\phantomsection\label{index:fourier-transform}Fourier transform & \hyperref[topic:fourier-transform]{Worked treatment, p.~\pageref*{topic:fourier-transform}} \\
\bottomrule
\end{longtable}
\section*{Quantum states and subsystem structure}
\begin{longtable}{p{0.57\linewidth}p{0.35\linewidth}}
\toprule Subject & Treatment in this edition \\
\midrule\endhead
\phantomsection\label{index:density-matrix}Density matrix & \hyperref[topic:density-matrix]{Worked treatment, p.~\pageref*{topic:density-matrix}} \\
\phantomsection\label{index:quantum-superposition}Quantum superposition & Overview entry; no individual derivation \\
\phantomsection\label{index:quantum-decoherence}Quantum decoherence & \hyperref[topic:quantum-decoherence]{Worked treatment, p.~\pageref*{topic:quantum-decoherence}} \\
\phantomsection\label{index:coherence-physics}Coherence (physics) & Overview entry; no individual derivation \\
\phantomsection\label{index:superposition-principle}Superposition principle & Alternative name: Quantum superposition \\
\phantomsection\label{index:wave-function}Wave function & \hyperref[topic:wave-function]{Worked treatment, p.~\pageref*{topic:wave-function}} \\
\phantomsection\label{index:quantum-state}Quantum state & \hyperref[topic:quantum-state]{Worked treatment, p.~\pageref*{topic:quantum-state}} \\
\phantomsection\label{index:two-state-quantum-system}Two-state quantum system & \hyperref[topic:two-state-quantum-system]{Worked treatment, p.~\pageref*{topic:two-state-quantum-system}} \\
\phantomsection\label{index:qubit}Qubit & \hyperref[topic:qubit]{Worked treatment, p.~\pageref*{topic:qubit}} \\
\phantomsection\label{index:quantum-number}Quantum number & Overview entry; no individual derivation \\
\phantomsection\label{index:spin-physics}Spin (physics) & Overview entry; no individual derivation \\
\phantomsection\label{index:schr-dinger-s-cat}Schrödinger's cat & Overview entry; no individual derivation \\
\phantomsection\label{index:quantum-entanglement}Quantum entanglement & \hyperref[topic:quantum-entanglement]{Worked treatment, p.~\pageref*{topic:quantum-entanglement}} \\
\phantomsection\label{index:quantum-biology}Quantum biology & Overview entry; no individual derivation \\
\phantomsection\label{index:macroscopic-quantum-phenomena}Macroscopic quantum phenomena & Overview entry; no individual derivation \\
\phantomsection\label{index:quantum-fluctuation}Quantum fluctuation & Overview entry; no individual derivation \\
\bottomrule
\end{longtable}
\section*{Dynamics and transformations}
\begin{longtable}{p{0.57\linewidth}p{0.35\linewidth}}
\toprule Subject & Treatment in this edition \\
\midrule\endhead
\phantomsection\label{index:unitary-operator}Unitary operator & \hyperref[topic:unitary-operator]{Worked treatment, p.~\pageref*{topic:unitary-operator}} \\
\phantomsection\label{index:perturbation-theory}Perturbation theory & Overview entry; no individual derivation \\
\phantomsection\label{index:quantum-dynamics}Quantum dynamics & Overview entry; no individual derivation \\
\phantomsection\label{index:path-integral-formulation}Path integral formulation & \hyperref[topic:path-integral-formulation]{Worked treatment, p.~\pageref*{topic:path-integral-formulation}} \\
\phantomsection\label{index:hamiltonian-mechanics}Hamiltonian mechanics & Overview entry; no individual derivation \\
\phantomsection\label{index:path-integral}Path integral & \hyperref[topic:path-integral]{Worked treatment, p.~\pageref*{topic:path-integral}} \\
\phantomsection\label{index:hamiltonian-quantum-mechanics}Hamiltonian (quantum mechanics) & \hyperref[topic:hamiltonian-quantum-mechanics]{Worked treatment, p.~\pageref*{topic:hamiltonian-quantum-mechanics}} \\
\phantomsection\label{index:perturbation-theory-quantum-mechanics}Perturbation theory (quantum mechanics) & Alternative name: Perturbation theory \\
\phantomsection\label{index:quantum-chaos}Quantum chaos & Overview entry; no individual derivation \\
\phantomsection\label{index:schr-dinger-equation}Schrödinger equation & \hyperref[topic:schr-dinger-equation]{Worked treatment, p.~\pageref*{topic:schr-dinger-equation}} \\
\phantomsection\label{index:symmetry-in-quantum-mechanics}Symmetry in quantum mechanics & \hyperref[topic:symmetry-in-quantum-mechanics]{Worked treatment, p.~\pageref*{topic:symmetry-in-quantum-mechanics}} \\
\phantomsection\label{index:heisenberg-picture}Heisenberg picture & \hyperref[topic:heisenberg-picture]{Worked treatment, p.~\pageref*{topic:heisenberg-picture}} \\
\phantomsection\label{index:quantum-stochastic-calculus}Quantum stochastic calculus & Overview entry; no individual derivation \\
\phantomsection\label{index:relativistic-quantum-mechanics}Relativistic quantum mechanics & Overview entry; no individual derivation \\
\phantomsection\label{index:schr-dinger-picture}Schrödinger picture & Overview entry; no individual derivation \\
\phantomsection\label{index:creation-and-annihilation-operators}Creation and annihilation operators & \hyperref[topic:creation-and-annihilation-operators]{Worked treatment, p.~\pageref*{topic:creation-and-annihilation-operators}} \\
\bottomrule
\end{longtable}
\section*{Observables and spectra}
\begin{longtable}{p{0.57\linewidth}p{0.35\linewidth}}
\toprule Subject & Treatment in this edition \\
\midrule\endhead
\phantomsection\label{index:spectral-theory}Spectral theory & \hyperref[topic:spectral-theory]{Worked treatment, p.~\pageref*{topic:spectral-theory}} \\
\phantomsection\label{index:pauli-matrices}Pauli matrices & \hyperref[topic:pauli-matrices]{Worked treatment, p.~\pageref*{topic:pauli-matrices}} \\
\phantomsection\label{index:angular-momentum-operator}Angular momentum operator & \hyperref[topic:angular-momentum-operator]{Worked treatment, p.~\pageref*{topic:angular-momentum-operator}} \\
\phantomsection\label{index:observable}Observable & \hyperref[topic:observable]{Worked treatment, p.~\pageref*{topic:observable}} \\
\phantomsection\label{index:self-adjoint-operator}Self-adjoint operator & \hyperref[topic:self-adjoint-operator]{Worked treatment, p.~\pageref*{topic:self-adjoint-operator}} \\
\phantomsection\label{index:operator-theory}Operator theory & Overview entry; no individual derivation \\
\phantomsection\label{index:operator-physics}Operator (physics) & Alternative name: Operator theory \\
\phantomsection\label{index:eigenvalues-and-eigenvectors}Eigenvalues and eigenvectors & Overview entry; no individual derivation \\
\phantomsection\label{index:heisenberg-group}Heisenberg group & \hyperref[topic:heisenberg-group]{Worked treatment, p.~\pageref*{topic:heisenberg-group}} \\
\phantomsection\label{index:von-neumann-algebra}Von Neumann algebra & Overview entry; no individual derivation \\
\phantomsection\label{index:quantum-logic}Quantum logic & Overview entry; no individual derivation \\
\bottomrule
\end{longtable}
\section*{Noncommuting observables}
\begin{longtable}{p{0.57\linewidth}p{0.35\linewidth}}
\toprule Subject & Treatment in this edition \\
\midrule\endhead
\phantomsection\label{index:commutator}Commutator & \hyperref[topic:commutator]{Worked treatment, p.~\pageref*{topic:commutator}} \\
\phantomsection\label{index:uncertainty-principle}Uncertainty principle & \hyperref[topic:uncertainty-principle]{Worked treatment, p.~\pageref*{topic:uncertainty-principle}} \\
\phantomsection\label{index:canonical-commutation-relation}Canonical commutation relation & \hyperref[topic:canonical-commutation-relation]{Worked treatment, p.~\pageref*{topic:canonical-commutation-relation}} \\
\bottomrule
\end{longtable}
\section*{Measurement, instruments, and probabilities}
\begin{longtable}{p{0.57\linewidth}p{0.35\linewidth}}
\toprule Subject & Treatment in this edition \\
\midrule\endhead
\phantomsection\label{index:povm}POVM & \hyperref[topic:povm]{Worked treatment, p.~\pageref*{topic:povm}} \\
\phantomsection\label{index:measurement-in-quantum-mechanics}Measurement in quantum mechanics & \hyperref[topic:measurement-in-quantum-mechanics]{Worked treatment, p.~\pageref*{topic:measurement-in-quantum-mechanics}} \\
\phantomsection\label{index:quantum-jump}Quantum jump & \hyperref[topic:quantum-jump]{Worked treatment, p.~\pageref*{topic:quantum-jump}} \\
\phantomsection\label{index:wave-function-collapse}Wave function collapse & Overview entry; no individual derivation \\
\phantomsection\label{index:born-rule}Born rule & \hyperref[topic:born-rule]{Worked treatment, p.~\pageref*{topic:born-rule}} \\
\phantomsection\label{index:measurement-problem}Measurement problem & \hyperref[topic:measurement-problem]{Worked treatment, p.~\pageref*{topic:measurement-problem}} \\
\phantomsection\label{index:quantum-eraser-experiment}Quantum eraser experiment & Overview entry; no individual derivation \\
\phantomsection\label{index:projection-valued-measure}Projection-valued measure & \hyperref[topic:projection-valued-measure]{Worked treatment, p.~\pageref*{topic:projection-valued-measure}} \\
\phantomsection\label{index:quantum-amplifier}Quantum amplifier & \hyperref[topic:quantum-amplifier]{Worked treatment, p.~\pageref*{topic:quantum-amplifier}} \\
\phantomsection\label{index:bell-s-theorem}Bell's theorem & \hyperref[topic:bell-s-theorem]{Worked treatment, p.~\pageref*{topic:bell-s-theorem}} \\
\phantomsection\label{index:quantum-nonlocality}Quantum nonlocality & Overview entry; no individual derivation \\
\phantomsection\label{index:delayed-choice-quantum-eraser}Delayed-choice quantum eraser & Overview entry; no individual derivation \\
\phantomsection\label{index:quantum-imaging}Quantum imaging & Overview entry; no individual derivation \\
\phantomsection\label{index:wave-particle-duality}Wave–particle duality & Overview entry; no individual derivation \\
\phantomsection\label{index:einstein-podolsky-rosen-paradox}Einstein–Podolsky–Rosen paradox & Overview entry; no individual derivation \\
\phantomsection\label{index:electron-microscope}Electron microscope & Overview entry; no individual derivation \\
\bottomrule
\end{longtable}
\section*{Control sequences and quantum channels}
\begin{longtable}{p{0.57\linewidth}p{0.35\linewidth}}
\toprule Subject & Treatment in this edition \\
\midrule\endhead
\phantomsection\label{index:quantum-computing}Quantum computing & Overview entry; no individual derivation \\
\phantomsection\label{index:quantum-information-science}Quantum information science & Alternative name: Quantum information \\
\phantomsection\label{index:quantum-network}Quantum network & Overview entry; no individual derivation \\
\phantomsection\label{index:quantum-algorithm}Quantum algorithm & Overview entry; no individual derivation \\
\phantomsection\label{index:quantum-error-correction}Quantum error correction & \hyperref[topic:quantum-error-correction]{Worked treatment, p.~\pageref*{topic:quantum-error-correction}} \\
\phantomsection\label{index:quantum-channel}Quantum channel & \hyperref[topic:quantum-channel]{Worked treatment, p.~\pageref*{topic:quantum-channel}} \\
\phantomsection\label{index:quantum-logic-gate}Quantum logic gate & Overview entry; no individual derivation \\
\phantomsection\label{index:quantum-neural-network}Quantum neural network & Overview entry; no individual derivation \\
\phantomsection\label{index:quantum-circuit}Quantum circuit & \hyperref[topic:quantum-circuit]{Worked treatment, p.~\pageref*{topic:quantum-circuit}} \\
\phantomsection\label{index:quantum-information}Quantum information & Overview entry; no individual derivation \\
\phantomsection\label{index:quantum-key-distribution}Quantum key distribution & \hyperref[topic:quantum-key-distribution]{Worked treatment, p.~\pageref*{topic:quantum-key-distribution}} \\
\phantomsection\label{index:quantum-sensor}Quantum sensor & Overview entry; no individual derivation \\
\phantomsection\label{index:quantum-machine}Quantum machine & Overview entry; no individual derivation \\
\phantomsection\label{index:quantum-teleportation}Quantum teleportation & \hyperref[topic:quantum-teleportation]{Worked treatment, p.~\pageref*{topic:quantum-teleportation}} \\
\phantomsection\label{index:quantum-finite-automaton}Quantum finite automaton & Overview entry; no individual derivation \\
\phantomsection\label{index:quantum-machine-learning}Quantum machine learning & Overview entry; no individual derivation \\
\phantomsection\label{index:quantum-programming}Quantum programming & Overview entry; no individual derivation \\
\phantomsection\label{index:quantum-simulator}Quantum simulator & \hyperref[topic:quantum-simulator]{Worked treatment, p.~\pageref*{topic:quantum-simulator}} \\
\phantomsection\label{index:quantum-cryptography}Quantum cryptography & Overview entry; no individual derivation \\
\phantomsection\label{index:quantum-metrology}Quantum metrology & \hyperref[topic:quantum-metrology]{Worked treatment, p.~\pageref*{topic:quantum-metrology}} \\
\phantomsection\label{index:quantum-bus}Quantum bus & Overview entry; no individual derivation \\
\phantomsection\label{index:quantum-image-processing}Quantum image processing & Overview entry; no individual derivation \\
\phantomsection\label{index:quantum-engineering}Quantum engineering & Overview entry; no individual derivation \\
\phantomsection\label{index:quantum-cellular-automaton}Quantum cellular automaton & Overview entry; no individual derivation \\
\phantomsection\label{index:quantum-complexity-theory}Quantum complexity theory & Overview entry; no individual derivation \\
\bottomrule
\end{longtable}
\section*{Boundaries and operator domains}
\begin{longtable}{p{0.57\linewidth}p{0.35\linewidth}}
\toprule Subject & Treatment in this edition \\
\midrule\endhead
\phantomsection\label{index:potential-well}Potential well & Overview entry; no individual derivation \\
\phantomsection\label{index:scattering}Scattering & Overview entry; no individual derivation \\
\phantomsection\label{index:particle-in-a-box}Particle in a box & \hyperref[topic:particle-in-a-box]{Worked treatment, p.~\pageref*{topic:particle-in-a-box}} \\
\phantomsection\label{index:quantum-optics}Quantum optics & Overview entry; no individual derivation \\
\phantomsection\label{index:wave-interference}Wave interference & Overview entry; no individual derivation \\
\phantomsection\label{index:spectral-line}Spectral line & Overview entry; no individual derivation \\
\phantomsection\label{index:quantum-tunnelling}Quantum tunnelling & \hyperref[topic:quantum-tunnelling]{Worked treatment, p.~\pageref*{topic:quantum-tunnelling}} \\
\phantomsection\label{index:s-matrix}S-matrix & Overview entry; no individual derivation \\
\phantomsection\label{index:quantum-harmonic-oscillator}Quantum harmonic oscillator & \hyperref[topic:quantum-harmonic-oscillator]{Worked treatment, p.~\pageref*{topic:quantum-harmonic-oscillator}} \\
\phantomsection\label{index:quantum-metamaterial}Quantum metamaterial & Overview entry; no individual derivation \\
\bottomrule
\end{longtable}
\section*{Fields, constraints, and scale}
\begin{longtable}{p{0.57\linewidth}p{0.35\linewidth}}
\toprule Subject & Treatment in this edition \\
\midrule\endhead
\phantomsection\label{index:gauge-theory}Gauge theory & \hyperref[topic:gauge-theory]{Worked treatment, p.~\pageref*{topic:gauge-theory}} \\
\phantomsection\label{index:photon}Photon & \hyperref[topic:photon]{Worked treatment, p.~\pageref*{topic:photon}} \\
\phantomsection\label{index:renormalization}Renormalization & \hyperref[topic:renormalization]{Worked treatment, p.~\pageref*{topic:renormalization}} \\
\phantomsection\label{index:quantum-field-theory}Quantum field theory & \hyperref[topic:quantum-field-theory]{Worked treatment, p.~\pageref*{topic:quantum-field-theory}} \\
\phantomsection\label{index:fermi-dirac-statistics}Fermi–Dirac statistics & \hyperref[topic:fermi-dirac-statistics]{Worked treatment, p.~\pageref*{topic:fermi-dirac-statistics}} \\
\phantomsection\label{index:dirac-equation}Dirac equation & Overview entry; no individual derivation \\
\phantomsection\label{index:quantum-electrodynamics}Quantum electrodynamics & \hyperref[topic:quantum-electrodynamics]{Worked treatment, p.~\pageref*{topic:quantum-electrodynamics}} \\
\phantomsection\label{index:boson}Boson & \hyperref[topic:boson]{Worked treatment, p.~\pageref*{topic:boson}} \\
\phantomsection\label{index:quantum-chromodynamics}Quantum chromodynamics & \hyperref[topic:quantum-chromodynamics]{Worked treatment, p.~\pageref*{topic:quantum-chromodynamics}} \\
\phantomsection\label{index:fermion}Fermion & \hyperref[topic:fermion]{Worked treatment, p.~\pageref*{topic:fermion}} \\
\phantomsection\label{index:electron}Electron & \hyperref[topic:electron]{Worked treatment, p.~\pageref*{topic:electron}} \\
\phantomsection\label{index:spin-foam}Spin foam & Overview entry; no individual derivation \\
\phantomsection\label{index:ads-cft-correspondence}AdS/CFT correspondence & \hyperref[topic:ads-cft-correspondence]{Worked treatment, p.~\pageref*{topic:ads-cft-correspondence}} \\
\phantomsection\label{index:loop-quantum-gravity}Loop quantum gravity & Overview entry; no individual derivation \\
\phantomsection\label{index:spin-network}Spin network & Overview entry; no individual derivation \\
\phantomsection\label{index:quantum-chemistry}Quantum chemistry & Overview entry; no individual derivation \\
\phantomsection\label{index:standard-model}Standard Model & Overview entry; no individual derivation \\
\phantomsection\label{index:klein-gordon-equation}Klein–Gordon equation & Overview entry; no individual derivation \\
\phantomsection\label{index:bose-einstein-statistics}Bose–Einstein statistics & \hyperref[topic:bose-einstein-statistics]{Worked treatment, p.~\pageref*{topic:bose-einstein-statistics}} \\
\phantomsection\label{index:string-theory}String theory & Overview entry; no individual derivation \\
\phantomsection\label{index:quantum-statistical-mechanics}Quantum statistical mechanics & \hyperref[topic:quantum-statistical-mechanics]{Worked treatment, p.~\pageref*{topic:quantum-statistical-mechanics}} \\
\phantomsection\label{index:quantum-gravity}Quantum gravity & \hyperref[topic:quantum-gravity]{Worked treatment, p.~\pageref*{topic:quantum-gravity}} \\
\phantomsection\label{index:quantum-geometry}Quantum geometry & \hyperref[topic:quantum-geometry]{Worked treatment, p.~\pageref*{topic:quantum-geometry}} \\
\phantomsection\label{index:quantum-spacetime}Quantum spacetime & Overview entry; no individual derivation \\
\phantomsection\label{index:fock-space}Fock space & \hyperref[topic:fock-space]{Worked treatment, p.~\pageref*{topic:fock-space}} \\
\bottomrule
\end{longtable}
\section*{Annotations: history, interpretations, and popular frames}
\begin{longtable}{p{0.57\linewidth}p{0.35\linewidth}}
\toprule Subject & Treatment in this edition \\
\midrule\endhead
\phantomsection\label{index:introduction-to-quantum-mechanics-book}Introduction to Quantum Mechanics (book) & Historical or interpretive entry \\
\phantomsection\label{index:quantum-mind}Quantum mind & Historical or interpretive entry \\
\phantomsection\label{index:erwin-schr-dinger}Erwin Schrödinger & Historical or interpretive entry \\
\phantomsection\label{index:history-of-quantum-mechanics}History of quantum mechanics & Historical or interpretive entry \\
\phantomsection\label{index:introduction-to-quantum-mechanics}Introduction to quantum mechanics & Historical or interpretive entry \\
\phantomsection\label{index:history-of-quantum-field-theory}History of quantum field theory & Historical or interpretive entry \\
\phantomsection\label{index:david-hilbert}David Hilbert & Historical or interpretive entry \\
\phantomsection\label{index:quantum-computing-a-gentle-introduction}Quantum Computing: A Gentle Introduction & Historical or interpretive entry \\
\phantomsection\label{index:quantum-mysticism}Quantum mysticism & Historical or interpretive entry \\
\phantomsection\label{index:interpretations-of-quantum-mechanics}Interpretations of quantum mechanics & Historical or interpretive entry \\
\phantomsection\label{index:quantum-theory-concepts-and-methods}Quantum Theory: Concepts and Methods & Historical or interpretive entry \\
\phantomsection\label{index:qbism}QBism & Historical or interpretive entry \\
\phantomsection\label{index:relational-quantum-mechanics}Relational quantum mechanics & Historical or interpretive entry \\
\phantomsection\label{index:modern-quantum-mechanics}Modern Quantum Mechanics & Historical or interpretive entry \\
\phantomsection\label{index:applications-of-quantum-mechanics}Applications of quantum mechanics & Historical or interpretive entry \\
\phantomsection\label{index:werner-heisenberg}Werner Heisenberg & Historical or interpretive entry \\
\phantomsection\label{index:the-quantum-universe}The Quantum Universe & Historical or interpretive entry \\
\bottomrule
\end{longtable}

\backmatter

\chapter{Sources And Reproduction}

Equation Sources links corpus-aligned relations for 6 topics. Relations In The Original Papers links 16 displays from 7 articles across 15 topics. The cited relations and their conditions appear beside the relevant treatment. The mechanism tree and cross-topic comparisons are the authors' explanatory organization of established theory. The edition manifest records the individual equation and article identifiers. Independent scientific review of every chapter remains to be completed.

\end{document}